% Aspect ondulatoire de la lumière : diffraction, interférences et effet Doppler — Physique Tle D — Cours signé OnBuch+
% Programme officiel MINESEC. Compiler avec : tectonic P16-aspect-ondulatoire-lumiere.tex
\newcommand{\DOCMATIERE}{Physique}
\newcommand{\DOCNIVEAU}{Tle D}
\newcommand{\DOCMODULE}{Module 4 : Onde, matière et transformations nucléaires}
\newcommand{\DOCLECON}{P16}
\newcommand{\DOCTITRE}{Aspect ondulatoire de la lumière : diffraction, interférences et effet Doppler}
% OnBuch+ — préambule « cours signés » ENRICHI (compilé avec tectonic / XeLaTeX)
% Système visuel « Pop » d'OnBuch+ : fond crème (jamais blanc), bordures encre
% épaisses, ombres « dures » décalées, titres Archivo Black en capitales.
% Version MATHÉMATIQUES : graphes (pgfplots/tikz), boîtes pédagogiques
% (définition, propriété, méthode, attention, le savais-tu, exercice résolu,
% exercice, TP, bilan), page de garde et signature OnBuch+ ancrée en bas.
%
% Macros attendues dans main.tex AVANT \input{preamble} :
%   \DOCMATIERE  \DOCNIVEAU  \DOCMODULE  \DOCLECON (numéro)  \DOCTITRE
\documentclass[11pt]{article}

\usepackage{fontspec}
\usepackage[french]{babel}
\usepackage[a4paper,margin=2cm,top=3.4cm,bottom=2.8cm,headheight=1.4cm,footskip=1.3cm]{geometry}
\usepackage{amsmath,amssymb,amsfonts}
\usepackage{mathtools} % \xmapsto, \coloneqq... souvent utilisés par les modèles
\usepackage{cancel} % simplifications barrées (bilans de forces)
% \abs{z} (module, valeur absolue) et \norm{u} (norme), utilisés par les modèles
\providecommand{\abs}{}\let\abs\relax\DeclarePairedDelimiter{\abs}{\lvert}{\rvert}
\providecommand{\norm}{}\let\norm\relax\DeclarePairedDelimiter{\norm}{\lVert}{\rVert}
\usepackage[table]{xcolor} % [table] : fournit aussi \rowcolors (alternance de lignes)
\usepackage{pagecolor}
\usepackage{tikz}
\usetikzlibrary{arrows.meta,positioning,calc,shapes.geometric,shapes.misc,shapes.arrows,decorations.pathmorphing,decorations.markings,patterns,shadows,mindmap,trees,fit,backgrounds,matrix,intersections,angles,quotes}
\usepackage{pgfplots}
\usepackage[version=4]{mhchem} % équations nucléaires \ce{^{235}_{92}U}
\usepackage[european,siunitx]{circuitikz} % schémas électriques
\pgfplotsset{compat=1.18}
\usepgfplotslibrary{fillbetween} % aires entre courbes (calcul intégral)
\usepackage{enumitem}
\usepackage{fancyhdr}
% [most] charge toutes les bibliothèques tcolorbox (listings, minted, raster,
% documentation...) et gonfle inutilement la mémoire TeX sur un document long
% (~300 boîtes) : on ne charge que ce qui est réellement utilisé.
\usepackage[skins,breakable]{tcolorbox}
\usepackage{graphicx}
\usepackage{colortbl}
\usepackage{booktabs}
\usepackage{tabularx}
\usepackage{diagbox} % case d'en-tête diagonale des tableaux à double entrée
% Colonnes extensibles centrée / gauche / droite (C, L, R), souvent utilisées par les modèles
\newcolumntype{C}{>{\centering\arraybackslash}X}
\newcolumntype{L}{>{\raggedright\arraybackslash}X}
\newcolumntype{R}{>{\raggedleft\arraybackslash}X}
\usepackage{array}
\usepackage{multicol}
\usepackage{siunitx}
% parse-numbers=false : accepte aussi \SI{2,5 \times 10^{-3}}{...}
\sisetup{locale=FR,output-decimal-marker={,},group-minimum-digits=4,parse-numbers=false}
\DeclareSIUnit\ohm{\ensuremath{\symup{\Omega}}} % U+2126 absent de la police
\DeclareSIUnit\division{div} % réglages d'oscilloscope (ms/div, V/div)
\usepackage{qrcode}
% \degree : commande fréquemment utilisée par les modèles pour un angle ou une
% température en dehors de \SI{}{} (non fournie par les paquets chargés).
\providecommand{\degree}{^{\circ}}
% \qed (fin de démonstration), utilisé par les modèles sans amsthm
\providecommand{\qed}{\hfill$\square$}
% \barre : double barre des valeurs interdites dans un tableau de variations (tkz-tab)
\providecommand{\barre}{\ensuremath{\|}}
% environnement proof (amsthm non chargé) utilisé par les modèles
\ifdefined\proof\else\newenvironment{proof}[1][Démonstration]{\par\noindent\textit{#1.}\ }{\qed\par}\fi
\usepackage{lastpage}
\usepackage{hyperref}
\hypersetup{hidelinks}

% ============================================================
% Palette « Pop » OnBuch+ — mêmes hex que la classe P (plus_ui.dart)
% ============================================================
\definecolor{poporange}{HTML}{F59321}
\definecolor{popdark}{HTML}{E07A0C}
\definecolor{popcream}{HTML}{FFF6E8}
\definecolor{popink}{HTML}{1C1714}
\definecolor{poppurple}{HTML}{7A5AE0}
\definecolor{popgreen}{HTML}{2E9E6A}
\definecolor{poppink}{HTML}{E0457B}
\definecolor{popblue}{HTML}{2F7DE1}
\definecolor{popgold}{HTML}{C98A12}
\definecolor{popmuted}{HTML}{8A7F76}
\definecolor{popline}{HTML}{EADFD2}
\definecolor{popgray}{HTML}{8A7F76}
\colorlet{popgrey}{popgray}
% Alias tolérants (le modèle nomme parfois les couleurs différemment)
\colorlet{popwhite}{white}
\colorlet{popblack}{popink}
\colorlet{popred}{poppink}
\colorlet{popyellow}{popgold}
% Teintes claires pour fonds de tableaux / zones
\colorlet{popblueL}{popblue!10!white}
\colorlet{poporangeL}{poporange!14!white}
\colorlet{popgreenL}{popgreen!12!white}
\colorlet{poppurpleL}{poppurple!12!white}
\colorlet{poppinkL}{poppink!12!white}
\colorlet{popgoldL}{popgold!14!white}

\pagecolor{popcream}

% ============================================================
% Typographie
% ============================================================
\defaultfontfeatures{Ligatures=TeX}
\setmainfont{PlusJakartaSans-Medium.ttf}[Path=../../../fonts/,BoldFont=PlusJakartaSans-Bold.ttf]
% Maths en sans-serif (Fira Math) avec chiffres et lettres droites de Plus
% Jakarta, pour que les formules se fondent dans le texte.
\usepackage[math-style=ISO,bold-style=ISO]{unicode-math}
\setmathfont{FiraMath-Regular.otf}[Path=../../../fonts/]
\setmathfont{PlusJakartaSans-Medium.ttf}[Path=../../../fonts/,range={up/num,up/{Latin,latin},\mathup}]
\usepackage{icomma} % virgule décimale sans espace en mode math
% Caractères mathématiques Unicode absents des polices de texte (Plus Jakarta,
% Archivo Black) : rendus par leur équivalent mathématique, en texte comme en
% formule — sinon ils s'affichent en carré vide (ex. « Arithmétique dans ℤ »).
\usepackage{newunicodechar}
\newunicodechar{ℝ}{\ensuremath{\mathbb{R}}}
\newunicodechar{ℤ}{\ensuremath{\mathbb{Z}}}
\newunicodechar{ℂ}{\ensuremath{\mathbb{C}}}
\newunicodechar{ℕ}{\ensuremath{\mathbb{N}}}
\newunicodechar{ℚ}{\ensuremath{\mathbb{Q}}}
\newunicodechar{𝔻}{\ensuremath{\mathbb{D}}}
\newunicodechar{α}{\ensuremath{\alpha}}
\newunicodechar{β}{\ensuremath{\beta}}
\newunicodechar{γ}{\ensuremath{\gamma}}
\newunicodechar{δ}{\ensuremath{\delta}}
\newunicodechar{ε}{\ensuremath{\varepsilon}}
\newunicodechar{θ}{\ensuremath{\theta}}
\newunicodechar{λ}{\ensuremath{\lambda}}
\newunicodechar{μ}{\ensuremath{\mu}}
\newunicodechar{σ}{\ensuremath{\sigma}}
\newunicodechar{τ}{\ensuremath{\tau}}
\newunicodechar{φ}{\ensuremath{\varphi}}
\newunicodechar{ψ}{\ensuremath{\psi}}
\newunicodechar{ω}{\ensuremath{\omega}}
\newunicodechar{Σ}{\ensuremath{\Sigma}}
% majuscules produites par \MakeUppercase dans les titres (α-aminés -> Α)
\newunicodechar{Α}{\ensuremath{\alpha}}
\newunicodechar{Β}{\ensuremath{\beta}}
\newunicodechar{Γ}{\ensuremath{\Gamma}}
\newunicodechar{Δ}{\ensuremath{\Delta}}
\newunicodechar{Ε}{\ensuremath{\varepsilon}}
\newunicodechar{Θ}{\ensuremath{\Theta}}
\newunicodechar{Λ}{\ensuremath{\Lambda}}
\newunicodechar{Μ}{\ensuremath{\mu}}
\newunicodechar{Τ}{\ensuremath{\tau}}
\newunicodechar{Φ}{\ensuremath{\Phi}}
\newunicodechar{Ψ}{\ensuremath{\Psi}}
\newunicodechar{Ω}{\ensuremath{\Omega}}
\newunicodechar{↦}{\ensuremath{\mapsto}}
\newunicodechar{∈}{\ensuremath{\in}}
\newunicodechar{∉}{\ensuremath{\notin}}
\newunicodechar{∧}{\ensuremath{\wedge}}
\newunicodechar{⇔}{\ensuremath{\Leftrightarrow}}
\newunicodechar{⟺}{\ensuremath{\Longleftrightarrow}}
\newunicodechar{⇒}{\ensuremath{\Rightarrow}}
\newunicodechar{⟹}{\ensuremath{\Longrightarrow}}
\newunicodechar{∘}{\ensuremath{\circ}}
\newunicodechar{∪}{\ensuremath{\cup}}
\newunicodechar{∩}{\ensuremath{\cap}}
\newunicodechar{≡}{\ensuremath{\equiv}}
\newunicodechar{⊂}{\ensuremath{\subset}}
\newunicodechar{⊥}{\ensuremath{\perp}}
\newunicodechar{∃}{\ensuremath{\exists}}
\newunicodechar{∀}{\ensuremath{\forall}}
\newunicodechar{∛}{\ensuremath{\sqrt[3]{\,}}}
\newunicodechar{∜}{\ensuremath{\sqrt[4]{\,}}}
\newunicodechar{⃗}{\ensuremath{{}^{\to}}}
\newunicodechar{ⁿ}{\ifmmode^{n}\else\textsuperscript{\ensuremath{n}}\fi}
\newunicodechar{ˣ}{\ifmmode^{x}\else\textsuperscript{\ensuremath{x}}\fi}
\newunicodechar{ᵇ}{\ifmmode^{b}\else\textsuperscript{\ensuremath{b}}\fi}
\newunicodechar{ʳ}{\ifmmode^{r}\else\textsuperscript{\ensuremath{r}}\fi}
\newunicodechar{ᵘ}{\ifmmode^{u}\else\textsuperscript{\ensuremath{u}}\fi}
\newunicodechar{ʸ}{\ifmmode^{y}\else\textsuperscript{\ensuremath{y}}\fi}
\newunicodechar{ᵃ}{\ifmmode^{a}\else\textsuperscript{\ensuremath{a}}\fi}
\newunicodechar{ᵉ}{\ifmmode^{e}\else\textsuperscript{\ensuremath{e}}\fi}
\newunicodechar{ᵏ}{\ifmmode^{k}\else\textsuperscript{\ensuremath{k}}\fi}
\newunicodechar{ᵖ}{\ifmmode^{p}\else\textsuperscript{\ensuremath{p}}\fi}
\newunicodechar{ᴺ}{\ifmmode^{N}\else\textsuperscript{\ensuremath{N}}\fi}
\newunicodechar{ᶜ}{\ifmmode^{c}\else\textsuperscript{\ensuremath{c}}\fi}
\newunicodechar{ʰ}{\ifmmode^{h}\else\textsuperscript{\ensuremath{h}}\fi}
\newunicodechar{ᵠ}{\ifmmode^{q}\else\textsuperscript{\ensuremath{q}}\fi}
\newunicodechar{ₙ}{\ifmmode_{n}\else\textsubscript{\ensuremath{n}}\fi}
\newunicodechar{ₐ}{\ifmmode_{a}\else\textsubscript{\ensuremath{a}}\fi}
\newunicodechar{ᵢ}{\ifmmode_{i}\else\textsubscript{\ensuremath{i}}\fi}
\newunicodechar{ᵣ}{\ifmmode_{r}\else\textsubscript{\ensuremath{r}}\fi}
\newunicodechar{ₖ}{\ifmmode_{k}\else\textsubscript{\ensuremath{k}}\fi}
\newunicodechar{ᵦ}{\ifmmode_{\beta}\else\textsubscript{\ensuremath{\beta}}\fi}
\newunicodechar{ₓ}{\ifmmode_{x}\else\textsubscript{\ensuremath{x}}\fi}
\newfontfamily\popdisplay{ArchivoBlack-Regular.ttf}[Path=../../../fonts/]
\newfontfamily\poplogo{SpaceGrotesk-Bold.ttf}[Path=../../../fonts/]
\newfontfamily\popsemi{PlusJakartaSans-SemiBold.ttf}[Path=../../../fonts/]
\newfontfamily\popxbold{PlusJakartaSans-ExtraBold.ttf}[Path=../../../fonts/]
\newfontfamily\popxboldls{PlusJakartaSans-ExtraBold.ttf}[Path=../../../fonts/,LetterSpace=3.0]

\setlist[enumerate]{leftmargin=1.7em,itemsep=5pt,topsep=4pt}
\setlist[itemize]{leftmargin=1.7em,itemsep=4pt,topsep=4pt,label={\color{poporange}\textbullet}}
\setlength{\parindent}{0pt}
\setlength{\parskip}{5pt}
\renewcommand{\arraystretch}{1.25}

% pgfplots : style Pop par défaut
\pgfplotsset{
  every axis/.append style={axis line style={popink,line width=1pt},tick style={popink},
    grid style={popline,line width=0.5pt},label style={font=\small},tick label style={font=\footnotesize}},
  popcurve/.style={poporange,line width=1.8pt,no markers,smooth},
}
% Même style disponible dans un \draw TikZ simple (hors pgfplots)
\tikzset{popcurve/.style={poporange,line width=1.8pt}}
\tikzset{popgrid/.style={popline,very thin}}
\colorlet{popcurve}{poporange} % le nom du style est parfois utilisé comme couleur

% ============================================================
% Pill / squiggle
% ============================================================
\newcommand{\poppill}[3]{%
  \tikz[baseline=-0.55ex]\node[draw=popink,line width=1.2pt,rounded corners=2.6pt,%
    fill=#2,inner xsep=6.5pt,inner ysep=3pt]{%
    \popxboldls\fontsize{7}{7}\selectfont\color{#3}\MakeUppercase{#1}};%
}
\newcommand{\popsquiggle}[1]{%
  \tikz[baseline=-0.4ex]{
    \draw[#1,line width=2.6pt,line cap=round]
      plot [smooth] coordinates {(0,0.09) (0.34,0.01) (0.68,0.21) (1.02,0.01) (1.36,0.10)};
  }%
}
% Mot-clé surligné (marqueur orange sous le texte)
\newcommand{\cle}[1]{{\popxbold\color{popdark}#1}}

% ============================================================
% En-tête / pied
% ============================================================
\pagestyle{fancy}
\fancyhf{}
\renewcommand{\headrulewidth}{0pt}
\renewcommand{\footrulewidth}{0pt}
\lhead{\raisebox{-0.15ex}{\poplogo\fontsize{13}{13}\selectfont\color{popink}OnBuch\color{poporange}+}}
\rhead{\poppill{\DOCMATIERE\ \textbullet\ \DOCNIVEAU\ \textbullet\ Leçon \DOCLECON}{white}{popink}}
\lfoot{\popxbold\fontsize{7.6}{7.6}\selectfont\color{popmuted}\MakeUppercase{Cours signé OnBuch+}\ \textbullet\ {\popsemi\DOCTITRE}}
\rfoot{\tikz[baseline=-0.6ex]\node[rounded corners=8.5pt,fill=popink,minimum height=17pt,minimum width=17pt,inner xsep=5pt,inner sep=0]{\popxbold\fontsize{8}{8}\selectfont\color{popcream}\thepage\,/\,\pageref*{LastPage}};}

% ============================================================
% Titres
% ============================================================
\newcommand{\chaptitre}[2]{%
  \begin{tcolorbox}[enhanced,colback=poporange,colframe=popink,boxrule=2.6pt,arc=11pt,
      drop shadow=popink,left=15pt,right=15pt,top=12pt,bottom=11pt,before skip=3pt,after skip=18pt]
    \poppill{#2}{popink}{popcream}\par\vspace{9pt}
    {\raggedright\hyphenpenalty=10000\exhyphenpenalty=10000\popdisplay\fontsize{22}{24}\selectfont\color{popcream}\MakeUppercase{#1}\par}
  \end{tcolorbox}
}
% Section numérotée : pastille orange avec le numéro + titre Archivo
\newcounter{popsec}
\newcommand{\coursec}[1]{%
  \stepcounter{popsec}\setcounter{popsub}{0}%
  \par\vspace{16pt}\noindent
  \begin{minipage}{\linewidth}
  \tikz[baseline=-0.7ex]\node[draw=popink,line width=1.6pt,fill=poporange,rounded corners=4pt,minimum size=22pt,inner sep=0]{\popdisplay\fontsize{12}{12}\selectfont\color{popcream}\thepopsec};%
  \hspace{8pt}{\popdisplay\fontsize{15}{17}\selectfont\color{popink}\MakeUppercase{#1}}\par
  \vspace{4pt}\noindent{\color{poporange}\rule{36pt}{3.6pt}}
  \end{minipage}\par\nopagebreak\vspace{8pt}
}
\newcounter{popsub}
\newcommand{\courssub}[1]{%
  \stepcounter{popsub}%
  \par\vspace{9pt}\noindent{\popxbold\fontsize{11.5}{13}\selectfont\color{popink}{\color{poporange}\thepopsec.\thepopsub}\hspace{6pt}#1}\par\nopagebreak\vspace{4pt}
}

% ============================================================
% Boîtes pédagogiques — même recette : fond blanc, bordure épaisse colorée,
% ombre dure encre, badge plein. Titre optionnel [#1].
% ============================================================
\tcbset{popbase/.style={breakable,enhanced,colback=white,boxrule=2.3pt,arc=9pt,
  shadow={3pt}{-3pt}{0pt}{popink},left=14pt,right=14pt,top=11pt,bottom=11pt,before skip=11pt,after skip=15pt,
  fontupper=\popsemi\fontsize{10.6}{14.5}\selectfont\color{popink},
  fontlower=\popsemi\fontsize{10.6}{14.5}\selectfont\color{popink}}}
% Coche dessinée (le glyphe ✓ n'existe pas dans les polices du document)
\newcommand{\popcheck}[1]{\tikz[baseline=-0.5ex]\draw[#1,line width=1.6pt,line cap=round,line join=round](-0.13,0.0)--(-0.04,-0.09)--(0.13,0.1);}
\providecommand{\coche}{\popcheck{popgreen}}
\providecommand{\resultat}[1]{\par\smallskip\noindent\fcolorbox{popgreen}{popgreenL}{\parbox{\dimexpr\linewidth-2\fboxsep-2\fboxrule\relax}{\textbf{Résultat.} #1}}\par\smallskip}
\newcommand{\popboxhead}[3]{\poppill{#1}{#2}{white}\ifx\relax#3\relax\else\hspace{6pt}{\popxbold\fontsize{10.6}{12}\selectfont\color{popink}#3}\fi\par\vspace{7pt}}

\newtcolorbox{aretenir}[1][]{popbase,colframe=popgreen,before upper={\popboxhead{À retenir}{popgreen}{#1}}}
\newtcolorbox{exemplebox}[1][]{popbase,colframe=poporange,before upper={\popboxhead{Exemple}{poporange}{#1}}}
\newtcolorbox{definition}[1][]{popbase,colframe=popblue,before upper={\popboxhead{Définition}{popblue}{#1}}}
\newtcolorbox{propriete}[1][]{popbase,colframe=poppurple,before upper={\popboxhead{Propriété}{poppurple}{#1}}}
\newtcolorbox{methode}[1][]{popbase,colframe=popgold,before upper={\popboxhead{Méthode}{popgold}{#1}}}
\newtcolorbox{attention}[1][]{popbase,colframe=poppink,before upper={\popboxhead{Attention}{poppink}{#1}}}
\newtcolorbox{experience}[1][]{popbase,colframe=popblue,colback=popblueL,before upper={\popboxhead{Expérience}{popblue}{#1}}}
% « Le savais-tu ? » : carte violette pleine (contexte, culture, Cameroun)
\newtcolorbox{savaistu}[1][]{popbase,colback=poppurple,colframe=popink,
  fontupper=\popsemi\fontsize{10.4}{14.2}\selectfont\color{white},
  before upper={\poppill{Le savais-tu\,?}{white}{poppurple}\ifx\relax#1\relax\else\hspace{6pt}{\popxbold\color{white}#1}\fi\par\vspace{7pt}}}
% Exercice résolu : énoncé en haut, \tcblower puis la solution sur fond crème
\newcounter{popexo}\newcounter{popexr}
\newtcolorbox{exoresolu}[1][]{popbase,colframe=poporange,colbacklower=popcream,
  segmentation style={popink,line width=1.2pt,dashed},
  before upper={\stepcounter{popexr}\popboxhead{Exercice résolu \thepopexr}{poporange}{#1}},
  before lower={\poppill{Solution}{popink}{popcream}\par\vspace{6pt}}}
% Exercice (non résolu en place) — la correction est dans la partie Corrigés
\newtcolorbox{exercice}[1][]{popbase,colframe=popink,
  before upper={\stepcounter{popexo}\popboxhead{Exercice \thepopexo}{popink}{#1}}}
% Corrigé
\newtcolorbox{corrige}[1][]{popbase,colframe=popgreen,colback=popgreenL,
  before upper={\popboxhead{Corrigé}{popgreen}{#1}}}
% Objectifs / prérequis / bilan
\newtcolorbox{objectifs}{popbase,colframe=popink,colback=poporangeL,
  before upper={\popboxhead{Objectifs de la leçon}{popink}{}}}
\newtcolorbox{prerequis}{popbase,colframe=popink,colback=popblueL,
  before upper={\popboxhead{Prérequis}{popblue}{}}}
\newtcolorbox{bilan}{popbase,colframe=popink,colback=white,boxrule=2.8pt,
  before upper={\popboxhead{Fiche bilan}{poporange}{L'essentiel à connaître pour l'examen}}}
% Figure encadrée avec légende
\newtcolorbox{popfigure}[1][]{enhanced,breakable=false,colback=white,colframe=popline,boxrule=1.4pt,arc=8pt,
  left=10pt,right=10pt,top=10pt,bottom=8pt,before skip=10pt,after skip=12pt,halign=center,
  lower separated=false,halign lower=center,
  fontlower=\popsemi\fontsize{9}{11.5}\selectfont\color{popmuted},
  #1}
\newcommand{\legende}[1]{\par\vspace{6pt}{\popsemi\fontsize{9}{11.5}\selectfont\color{popmuted}#1}}

% ============================================================
% Signature OnBuch+
% ============================================================
\newcommand{\poplogoinline}[1]{{\poplogo\fontsize{#1}{#1}\selectfont\color{popink}OnBuch\color{poporange}+}}
\newcommand{\signatureonbuch}{%
  \begin{tcolorbox}[enhanced,colback=popink,colframe=popink,boxrule=2.2pt,arc=10pt,
      drop shadow=poporange,left=14pt,right=14pt,top=10pt,bottom=10pt,before skip=0pt,after skip=0pt]
    \tikz[baseline=-0.7ex]\node[circle,fill=popgreen,draw=popcream,line width=1.3pt,minimum size=22pt,inner sep=0]{\popcheck{white}};%
    \hspace{9pt}\begin{minipage}[c]{0.62\linewidth}
      {\popxbold\fontsize{12}{13}\selectfont\color{popcream}Signé\ \textbullet\ {\poplogo OnBuch\color{poporange}+}}\par\vspace{2pt}
      {\raggedright\popsemi\fontsize{8}{10.5}\selectfont\color{popline}\MakeUppercase{Équipe pédagogique OnBuch+}\\\MakeUppercase{Conforme au programme officiel MINESEC}\par}
    \end{minipage}\hfill
    {\poplogo\fontsize{22}{22}\selectfont\color{popcream}OnBuch\color{poporange}+}
  \end{tcolorbox}%
}

% ============================================================
% Page de garde
% #1 titre · #2 matière · #3 niveau · #4 module · #5 n° leçon · #6 durée ·
% #7 description / objectifs (texte ou itemize)
% ============================================================
\newcommand{\pagegarde}[7]{%
  \thispagestyle{empty}
  \newgeometry{a4paper,margin=1.7cm,top=1.6cm,bottom=1.4cm}%
  \begin{tikzpicture}[remember picture,overlay]
    \fill[poporange!18] ([xshift=-3.2cm,yshift=-3.6cm]current page.north east) circle (4.6cm);
    \fill[poppurple!14] ([xshift=2.4cm,yshift=7.6cm]current page.south west) circle (3.8cm);
  \end{tikzpicture}%
  {\poplogo\fontsize{30}{30}\selectfont\color{popink}OnBuch\color{poporange}+}\hfill
  \raisebox{0.6ex}{\poppill{Cours signé \textbullet\ Premium}{popink}{popcream}}\par
  \vspace{1pt}\popsquiggle{poppurple}\par
  \vspace{8pt}
  \begin{tcolorbox}[enhanced,colback=poporange,colframe=popink,boxrule=2.8pt,arc=13pt,
      drop shadow=popink,left=18pt,right=18pt,top=12pt,bottom=13pt,after skip=10pt]
    \poppill{#2 \textbullet\ #3}{popink}{popcream}\hspace{5pt}\poppill{#4}{white}{popink}\par\vspace{8pt}
    {\popdisplay\fontsize{14}{14}\selectfont\color{popink}LEÇON #5}\par\vspace{4pt}
    {\raggedright\hyphenpenalty=10000\exhyphenpenalty=10000\popdisplay\fontsize{25}{27}\selectfont\color{popcream}\MakeUppercase{#1}\par}
    \vspace{8pt}
    {\popxbold\fontsize{9.5}{9.5}\selectfont\color{popink}\MakeUppercase{Durée conseillée : #6}}
  \end{tcolorbox}
  \begin{tcolorbox}[enhanced,colback=white,colframe=popink,boxrule=2.2pt,arc=10pt,
      drop shadow=popink,left=16pt,right=16pt,top=10pt,bottom=10pt,after skip=8pt]
    \poppill{Dans cette leçon}{poppurple}{white}\par\vspace{5pt}
    \popsemi\fontsize{9.3}{12.6}\selectfont\color{popink}#7
  \end{tcolorbox}
  \noindent\begin{minipage}[t]{0.487\linewidth}
    \begin{tcolorbox}[enhanced,colback=popgreen,colframe=popink,boxrule=2.2pt,arc=10pt,
        drop shadow=popink,left=11pt,right=11pt,top=10pt,bottom=10pt,height=3.1cm,valign=center]
      \tcbox[colback=white,colframe=popink,boxrule=1.3pt,arc=5pt,boxsep=0pt,nobeforeafter,
        left=3pt,right=3pt,top=3pt,bottom=3pt,valign=center]{\qrcode[height=1.9cm]{https://play.google.com/store/apps/details?id=com.luvvix.onbuch&hl=fr}}%
      \hspace{8pt}\begin{minipage}[c]{0.5\linewidth}
        {\popdisplay\fontsize{11}{12.5}\selectfont\color{white}\MakeUppercase{Télécharge OnBuch}}\par\vspace{4pt}
        {\raggedright\popsemi\fontsize{8.4}{11}\selectfont\color{white}Gratuit \textbullet\ Google Play\\Tuteur IA, annales, concours, résultats\par}
      \end{minipage}
    \end{tcolorbox}
  \end{minipage}\hfill
  \begin{minipage}[t]{0.487\linewidth}
    \begin{tcolorbox}[enhanced,colback=poppurple,colframe=popink,boxrule=2.2pt,arc=10pt,
        drop shadow=popink,left=11pt,right=11pt,top=10pt,bottom=10pt,height=3.1cm,valign=center]
      \tikz[baseline=-0.7ex]\node[circle,fill=white,draw=popink,line width=1.3pt,minimum size=1.6cm,inner sep=0]{\popdisplay\fontsize{24}{24}\selectfont\color{poppurple}+};%
      \hspace{8pt}\begin{minipage}[c]{0.6\linewidth}
        {\popdisplay\fontsize{11}{12.5}\selectfont\color{white}\MakeUppercase{Passe à OnBuch+}}\par\vspace{4pt}
        {\raggedright\popsemi\fontsize{8.4}{11}\selectfont\color{white}Tous les cours signés, Tuteur IA illimité, fiches PDF — onglet OnBuch+ de l'app\par}
      \end{minipage}
    \end{tcolorbox}
  \end{minipage}
  \vspace{8pt}
  \popcontenu
  \vfill
  \signatureonbuch
  \restoregeometry
  \newpage
}

% Rangée « Ce cours contient » (page de garde)
\newcommand{\poptile}[3]{% #1 couleur #2 gros texte #3 légende
  \begin{tcolorbox}[enhanced,colback=#1,colframe=popink,boxrule=2pt,arc=9pt,shadow={2.5pt}{-2.5pt}{0pt}{popink},
      left=6pt,right=6pt,top=9pt,bottom=9pt,halign=center,valign=center,height=1.75cm,width=\linewidth,nobeforeafter]
    {\popdisplay\fontsize{12.5}{13}\selectfont\color{white}\MakeUppercase{#2}}\par\vspace{3pt}
    {\centering\popsemi\fontsize{7.8}{9.5}\selectfont\color{white}#3\par}
  \end{tcolorbox}}
\newcommand{\popcontenu}{%
  \par\noindent{\popxboldls\fontsize{7.5}{7.5}\selectfont\color{popmuted}CE COURS SIGNÉ CONTIENT}\par\vspace{7pt}
  \noindent\begin{minipage}[t]{0.235\linewidth}\poptile{popblue}{Cours}{complet, illustré, pas à pas}\end{minipage}\hfill
  \begin{minipage}[t]{0.235\linewidth}\poptile{poporange}{Méthodes}{et exercices résolus}\end{minipage}\hfill
  \begin{minipage}[t]{0.235\linewidth}\poptile{poppink}{Exercices}{type examen + corrigés}\end{minipage}\hfill
  \begin{minipage}[t]{0.235\linewidth}\poptile{popgreen}{Fiche bilan}{l'essentiel à retenir}\end{minipage}\par}

% Sommaire
\newtcolorbox{sommaire}{enhanced,colback=white,colframe=popink,boxrule=2.2pt,arc=10pt,
  drop shadow=popink,left=18pt,right=18pt,top=13pt,bottom=13pt,before skip=6pt,after skip=20pt,
  before upper={\poppill{Sommaire}{poppurple}{white}\par\vspace{9pt}\popsemi\fontsize{10.6}{14.5}\selectfont\color{popink}}}

% Signature de fin de cours — poussée tout en bas de la dernière page
\newcommand{\findecours}{%
  \par\vfill\nopagebreak
  \begin{center}\popsquiggle{poporange}\end{center}\vspace{4pt}
  \signatureonbuch
}

%% FIGFIX-BEGIN v5 — correctifs globaux des figures (docs/onbuch-generation/tools/figfix)
\usepackage{adjustbox}\usepackage{etoolbox}\tcbuselibrary{raster}
% toute figure TikZ/pgfplots plus large que la ligne est réduite à la largeur disponible
\BeforeBeginEnvironment{tikzpicture}{\begin{adjustbox}{max width=\linewidth}}
\AfterEndEnvironment{tikzpicture}{\end{adjustbox}}
% légendes pgfplots lisibles par-dessus les courbes
\pgfplotsset{every axis/.append style={legend style={fill=white,fill opacity=0.92,text opacity=1,draw=black!15,inner sep=3pt}}}
% étiquettes d'arêtes (syntaxe edge["..."]) avec fond blanc
\tikzset{every edge quotes/.append style={fill=white,inner sep=1.2pt}}
%% FIGFIX-END
\begin{document}

\pagegarde{Aspect ondulatoire de la lumière : diffraction, interférences et effet Doppler}{Physique}{Tle D}{Module 4 : Onde, matière et transformations nucléaires}{P16}{5 h}{Cette leçon établit le caractère ondulatoire de la lumière à travers deux phénomènes fondamentaux : la diffraction par une fente ou un fil, et les interférences à deux ondes avec le dispositif des fentes de Young. Elle se conclut par une étude qualitative de l'effet Doppler et de ses applications, du radar routier à l'expansion de l'Univers.
\vspace{4pt}
\begin{multicols}{2}\small
\begin{itemize}
\item À la fin de cette leçon, je sais décrire et modéliser le phénomène de diffraction de la lumière par une fente ou un fil
\item À la fin de cette leçon, je sais établir et exploiter la relation θ = λ/a pour l'écart angulaire de diffraction
\item À la fin de cette leçon, je sais montrer la nature ondulatoire de la lumière à partir d'observations expérimentales
\item À la fin de cette leçon, je sais définir la différence de marche et l'ordre d'interférence, et prévoir les franges brillantes et sombres
\item À la fin de cette leçon, je sais établir l'expression de l'interfrange i = λD/a et l'exploiter pour mesurer une longueur d'onde
\item À la fin de cette leçon, je sais exploiter la coïncidence des franges en lumière polychromatique et décrire la figure en lumière blanche
\end{itemize}\end{multicols}}

\begin{sommaire}
\begin{multicols}{2}
\begin{enumerate}
\item La diffraction de la lumière
\item Interférences lumineuses : conditions et différence de marche
\item Interfrange : établissement et exploitation de i = λD/a
\item Interférences en lumière polychromatique et en lumière blanche
\item Effet Doppler : étude qualitative et applications
\item Activité d'intégration
\item Méthodes et exercices résolus
\item Exercices
\item Corrigés des exercices
\item Fiche bilan
\end{enumerate}
\end{multicols}
\end{sommaire}

\begin{objectifs}
\begin{itemize}
\item À la fin de cette leçon, je sais décrire et modéliser le phénomène de diffraction de la lumière par une fente ou un fil
\item À la fin de cette leçon, je sais établir et exploiter la relation θ = λ/a pour l'écart angulaire de diffraction
\item À la fin de cette leçon, je sais montrer la nature ondulatoire de la lumière à partir d'observations expérimentales
\item À la fin de cette leçon, je sais définir la différence de marche et l'ordre d'interférence, et prévoir les franges brillantes et sombres
\item À la fin de cette leçon, je sais établir l'expression de l'interfrange i = λD/a et l'exploiter pour mesurer une longueur d'onde
\item À la fin de cette leçon, je sais exploiter la coïncidence des franges en lumière polychromatique et décrire la figure en lumière blanche
\item À la fin de cette leçon, je sais expliquer qualitativement l'effet Doppler et citer des applications (radar, astrophysique)
\end{itemize}
\end{objectifs}

\begin{prerequis}
\begin{itemize}
\item Ondes mécaniques progressives et ondes sinusoïdales (Première) : période, fréquence, longueur d'onde, célérité
\item Relation fondamentale λ = cT = c/ν pour une onde lumineuse
\item Modèle du rayon lumineux et propagation rectiligne de la lumière (optique géométrique)
\item Notions de trigonométrie : approximation tan θ ≈ θ pour les petits angles (en radians)
\item Spectre de la lumière blanche et dispersion par un prisme (longueurs d'onde du visible : 400 à 800 nm)
\end{itemize}
\end{prerequis}

\newpage

\coursec{La diffraction de la lumière}

Lors d'un concert au Palais des Sports de Yaoundé, on entend la musique même caché derrière un pilier, mais on ne voit pas la scène : le son contourne les obstacles, la lumière semble aller tout droit. Pourtant, en éclairant un fil très fin avec un laser, un élève de Terminale observe sur l'écran une tache centrale large bordée de taches latérales : la lumière aussi se diffracte ! De même, les couleurs irisées d'une flaque d'huile sur la route de Douala et le décalage vers le rouge de la lumière des galaxies lointaines s'expliquent par la nature ondulatoire de la lumière. Comment modéliser ces phénomènes et en tirer des mesures précises ? C'est tout l'objet de cette leçon, et nous commençons par le phénomène fondamental : la \cle{diffraction}.

\courssub{Mise en évidence expérimentale}

\textbf{Intuition d'abord.} Tu sais depuis la classe de Seconde que la lumière se propage en ligne droite dans un milieu homogène : c'est le principe de propagation rectiligne, qui explique les ombres et la formation des images dans une chambre noire. Si la lumière allait \emph{toujours} tout droit, une fente fine éclairée par un laser devrait projeter sur un écran une simple ligne lumineuse de même largeur que la fente. L'expérience montre qu'il n'en est rien.

\begin{experience}[Diffraction d'un faisceau laser par une fente fine]
\textbf{Matériel :} un laser (pointeur rouge, $\lambda = \SI{633}{nm}$), une fente réglable de largeur $a$ (ou un fil très fin), un écran blanc placé à la distance $D$ de la fente, une salle plongée dans la pénombre.

\textbf{Protocole :}
\begin{enumerate}
\item Placer la fente verticale de largeur $a$ sur le trajet du faisceau laser.
\item Placer l'écran perpendiculairement au faisceau, à la distance $D \approx \SI{2}{m}$ de la fente.
\item Observer la figure obtenue, puis réduire progressivement la largeur $a$ de la fente.
\end{enumerate}

\textbf{Observations :}
\begin{itemize}
\item Au lieu d'une ligne fine, on observe une \textbf{tache centrale brillante et large}, étalée \emph{perpendiculairement} à la direction de la fente (horizontalement si la fente est verticale).
\item De part et d'autre, on distingue des \textbf{taches latérales} alternées sombres et claires, deux fois moins larges que la tache centrale et de luminosité rapidement décroissante.
\item Plus la fente est fine ($a$ petit), plus la tache centrale s'élargit : la figure s'étale d'autant plus que l'ouverture est petite.
\end{itemize}

\textbf{Interprétation :} la lumière s'étale au-delà de l'ombre géométrique de la fente. Le principe de propagation rectiligne est mis en défaut : la lumière se comporte comme une \cle{onde}.
\end{experience}

\begin{popfigure}
\begin{center}
\begin{tikzpicture}[scale=1.0]
% Laser
\draw[fill=popdark,rounded corners=2pt] (-4.6,-0.25) rectangle (-3.4,0.25);
\node[below] at (-4,-0.3) {\small laser};
% Faisceau avant la fente
\draw[poporange,very thick] (-3.4,0) -- (-1.1,0);
\draw[poporange,thick,->] (-2.6,0.08) -- (-2.2,0.08);
% Fente (obstacle avec ouverture)
\draw[fill=popink] (-1.1,0.25) rectangle (-0.9,1.5);
\draw[fill=popink] (-1.1,-0.25) rectangle (-0.9,-1.5);
\node[above] at (-1,1.55) {\small fente};
% Cote a
\draw[<->,popblue,thick] (-0.75,-0.25) -- (-0.75,0.25) node[midway,right] {\small $a$};
% Rayons diffractés
\draw[poporange,thick] (-1,0) -- (4.2,1.3);
\draw[poporange,thick] (-1,0) -- (4.2,-1.3);
\draw[poporange,thick,dashed] (-1,0) -- (4.2,0);
% Angle theta
\draw[popblue,thick] (1.4,0) arc (0:17.2:2.4);
\node[popblue] at (1.9,0.42) {\small $\theta$};
% Ecran
\draw[fill=popcream,draw=popdark,thick] (4.2,-2.2) rectangle (4.45,2.2);
\node[above] at (4.3,2.25) {\small écran};
% Tache centrale
\draw[fill=poporange!70] (4.2,-1.3) rectangle (4.45,1.3);
\draw[fill=poporange!30] (4.2,1.3) rectangle (4.45,1.95);
\draw[fill=poporange!30] (4.2,-1.95) rectangle (4.2,-1.3);
\draw[fill=poporange!30] (4.2,-1.3) rectangle (4.45,-1.95);
% Cote L
\draw[<->,popblue,thick] (4.7,-1.3) -- (4.7,1.3) node[midway,right] {\small $L$};
% Cote D
\draw[<->,popdark,thick] (-1,-2.0) -- (4.2,-2.0) node[midway,below] {\small $D$};
\end{tikzpicture}
\end{center}
\legende{Montage de la diffraction : le faisceau laser traverse une fente de largeur $a$ ; la tache centrale de largeur $L$ est observée sur un écran à la distance $D$. L'écart angulaire $\theta$ est le demi-angle sous lequel on voit la tache centrale depuis la fente.}
\end{popfigure}

\courssub{Définition et interprétation ondulatoire}

\begin{definition}[Diffraction d'une onde]
La \cle{diffraction} est l'étalement d'une onde au-delà de l'ouverture ou de l'obstacle qu'elle rencontre, lorsque les dimensions de cette ouverture ou de cet obstacle sont de l'\textbf{ordre de grandeur de la longueur d'onde} $\lambda$ (ou plus petites). L'onde ne se propage plus en ligne droite : elle « contourne » l'obstacle ou s'épanouit après l'ouverture.
\end{definition}

Pourquoi ce phénomène ? Le \textbf{principe de Huygens} (1678) fournit l'explication qualitative : chaque point d'un front d'onde atteint par une onde se comporte comme une source secondaire émettant une ondelette sphérique, et le front d'onde suivant est l'enveloppe de ces ondelettes. Quand l'ouverture est grande devant $\lambda$, les ondelettes se recombinent en un faisceau quasiment rectiligne : on retrouve l'optique géométrique. Mais quand l'ouverture $a$ devient comparable à $\lambda$, seules quelques ondelettes passent, et elles rayonnent dans toutes les directions : l'onde s'étale. C'est exactement ce qui se passe pour le son du concert ($\lambda_{\text{son}} \approx \SI{1}{m}$, comparable à la largeur d'une porte) et ce qui arrive à la lumière ($\lambda \approx \SI{0,5}{\micro m}$) face à une fente de quelques dixièmes de millimètre.

\begin{attention}[La diffraction existe pour toutes les ondes]
La diffraction n'est pas propre à la lumière : les vagues à l'entrée du port de Kribi s'étalent derrière la digue, les ondes radio contournent les collines. Mais elle ne devient \emph{observable} que si $a \lesssim \lambda$ environ. Comme $\lambda_{\text{lumière}}$ est de l'ordre du micromètre, il faut des fentes ou des fils extrêmement fins pour la voir : voilà pourquoi la lumière semble « aller tout droit » dans la vie quotidienne.
\end{attention}

\courssub{Écart angulaire et largeur de la tache centrale}

\begin{propriete}[Écart angulaire de la tache centrale (admise)]
Pour une fente de largeur $a$ éclairée en lumière monochromatique de longueur d'onde $\lambda$ (dans l'air, assimilé au vide), l'écart angulaire (demi-angle) de la tache centrale de diffraction est :
\[ \theta = \frac{\lambda}{a} \]
avec $\theta$ en \textbf{radians}, $\lambda$ et $a$ dans la \textbf{même unité} (le mètre en SI). Cette relation est admise ; elle se justifie qualitativement par le principe de Huygens : la première extinction (bord de la tache centrale) apparaît dans la direction où les ondelettes issues des deux bords de la fente sont en opposition de phase, ce qui conduit à $\sin\theta = \lambda/a \approx \theta$ aux petits angles.
\end{propriete}

Vérifions l'homogénéité : $\lambda$ et $a$ sont des longueurs, donc $\lambda/a$ est sans dimension, comme un angle en radians. La relation est cohérente.

\textbf{Largeur de la tache centrale sur l'écran.} Plaçons-nous dans le triangle formé par la fente, le centre de l'écran et le bord de la tache centrale. Si $D$ est la distance fente--écran, le demi-côté de la tache vaut $D\tan\theta$. Aux petits angles, $\tan\theta \approx \theta$, donc :
\[ \frac{L}{2} = D\tan\theta \approx D\,\theta = D\,\frac{\lambda}{a} \quad\Longrightarrow\quad L \approx \frac{2\lambda D}{a} \]

\begin{aretenir}[Les deux relations clés]
\begin{itemize}
\item Écart angulaire (demi-angle) : $\theta = \dfrac{\lambda}{a}$ (en rad).
\item Largeur totale de la tache centrale sur l'écran : $L \approx \dfrac{2\lambda D}{a}$ (aux petits angles).
\end{itemize}
La tache centrale est deux fois plus large que chaque tache latérale.
\end{aretenir}

\begin{popfigure}
\begin{center}
\begin{tikzpicture}
\begin{axis}[
width=0.85\linewidth, height=5.2cm,
xlabel={position $x$ sur l'écran (cm)},
ylabel={intensité $I$},
xmin=-4, xmax=4, ymin=0, ymax=1.15,
xtick={-2.5,0,2.5}, xticklabels={$-\frac{L}{2}$,$0$,$+\frac{L}{2}$},
ytick=\empty,
axis lines=left,
every axis x label/.style={at={(ticklabel* cs:1.02)},anchor=west},
]
\addplot[popcurve,domain=-4:4,samples=400,color=popblue,very thick] {(x==0) ? 1 : (sin(deg(3.14159*x/2.5))/(3.14159*x/2.5))^2};
\draw[dashed,popmuted] (axis cs:-2.5,0) -- (axis cs:-2.5,0.05);
\draw[dashed,popmuted] (axis cs:2.5,0) -- (axis cs:2.5,0.05);
\node[popblue] at (axis cs:0,1.08) {\small tache centrale};
\node[popmuted] at (axis cs:3.3,0.12) {\small taches latérales};
\end{axis}
\end{tikzpicture}
\end{center}
\legende{Répartition de l'intensité lumineuse sur l'écran : la tache centrale concentre l'essentiel de la lumière ; les taches latérales, deux fois moins larges, ont une intensité qui décroît rapidement.}
\end{popfigure}

\textbf{Influence des paramètres.} La relation $\theta = \lambda/a$ se lit comme un tableau de proportionnalité :
\begin{itemize}
\item \textbf{Largeur de la fente :} $\theta$ est inversement proportionnel à $a$. Plus la fente est fine, plus la figure est étalée. Si $a \gg \lambda$, alors $\theta \to 0$ : la diffraction devient invisible et on retrouve la propagation rectiligne.
\item \textbf{Longueur d'onde :} $\theta$ est proportionnel à $\lambda$. Le rouge ($\lambda \approx \SI{650}{nm}$) diffracte \emph{plus} que le bleu ($\lambda \approx \SI{450}{nm}$) : en lumière blanche, le bord de la tache centrale est irisé de rouge.
\item \textbf{Distance $D$ :} elle ne change pas l'angle $\theta$ mais agrandit la figure sur l'écran ($L$ proportionnel à $D$), ce qui facilite les mesures.
\end{itemize}

\begin{exemplebox}[Laser rouge et fente fine]
Un laser rouge de longueur d'onde $\lambda = \SI{633}{nm}$ éclaire une fente de largeur $a = \SI{0,10}{mm}$. L'écran est placé à $D = \SI{2,0}{m}$.
\begin{align*}
\theta &= \frac{\lambda}{a} = \frac{633\times 10^{-9}}{1,0\times 10^{-4}} = 6,33\times 10^{-3}\ \text{rad} \approx 0,36° \\
L &\approx \frac{2\lambda D}{a} = \frac{2\times 633\times 10^{-9}\times 2,0}{1,0\times 10^{-4}} = 2,53\times 10^{-2}\ \text{m} \approx \SI{2,5}{cm}
\end{align*}
Une fente d'un dixième de millimètre produit donc une tache de $2,5$ cm sur l'écran : l'étalement est $250$ fois plus grand que la fente ! La diffraction est un formidable « amplificateur » de mesure.
\end{exemplebox}

\begin{attention}[Les deux pièges classiques du Bac]
\begin{enumerate}
\item $\theta = \lambda/a$ donne le \textbf{demi-écart angulaire} : la tache centrale a une largeur angulaire \emph{totale} de $2\theta$, et sa largeur sur l'écran est $L = 2\lambda D/a$, pas $\lambda D/a$.
\item \textbf{Convertis toutes les longueurs en mètres} avant le calcul : $\lambda$ en nanomètres ($\SI{1}{nm} = 10^{-9}\ \text{m}$), $a$ souvent en millimètres. Un mélange nm/mm donne un résultat faux de plusieurs ordres de grandeur.
\end{enumerate}
\end{attention}

\courssub{Diffraction par un fil et mesure d'une longueur d'onde}

\begin{propriete}[Diffraction par un fil — théorème de Babinet (admis)]
Un fil de diamètre $a$ éclairé par une onde monochromatique donne, hors de l'image géométrique centrale, une figure de diffraction \textbf{identique} à celle d'une fente de même largeur $a$ : mêmes positions des extinctions, même écart angulaire $\theta = \lambda/a$. Ce résultat (théorème de Babinet) est admis au Bac.
\end{propriete}

C'est très pratique : un cheveu, un fil de pêche ou le fil d'une ampoule suffisent pour observer la diffraction, et réciproquement, la mesure de $L$ permet de déterminer le diamètre d'un fil (contrôle qualité industriel).

\begin{methode}[Déterminer une longueur d'onde inconnue par diffraction]
\begin{enumerate}
\item Mesurer la largeur $a$ de la fente (ou du fil) et la distance $D$ fente--écran, en \textbf{mètres}.
\item Mesurer la largeur $L$ de la tache centrale sur l'écran, du bord sombre au bord sombre opposé.
\item Appliquer $L = \dfrac{2\lambda D}{a}$, d'où $\lambda = \dfrac{L\,a}{2D}$.
\item Vérifier l'ordre de grandeur : une longueur d'onde visible est comprise entre $400$ et $\SI{800}{nm}$.
\end{enumerate}
\end{methode}

\begin{exoresolu}[Mesure du diamètre d'un fil de pêche]
Un fil de pêche de diamètre $a$ inconnu est éclairé par un laser vert ($\lambda = \SI{532}{nm}$). Sur un écran placé à $D = \SI{1,50}{m}$, la tache centrale mesure $L = \SI{3,2}{cm}$. Déterminer le diamètre du fil.
\tcblower
D'après le théorème de Babinet, le fil donne la même figure qu'une fente de même largeur, donc :
\[ L = \frac{2\lambda D}{a} \quad\Longrightarrow\quad a = \frac{2\lambda D}{L} = \frac{2\times 532\times 10^{-9}\times 1,50}{3,2\times 10^{-2}} = 4,99\times 10^{-5}\ \text{m} \]
Le diamètre du fil est $a \approx \SI{50}{\micro m}$ (soit $0,050$ mm), valeur typique d'un fil de pêche fin. Toutes les grandeurs ont été converties en mètres avant le calcul.
\end{exoresolu}

\begin{savaistu}[Pourquoi le tamis ne diffracte pas la lumière, mais le CD si]
C'est parce que les pores d'un tamis à farine (environ $\SI{100}{\micro m}$) sont bien plus grands que la longueur d'onde de la lumière ($\approx \SI{0,5}{\micro m}$) que la lumière n'y est pas diffractée de façon visible : le rapport $\lambda/a$ y est minuscule. En revanche, les rainures d'un CD, espacées d'environ $\SI{1,6}{\micro m}$, sont de l'ordre de grandeur de $\lambda$ : chaque rainure diffracte la lumière, et comme $\theta$ dépend de $\lambda$, chaque couleur est déviée différemment — d'où les arcs-en-ciel que tu observes en inclinant le disque. Le même principe explique les irisations des ailes de certains papillons de la forêt camerounaise, comme le \emph{Morpho} !
\end{savaistu}

\courssub{Conclusion : une signature de la nature ondulatoire}

La diffraction est une \textbf{signature irréfutable de la nature ondulatoire de la lumière} : un flux de particules se propageant en ligne droite ne pourrait jamais s'étaler derrière une fente ni produire des zones sombres là où l'optique géométrique prévoit de la lumière. La lumière se comporte comme une onde électromagnétique de très petite longueur d'onde (de $\SI{400}{nm}$ pour le violet à $\SI{800}{nm}$ pour le rouge), ce qui explique pourquoi sa diffraction n'apparaît qu'avec des objets micrométriques. Dans la suite de la leçon, une seconde expérience — les \emph{interférences} des fentes de Young — confirmera cette nature ondulatoire et nous donnera un outil de mesure encore plus précis.

\coursec{Interférences lumineuses : conditions et différence de marche}

Nous venons de voir que la lumière, en passant par une fente, ne se propage pas en ligne droite stricte mais s'épanouit : c'est la \cle{diffraction}. Ce phénomène, loin d'être une curiosité, est la clé qui permet de révéler l'interférence. En effet, pour observer des interférences lumineuses, il nous faut deux ondes \cle{cohérentes} issues de la même source primaire. La diffraction par deux fentes rapprochées (le dispositif de Young) réalise exactement cette condition : chaque fente devient une source secondaire, synchronisée avec l'autre, et c'est la superposition de leurs ondes diffractées qui dessine les franges sur l'écran.

\courssub{Des sources cohérentes : la condition sine qua non}

Rappelons-nous : deux sources sont dites \cle{cohérentes} si elles émettent des ondes de \textbf{même fréquence} (donc même longueur d'onde $\lambda$ dans le même milieu) et si leur \textbf{différence de phase reste constante dans le temps} (elles sont \emph{synchrones}). Deux lampes à incandescence indépendantes, même identiques, ne sont \textbf{jamais} cohérentes : leurs atomes émettent de façon aléatoire et indépendante, la phase relative fluctue en $\sim \SI{e-8}{s}$, et l'œil ou le capteur ne voit qu'une illumination uniforme, moyenne temporelle de l'interférence fugace.

\begin{attention}[Piège classique : deux lampes $\neq$ interférences]
Ne confonds pas \emph{monochromatique} (une seule $\lambda$) et \emph{cohérente} (phase constante). Deux lasers He–Ne distincts sont monochromatiques mais \textbf{non cohérents} entre eux (leurs cavités ne sont pas verrouillées en phase). En pratique, on obtient la cohérence spatiale en faisant éclater \textbf{une seule source primaire} sur deux ouvertures (fentes de Young, fente de Fresnel, miroir de Lloyd, lame d'air, etc.).
\end{attention}

\courssub{Le dispositif des fentes de Young : géométrie et sources secondaires}

\begin{popfigure}
\begin{tikzpicture}[scale=1.1, >=Stealth]
% Paramètres
\def\D{6}   % distance fentes-écran (cm)
\def\a{1.2} % séparation fentes (cm)
\def\h{2.5} % hauteur fentes/écran
\def\xM{1.5}% abscisse point M sur écran
% Source primaire S
\coordinate (S) at (-2,0);
\fill[popblue] (S) circle (2.5pt) node[left=2pt] {$S$};
% Fentes S1, S2
\coordinate (S1) at (0, \a/2);
\coordinate (S2) at (0, -\a/2);
\draw[thick, popdark] (-0.15,\h) -- (-0.15,-\h) node[midway, left=3pt] {Plan des fentes};
\fill[poporange] (S1) circle (2pt) node[left=2pt] {$S_1$};
\fill[poporange] (S2) circle (2pt) node[left=2pt] {$S_2$};
% Flèche a
\draw[<->, poppurple] (0.3,\a/2) -- (0.3,-\a/2) node[midway, right] {$a$};
% Écran
\draw[thick, popgreen] (\D,\h) -- (\D,-\h) node[midway, right=3pt] {Écran};
% Point M
\coordinate (M) at (\D, \xM);
\fill[poppink] (M) circle (2pt) node[right=2pt] {$M(x)$};
% Trajets S1M, S2M
\draw[popblue, dashed] (S1) -- (M) node[midway, above left, sloped] {$S_1M$};
\draw[popblue, dashed] (S2) -- (M) node[midway, below left, sloped] {$S_2M$};
% Distance D
\draw[<->, popdark] (0,-2.8) -- (\D,-2.8) node[midway, below] {$D$};
% Source vers fentes
\draw[popblue, ->] (S) -- (S1) node[midway, above, sloped] {Onde incidente};
\draw[popblue, ->] (S) -- (S2);
% Cote x
\draw[<->, popdark] (\D+0.3,0) -- (\D+0.3,\xM) node[midway, right] {$x$};
% Angle theta
\draw[poppurple, ->] (0,0) ++(0.8,0) arc (0:14:0.8) node[right=2pt] {$\theta$};
\end{tikzpicture}
\legende{Schéma géométrique du dispositif de Young. Une source ponctuelle $S$ (ou une fente éclairée par un laser) éclaire deux fentes fines $S_1$ et $S_2$, parallèles, distantes de $a$. L'écran est à une distance $D \gg a$. En un point $M$ d'abscisse $x$, les ondes issues de $S_1$ et $S_2$ se superposent.}
\end{popfigure}

La source primaire $S$ (souvent une fente fine éclairée par un laser ou une lampe à vapeur de sodium + filtre interférentiel) émet une onde sphérique (ou cylindrique si fente). Les deux fentes $S_1$ et $S_2$, très fines ($l \ll a$), sélectionnent chacune une portion de ce front d'onde. D'après le \cle{principe de Huygens–Fresnel}, chaque fente se comporte comme une \textbf{source secondaire} : elles émettent des ondes de même fréquence, de même phase initiale (synchrones) et de même amplitude (si l'éclairage est symétrique). Elles sont donc \textbf{cohérentes}.

La distance $a = S_1S_2$ est de l'ordre du millimètre (typiquement $\SI{0,2}{mm}$ à $\SI{1}{mm}$), tandis que la distance fentes–écran $D$ est de l'ordre du mètre ($\SI{1}{m}$ à $\SI{2}{m}$). On a donc $D \gg a$, ce qui justifie les approximations de \emph{petits angles} que nous utiliserons.

\courssub{Zone d'interférences et aspect des franges}

Les ondes diffractées par $S_1$ et $S_2$ s'étalent dans l'espace. La \cle{zone d'interférences} est la région où les deux faisceaux se superposent significativement. Comme les fentes sont parallèles (direction $Oz$), l'interférence est invariante par translation selon $z$ : les franges sont des \textbf{bandes rectilignes, parallèles aux fentes}.

\begin{popfigure}
\begin{tikzpicture}[scale=1.2]
% Franges : 9 bandes (4 sombres, 5 brillantes centrées sur k=0)
\def\frangeW{0.6} % largeur d'une frange (interfrange)
\foreach \k in {-4,-3,...,4} {
  \pgfmathsetmacro{\y}{\k*\frangeW};
  \pgfmathsetmacro{\ynext}{\y+\frangeW};
  \ifodd\numexpr\k+4\relax
    % Frange brillante
    \fill[popgold!80] (-3,\y) rectangle (3,\ynext);
    \node at (0, \y+0.5*\frangeW) [font=\bfseries\small, text=popdark] {$k=\k$};
  \else
    % Frange sombre
    \fill[popdark!90] (-3,\y) rectangle (3,\ynext);
    \node at (0, \y+0.5*\frangeW) [font=\bfseries\small, text=white] {$k=\k$};
  \fi
}
% Légendes
\draw[<->, thick, popblue] (-3.5,-2.4) -- (3.5,-2.4) node[midway, below] {Direction $Ox$ (perpendiculaire aux fentes)};
\draw[->, thick, poppurple] (0,-2.8) -- (0,-1.8) node[left] {$Oz \parallel$ fentes};
\node[align=center, font=\small] at (0,-3.3) {Frange centrale brillante ($k=0$) : $\delta=0$ \\ Franges brillantes : $k \in \mathbb{Z}$ \quad Franges sombres : $k+\frac12$};
\end{tikzpicture}
\legende{Figure d'interférences de Young (lumière monochromatique). Franges rectilignes, équidistantes, alternativement brillantes (claires) et sombres (sombres). L'ordre d'interférence $k$ est indiqué. La frange centrale ($k=0$) est toujours brillante car $\delta=0$ en $O$.}
\end{popfigure}

Sur un écran placé dans cette zone, on observe une alternance de franges \textbf{brillantes} (intensité maximale) et \textbf{sombres} (intensité minimale, idéalement nulle si les amplitudes sont égales). Elles sont \textbf{équidistantes} (dans l'approximation $D \gg a$ et petits angles) et \textbf{non localisées} : on les voit sur n'importe quel écran placé dans la zone commune, contrairement aux franges localisées (ex. lames minces).

\courssub{Différence de marche et ordre d'interférence}

Considérons un point $M$ sur l'écran, d'abscisse $x$ par rapport à l'axe central $O$ (milieu de $S_1S_2$). Les ondes issues de $S_1$ et $S_2$ parcourent les distances $S_1M$ et $S_2M$ pour atteindre $M$.

\begin{definition}[Différence de marche $\delta$]
La \cle{différence de marche} en un point $M$ est la différence algébrique des distances optiques parcourues par les deux ondes depuis les sources secondaires jusqu'à $M$ :
\[
\delta(M) = S_2M - S_1M
\]
C'est une \textbf{longueur} (unité SI : mètre $\text{m}$). Elle est positive si $S_2M > S_1M$ (point $M$ du côté de $S_1$), négative dans le cas inverse.
\end{definition}

Dans le triangle $S_1S_2M$, comme $D \gg a$ et $x \ll D$, les droites $(S_1M)$ et $(S_2M)$ sont quasi parallèles. La différence de marche s'approxime alors par la projection de $S_1S_2$ sur la direction d'arrivée :
\[
\delta \approx a \sin\theta \approx a \tan\theta = a\,\frac{x}{D}
\]
où $\theta$ est l'angle que fait la direction $OM$ avec l'axe central. Cette relation $\delta = a\,x/D$ est fondamentale pour établir l'interfrange (section 3).

\begin{definition}[Ordre d'interférence $p$]
L'\cle{ordre d'interférence} en $M$ est le nombre sans unité défini par :
\[
p(M) = \frac{\delta(M)}{\lambda}
\]
Il caractérise l'état vibratoire en $M$ : c'est le nombre de longueurs d'onde de retard (ou d'avance) d'une onde par rapport à l'autre.
\end{definition}

\courssub{Conditions d'interférences constructives et destructives}

Lorsque deux ondes sinusoïdales de même fréquence, même amplitude $A$ et de différence de phase $\Delta\varphi = \frac{2\pi}{\lambda}\delta$ se superposent, l'amplitude résultante est $A_R = 2A\left|\cos(\Delta\varphi/2)\right|$. L'intensité (proportionnelle au carré de l'amplitude) est maximale pour $\cos(\Delta\varphi/2) = \pm 1$ et nulle pour $\cos(\Delta\varphi/2) = 0$.

\begin{propriete}[Conditions d'interférences en lumière monochromatique (admise au Bac)]
En un point $M$ de la zone d'interférences, on a :
\begin{itemize}
  \item \textbf{Interférence constructive (frange brillante)} $\iff$ les ondes sont en phase $\iff \Delta\varphi = 2k\pi$ $\iff$ $\boxed{\delta = k\lambda}$ avec $k \in \mathbb{Z}$.
  \item \textbf{Interférence destructive (frange sombre)} $\iff$ les ondes sont en opposition de phase $\iff \Delta\varphi = (2k+1)\pi$ $\iff$ $\boxed{\delta = \left(k+\frac12\right)\lambda}$ avec $k \in \mathbb{Z}$.
\end{itemize}
L'entier $k$ est l'ordre de la frange brillante ; pour une frange sombre, l'ordre est le demi-entier $k+\frac12$.
\end{propriete}

La \textbf{frange centrale} ($M \equiv O$, $x=0$) correspond à $\delta = 0$, donc $k=0$ : elle est \textbf{toujours brillante} (ordre 0). Les franges brillantes d'ordres $k = \pm 1, \pm 2, \dots$ sont symétriques par rapport à l'axe central.

\begin{attention}[Ordre $p$ vs différence de marche $\delta$ : confusion fréquente]
\begin{itemize}
  \item $\delta$ est une \textbf{longueur} (en $\text{m}$, $\text{nm}$, etc.). $\delta = \SI{600}{nm}$.
  \item $p = \delta/\lambda$ est un \textbf{nombre pur} (sans unité). Si $\lambda = \SI{600}{nm}$, alors $p = 1$.
  \item Au Bac, on te demande souvent : « Déterminer la nature de la frange en $M$ ». La méthode rigoureuse : calculer $\delta$, puis $p = \delta/\lambda$.
  \begin{itemize}
    \item Si $p$ est un \textbf{entier} ($\dots, -2, -1, 0, 1, 2, \dots$) $\rightarrow$ \textbf{frange brillante} d'ordre $p$.
    \item Si $p$ est un \textbf{demi-entier} ($\dots, -3/2, -1/2, 1/2, 3/2, \dots$) $\rightarrow$ \textbf{frange sombre} d'ordre $p$.
    \item Sinon $\rightarrow$ \textbf{état intermédiaire} (intensité entre $0$ et $I_{\max}$).
  \end{itemize}
  \textbf{Ne conclus jamais « brillante » ou « sombre » sans avoir vérifié que $p$ est bien entier ou demi-entier.}
\end{itemize}
\end{attention}

\begin{methode}[Déterminer la nature d'une frange en un point $M$]
\begin{enumerate}
  \item Identifier la géométrie : $a$ (distance fentes), $D$ (distance fentes-écran), $\lambda$ (longueur d'onde dans le milieu), $x$ (abscisse de $M$ sur l'écran).
  \item Calculer la différence de marche (approximation petits angles) : $\delta \approx a\,\dfrac{x}{D}$.
  \item Calculer l'ordre d'interférence : $p = \dfrac{\delta}{\lambda} = \dfrac{a\,x}{\lambda D}$.
  \item Analyser $p$ :
  \begin{itemize}
    \item $p \in \mathbb{Z}$ $\rightarrow$ Frange brillante d'ordre $p$.
    \item $p \in \mathbb{Z} + \frac12$ $\rightarrow$ Frange sombre d'ordre $p$.
    \item Autre $\rightarrow$ Frange intermédiaire (intensité partielle).
  \end{itemize}
  \item Conclure par une phrase complète : « En $M$, $\delta = \dots$, $p = \dots$, c'est donc une frange \emph{brillante/sombre/intermédiaire} d'ordre $\dots$. »
\end{enumerate}
\end{methode}

\begin{exemplebox}[Application numérique : nature d'une frange]
Dans un montage de Young, on utilise un laser He–Ne ($\lambda_0 = \num{632,8}\,\text{nm}$ dans le vide, $n_{\text{air}} \approx 1$). Les fentes sont distantes de $a = \SI{0,50}{mm}$ et l'écran est à $D = \SI{1,20}{m}$. On observe un point $M$ d'abscisse $x = \SI{3,0}{mm}$.
\medskip
\textbf{Solution.}
\begin{align*}
\delta &= a\,\frac{x}{D} = 0.50 \times 10^{-3} \times \frac{3.0 \times 10^{-3}}{1.20} = 1.25 \times 10^{-6}\,\text{m} = \SI{1250}{nm}. \\
p &= \frac{\delta}{\lambda} = \frac{1250}{632.8} \approx 1.975.
\end{align*}
$p \approx \num{1,975}$ n'est ni un entier ($2$), ni un demi-entier ($1,5$ ou $2,5$). C'est une \textbf{frange intermédiaire}, proche de la frange brillante d'ordre $2$ (pour laquelle $x_2 = 2\lambda D/a = \SI{3,04}{mm}$).
\end{exemplebox}

\begin{savaistu}[Thomas Young (1773–1829) : l'expérience décisive]
En 1801, le polymathe anglais Thomas Young présente à la Royal Society son « \emph{On the Theory of Light and Colours} ». Il fait passer la lumière du soleil par une fente, puis par deux fentes fines ($a \approx \SI{0,2}{mm}$) et observe sur un écran des « \emph{fringes of colours} ». Il mesure l'interfrange, en déduit $\lambda \approx \SIrange{500}{600}{nm}$ pour la lumière visible et explique les couleurs des bulles de savon et des plumes de paon par interférences dans des lames minces. Cette expérience, reproductible et quantitative, a été l'argument \textbf{décisif} en faveur de la \cle{nature ondulatoire de la lumière} face au modèle corpusculaire de Newton qui dominait alors. Young a aussi découvert l'astigmatisme, contribué au déchiffrement des hiéroglyphes (pierre de Rosette) et formulé le module d'élasticité (module de Young). Au Cameroun, les télécommunications par fibre optique (MTN, Orange, CAMTEL) reposent sur la cohérence de la lumière laser — une application directe de la cohérence spatiale que Young a mise en évidence.
\end{savaistu}

\begin{exoresolu}[Détermination de la longueur d'onde par les franges de Young]
\textbf{Énoncé.} On réalise l'expérience de Young avec une source laser. Les fentes sont distantes de $a = \SI{0,40}{mm}$. L'écran est placé à $D = \SI{2,00}{m}$. On mesure la distance entre la frange brillante d'ordre $k=5$ et la frange centrale : $x_5 = \SI{7,5}{mm}$. Déterminer la longueur d'onde $\lambda$ du laser.
\tcblower
\textbf{Solution détaillée.}
\begin{enumerate}
  \item Pour une frange brillante d'ordre $k$, la condition constructive donne $\delta = k\lambda$.
  \item Avec l'approximation $\delta \approx a\,x/D$, on a $a\,\dfrac{x_k}{D} = k\lambda$.
  \item Donc $\lambda = \dfrac{a\,x_k}{k\,D}$.
  \item Application numérique :
  \begin{align*}
  \lambda &= \frac{0.40 \times 10^{-3} \times 7.5 \times 10^{-3}}{5 \times 2.00} \\
  &= \frac{3.0 \times 10^{-6}}{10} = 3.0 \times 10^{-7}\,\text{m} = \SI{300}{nm}.
  \end{align*}
  \item \textbf{Réponse :} $\lambda = \SI{300}{nm}$ (ultraviolet proche — ce laser hypothétique n'est pas un He–Ne standard !).
\end{enumerate}
\textbf{Remarque.} En pratique, on mesure l'interfrange $i = x_{k+1} - x_k = \lambda D/a$ (section 3), ce qui moyenne les erreurs sur plusieurs franges.
\end{exoresolu}

\begin{exercice}[Vérification des ordres d'interférence]
On éclaire deux fentes distantes de $a = \SI{0,60}{mm}$ par une lumière monochromatique $\lambda = \SI{500}{nm}$. L'écran est à $D = \SI{1,50}{m}$.
\begin{enumerate}
  \item Calculer l'abscisse $x_3$ de la frange brillante d'ordre $k=3$.
  \item Calculer l'abscisse $x'_{2}$ de la frange sombre d'ordre $k=2$ (c'est-à-dire $p = 2,5$).
  \item Vérifier que la distance entre ces deux franges vaut bien la moitié de l'interfrange.
\end{enumerate}
\end{exercice}

\begin{corrige}[Exercice : Vérification des ordres d'interférence]
\begin{enumerate}
  \item Frange brillante $k=3$ : $\delta = 3\lambda = a\,x_3/D \Rightarrow x_3 = 3\lambda D/a$.
  $x_3 = 3 \times 500 \times 10^{-9} \times 1.50 / (0.60 \times 10^{-3}) = 3.75 \times 10^{-3}\,\text{m} = \SI{3,75}{mm}$.
  \item Frange sombre d'ordre $p = 2,5$ : $\delta = 2.5\lambda = a\,x'_2/D \Rightarrow x'_2 = 2.5\lambda D/a$.
  $x'_2 = 2.5 \times 500 \times 10^{-9} \times 1.50 / (0.60 \times 10^{-3}) = 3.125 \times 10^{-3}\,\text{m} = \SI{3,125}{mm}$.
  \item Distance $x_3 - x'_2 = (3 - 2.5)\lambda D/a = 0.5 \lambda D/a = i/2$. CQFD.
\end{enumerate}
\end{corrige}

\coursec{Interfrange : établissement et exploitation de i = λD/a}

Après avoir mis en évidence les conditions d'interférences constructives et destructives dans le dispositif des fentes de Young (section précédente), nous allons maintenant quantifier la géométrie des franges observées sur l'écran. La grandeur clé est l'\cle{interfrange}, notée $i$, qui caractérise l'espacement régulier entre les franges brillantes (ou sombres). Son expression simple, $i = \lambda D / a$, est au cœur de nombreuses applications expérimentales, de la mesure de longueurs d'onde à la métrologie optique.

\courssub{Définition et repérage des franges sur l'écran}

\begin{definition}[Interfrange]
Dans un système d'interférences à franges droites, parallèles et équidistantes (comme celles de Young), on appelle \textbf{interfrange} $i$ la distance séparant les centres de deux franges consécutives de même nature (deux franges brillantes consécutives ou deux franges sombres consécutives). Elle s'exprime en mètres (\SI{}{m}).
\end{definition}

Rappelons le dispositif : deux fentes $S_1$ et $S_2$, écartées de la distance $a = S_1S_2$, sont éclairées par une source monochromatique de longueur d'onde $\lambda$. L'écran $\mathcal{E}$ est placé à une distance $D$ des fentes, avec $D \gg a$. On repère un point $M$ sur l'écran par son abscisse $x$ par rapport à l'axe central $OO'$ (point $O$ milieu de $S_1S_2$, $O'$ intersection de l'axe avec l'écran).

\begin{popfigure}
\begin{tikzpicture}[scale=1.1, >=Stealth]
% Coordonnées principales
\coordinate (O) at (0,0);
\coordinate (Oprime) at (5,0);
\coordinate (S1) at (0,0.4);
\coordinate (S2) at (0,-0.4);
\coordinate (M) at (5,1.5);
\coordinate (Xaxis) at (1,0); % Pour l'angle theta

% Fentes S1 et S2 (traits verticaux épais)
\draw[very thick, popdark] (-0.1, 0.4) -- (-0.1, -0.4) node[midway, left=2pt] {$S_1$};
\draw[very thick, popdark] (0.1, 0.4) -- (0.1, -0.4) node[midway, right=2pt] {$S_2$};
\draw[popmuted, <->] (0.1, 0.7) -- (-0.1, 0.7) node[midway, above] {$a$};

% Écran
\draw[very thick, popblue] (5, -2) -- (5, 2) node[midway, right=5pt] {$\mathcal{E}$};
\draw[popmuted, <->] (0, -2.3) -- (5, -2.3) node[midway, below] {$D$};

% Axes et points O, O'
\draw[dashed, popmuted] (-1, 0) -- (5.5, 0) node[right] {$Ox$};
\node at (O) [circle, fill, inner sep=1pt, label=below left:$O$] {};
\node at (Oprime) [circle, fill, inner sep=1pt, label=below right:$O'$] {};

% Point M et abscisse x
\draw[poporange, thick] (M) node[right=2pt] {$M$} -- (5,0);
\node at (M) [circle, fill=poporange, inner sep=1.5pt] {};
\draw[popmuted, <->] (5.3, 0) -- (5.3, 1.5) node[midway, right] {$x$};

% Chemins optiques S1M et S2M
\draw[popgreen, thick] (-0.1, 0.4) -- (M);
\draw[poppurple, thick] (0.1, -0.4) -- (M);

% Angle theta en O
\pic[draw, popred, thick, angle radius=1.2cm, angle eccentricity=1.3, "$\theta$"] {angle = Xaxis--O--M};

% Zoom sur le triangle au niveau des fentes (en bas à gauche)
\begin{scope}[shift={(-3, -3)}, scale=3]
% Fentes zoom
\draw[very thick, popdark] (0, 0.2) -- (0, -0.2) node[midway, left=1pt] {$S_1$} node[midway, right=1pt] {$S_2$};
% Rayons parallèles (direction theta)
\pgfmathsetmacro{\angT}{16.7} % approx atan(1.5/5)
\draw[popgreen, thick] (0, 0.2) -- (1.2, {0.2 + 1.2*tan(\angT)});
\draw[poppurple, thick] (0, -0.2) -- (1.2, {-0.2 + 1.2*tan(\angT)});
% Perpendiculaire de S1 sur le rayon S2M (difference de marche delta)
\draw[dashed, popmuted] (0, 0.2) -- (0.15, 0.2); % Horizontal from S1
\draw[dashed, popmuted] (0, -0.2) -- (0, 0.2); % Vertical a
% Delta double fleche
\draw[poporange, thick, <->] (0, -0.2) -- (0, 0.2) node[midway, right=2pt] {$\delta$};
% Angle theta au niveau de S2
\coordinate (S2z) at (0, -0.2);
\coordinate (S1z) at (0, 0.2);
\coordinate (Rayz) at (1, {-0.2 + 1*tan(\angT)});
\coordinate (Vertz) at (0, 0.2); % Vertical up from S2
\pic[draw, popred, angle radius=0.35cm, "$\theta$"] {angle = Vertz--S2z--Rayz};
\node[popmuted, font=\small] at (-0.6, 0) {Zoom};
\end{scope}
\end{tikzpicture}
\legende{Dispositif des fentes de Young : géométrie pour l'établissement de la différence de marche $\delta$ en un point $M$ d'abscisse $x$. Le zoom montre le triangle rectangle au voisinage des fentes justifiant $\delta = a \sin\theta$.}
\end{popfigure}

\courssub{Établissement de la différence de marche $\delta = \frac{ax}{D}$}

Considérons les ondes issues de $S_1$ et $S_2$ arrivant en $M$. La différence de marche est $\delta = S_2M - S_1M$ (en valeur absolue, l'ordre n'importe que pour le signe de l'interférence).
Comme $D \gg a$ et $D \gg x$ (observation au voisinage de l'axe), les droites $(S_1M)$ et $(S_2M)$ sont quasi parallèles. L'angle $\theta$ que fait la direction $OM$ avec l'axe $OO'$ est petit.

Dans le triangle rectangle construit au niveau des fentes (voir zoom sur la figure), l'hypoténuse est l'écart $a = S_1S_2$ et le côté opposé à l'angle $\theta$ est précisément la différence de marche $\delta$. On a donc la relation géométrique exacte (dans l'approximation des rayons parallèles) :
\[
\delta = a \sin\theta.
\]

Pour les petits angles ($\theta$ en radians), $\sin\theta \approx \tan\theta$. Or, dans le triangle rectangle $O O' M$, $\tan\theta = \frac{O'M}{OO'} = \frac{x}{D}$.
On obtient ainsi la relation fondamentale, valable pour $D \gg a$ et $x \ll D$ (petits angles) :
\[
\boxed{\delta = \frac{a x}{D}}
\]

\begin{attention}[Condition de validité : $D \gg a$ et petits angles]
La formule $\delta = ax/D$ repose sur deux approximations successives :
\begin{enumerate}
\item Les rayons $S_1M$ et $S_2M$ sont parallèles (ondes planes) : nécessite $D \gg a$.
\item $\sin\theta \approx \tan\theta \approx \theta$ : nécessite $\theta \ll 1$ rad, c'est-à-dire $x \ll D$.
\end{enumerate}
Au Bac, il faut citer explicitement l'hypothèse $D \gg a$ (souvent notée $D \gg a$ ou $D \gg x$) pour justifier l'écriture $\delta = a \sin\theta$ puis $\delta = ax/D$.
\end{attention}

\courssub{Position des franges et expression de l'interfrange}

La condition d'interférence constructive (frange brillante) d'ordre $k \in \mathbb{Z}$ est $\delta = k \lambda$.
En utilisant $\delta = ax/D$, on trouve l'abscisse $x_k$ de la frange brillante d'ordre $k$ :
\[
\frac{a x_k}{D} = k \lambda \quad \Rightarrow \quad \boxed{x_k = k \frac{\lambda D}{a}}.
\]

La frange centrale ($k=0$) se trouve en $x_0 = 0$ (point $O'$) : elle est toujours brillante (blanche en lumière blanche).

L'interfrange $i$ est la distance entre deux franges brillantes consécutives (ordres $k$ et $k+1$) :
\[
i = x_{k+1} - x_k = (k+1)\frac{\lambda D}{a} - k\frac{\lambda D}{a} = \frac{\lambda D}{a}.
\]

\begin{propriete}[Interfrange des fentes de Young — Démonstration exigible au Bac]
Dans le dispositif des fentes de Young éclairé par une onde monochromatique de longueur d'onde $\lambda$, avec un écart $a$ entre les fentes et une distance fentes-écran $D$ telle que $D \gg a$, les franges brillantes sont équidistantes. L'interfrange $i$ vaut :
\[
\boxed{i = \frac{\lambda D}{a}}
\]
où $i$ est en mètres (\SI{}{m}), $\lambda$ en mètres (\SI{}{m}), $D$ en mètres (\SI{}{m}) et $a$ en mètres (\SI{}{m}).
\end{propriete}

\begin{aretenir}[Analyse dimensionnelle et discussion]
Vérifions l'homogénéité : $[\lambda] = L$, $[D] = L$, $[a] = L$ $\Rightarrow$ $[i] = L/L \cdot L = L$. \textbf{Correct}.
Discussion physique :
\begin{itemize}
\item $i \propto \lambda$ : plus la longueur d'onde est grande (rouge), plus les franges sont écartées.
\item $i \propto D$ : plus l'écran est loin, plus les franges s'écartent (ouverture angulaire constante).
\item $i \propto 1/a$ : plus les fentes sont rapprochées, plus les franges sont larges (diffraction plus forte).
\end{itemize}
\end{aretenir}

\courssub{Position des franges sombres (complément Bac)}

La condition d'interférence destructive est $\delta = (k + \frac{1}{2})\lambda$ avec $k \in \mathbb{Z}$.
L'abscisse des franges sombres est donc :
\[
x'_k = \left(k + \frac{1}{2}\right) \frac{\lambda D}{a}.
\]
On vérifie que les franges sombres sont intercalées exactement au milieu des franges brillantes : $x'_k = \frac{x_k + x_{k+1}}{2}$. L'interfrange entre deux franges sombres consécutives vaut aussi $i = \lambda D/a$.

\begin{attention}[Piège classique : confusion sur le paramètre $a$]
Dans la formule $i = \lambda D / a$, la lettre \textbf{$a$ désigne l'écart entre les deux fentes $S_1$ et $S_2$} (distance $S_1S_2$).
Ne confondez \textbf{JAMAIS} cet $a$ avec la \textbf{largeur d'une fente} (souvent notée $a$ aussi dans la diffraction par une fente unique, où l'écart angulaire du premier minimum est $\theta = \lambda/a$).
\begin{itemize}
\item \textbf{Young (interférences)} : $a = S_1S_2$ (écart entre fentes). $i = \lambda D/a$.
\item \textbf{Fente unique (diffraction)} : $a =$ largeur de la fente. Premier minimum à $\theta \approx \lambda/a$.
\end{itemize}
Relisez toujours l'énoncé pour identifier ce que représente $a$.
\end{attention}

\courssub{Exploitation expérimentale : mesure de $\lambda$ ou $a$}

La formule $i = \lambda D / a$ permet de déterminer une inconnue si les trois autres sont connues. En pratique, on ne mesure pas une seule interfrange (trop imprécis), mais la longueur $\ell$ occupée par $N$ interfranges (distance entre la frange d'ordre $k$ et la frange d'ordre $k+N$).

\begin{methode}[Mesurer une longueur d'onde par les fentes de Young]
\begin{enumerate}
\item \textbf{Matériel} : Laser (source monochromatique, $\lambda$ inconnu), diaphragme à deux fentes (écart $a$ connu, ex. $a = \SI{0,50}{mm}$), écran (papier millimétré), règle, monture optique.
\item \textbf{Réglage} : Placer l'écran à une distance $D$ grande (ex. $D = \SI{1,50}{m}$) pour avoir des franges bien écartées. Vérifier $D \gg a$.
\item \textbf{Observation} : Repérer les franges brillantes. Compter $N$ interfranges (ex. $N = 10$ ou $20$) et mesurer la distance $\ell$ entre les centres des franges extrêmes (à la règle ou au pied à coulisse).
\item \textbf{Calcul} : Déduire l'interfrange moyenne $i = \ell / N$.
\item \textbf{Résultat} : Calculer $\lambda = \dfrac{i a}{D} = \dfrac{\ell a}{N D}$.
\item \textbf{Incertitude} : Estimer les incertitudes sur $\ell$, $D$, $a$ et propager pour donner $\lambda \pm \Delta\lambda$.
\end{enumerate}
\end{methode}

\begin{experience}[TP : Mesure de l'interfrange et détermination de $\lambda$]
\textbf{Objectif} : Mesurer la longueur d'onde d'un laser à l'aide des fentes de Young.
\textbf{Matériel} : Laser He-Ne ($\lambda \approx \SI{633}{nm}$), support à fentes ($a = \SI{0,40}{mm}$), écran gradué, règle métallique, monture optique, chambre noire ou local obscurci.
\textbf{Protocole} :
\begin{enumerate}
\item Aligner le laser, les fentes et l'écran. Distance fentes-écran $D = \SI{2,00}{m}$.
\item Observer les franges sur l'écran. Repérer la frange centrale ($k=0$).
\item Compter $N = 20$ interfranges de part et d'autre du centre (franges $k=-10$ à $k=+10$). Mesurer la distance totale $\ell$ entre ces deux franges extrêmes.
\item Répéter la mesure 3 fois pour estimer l'incertitude type $\Delta\ell$.
\end{enumerate}
\textbf{Analyse} : $i = \ell / (2N)$ (car $2N$ interfranges entre $k=-N$ et $k=+N$). $\lambda = i a / D$.
\textbf{Interprétation} : Comparer la valeur trouvée à la valeur de référence. Discuter les sources d'erreur (alignement, épaisseur traits, lecture règle).
\end{experience}

\begin{exemplebox}[Application numérique : lumière rouge He-Ne]
On utilise un laser He-Ne ($\lambda \approx \SI{632,8}{nm}$) avec des fentes d'écart $a = \SI{0,40}{mm}$. L'écran est à $D = \SI{2,00}{m}$.
\begin{enumerate}
\item Calculer l'interfrange $i$.
\item Quelle distance $\ell$ occupent $N = 20$ interfranges ?
\item Si on mesure $\ell = \SI{6,35}{cm}$ pour $N=20$, quelle valeur de $\lambda$ en déduit-on ?
\end{enumerate}
\textbf{Solution.}
\begin{enumerate}
\item $i = \dfrac{\lambda D}{a} = \dfrac{632,8 \times 10^{-9} \times 2,00}{0,40 \times 10^{-3}} = \dfrac{1265,6 \times 10^{-9}}{0,40 \times 10^{-3}} = \SI{3,164 \times 10^{-3}}{m} = \SI{3,16}{mm}$.
\item $\ell = N \times i = 20 \times \SI{3,164}{mm} = \SI{63,28}{mm} \approx \SI{6,33}{cm}$. Facilement mesurable à la règle.
\item $i_{\text{mes}} = \ell/N = \SI{63,5}{mm}/20 = \SI{3,175}{mm}$.
$\lambda_{\text{mes}} = \dfrac{i_{\text{mes}} a}{D} = \dfrac{3,175 \times 10^{-3} \times 0,40 \times 10^{-3}}{2,00} = \SI{635 \times 10^{-9}}{m} = \SI{635}{nm}$.
L'erreur relative est $\approx 0,3\%$, excellente pour une mesure à la règle.
\end{enumerate}
\end{exemplebox}

\courssub{Visualisation de la répartition de l'intensité}

Le profil d'intensité $I(x)$ sur l'écran est le produit de l'enveloppe de diffraction (par une fente) et du peigne d'interférences (deux fentes). Pour des fentes très fines (diffraction large), on observe des franges d'intensité quasi constante.

\begin{popfigure}
\begin{tikzpicture}
\begin{axis}[
    width=\linewidth,
    height=7cm,
    axis lines=middle,
    xlabel={$x$ (mm)},
    ylabel={$I(x)$ (u.a.)},
    xlabel style={anchor=north west},
    ylabel style={anchor=south east},
    domain=-10:10,
    samples=800,
    xmin=-10.5, xmax=10.5,
    ymin=-0.1, ymax=1.1,
    xtick={-10,-8,-6,-4,-2,0,2,4,6,8,10},
    ytick={0,0.5,1},
    tick label style={font=\small},
    clip=false,
    grid=major,
    grid style={popline, opacity=0.3},
]
% Paramètres
% Courbe intensité : cos^2 modulé par gaussienne
\addplot[popcurve, thick, domain=-10:10] {exp(-x*x/(2*20*20)) * (cos(180*x/2))^2};
% Marqueurs d'ordre k
\pgfplotsinvokeforeach{-8,-6,-4,-2,0,2,4,6,8}{
    \draw[popmuted, dashed, thin] (#1, 0) -- (#1, 1.05);
}
% Flèches interfrange
\draw[poporange, <->, thick] (0, -0.15) -- (2, -0.15) node[midway, below=2pt, font=\small] {$i$};
\draw[poporange, <->, thick] (2, -0.15) -- (4, -0.15) node[midway, below=2pt, font=\small] {$i$};
% Frange centrale
\node[popdark, font=\small] at (0, 1.05) {$k=0$};
\end{axis}
\end{tikzpicture}
\legende{Profil d'intensité $I(x)$ pour les fentes de Young (fentes fines). Les maxima (franges brillantes) sont équidistants de l'interfrange $i = \lambda D/a$. L'enveloppe gaussienne (traits pleins) modélise la diffraction par chaque fente individuelle.}
\end{popfigure}

\courssub{Exercice résolu : Détermination de l'écart des fentes}

\begin{exoresolu}[Mesure de l'écart $a$ par interférences]
On éclaire deux fentes $S_1$ et $S_2$ par un laser à diode de longueur d'onde $\lambda = \SI{650,0}{nm}$. L'écran est placé à $D = \SI{1,20}{m}$. On mesure la distance entre les franges brillantes d'ordres $k=5$ et $k=15$ : $\ell = \SI{18,7}{mm}$.
Déterminer l'écart $a$ entre les fentes et son incertitude si $\Delta\ell = \pm \SI{0,2}{mm}$, $\Delta D = \pm \SI{2}{mm}$, $\Delta\lambda$ négligeable.

\tcblower
\textbf{Analyse.}
La distance $\ell$ correspond à $N = 15 - 5 = 10$ interfranges.
L'interfrange moyenne est $i = \ell / N$.
La relation $i = \lambda D / a$ donne $a = \lambda D / i = \lambda D N / \ell$.

\textbf{Calcul de $a$.}
\[
i = \frac{\SI{18,7}{mm}}{10} = \SI{1,87}{mm} = 1,87 \times 10^{-3}\,\text{m}.
\]
\[
a = \frac{\lambda D}{i} = \frac{650,0 \times 10^{-9} \times 1,20}{1,87 \times 10^{-3}} = \frac{780 \times 10^{-9}}{1,87 \times 10^{-3}} = \SI{417 \times 10^{-6}}{m} = \SI{0,417}{mm}.
\]

\textbf{Incertitude relative sur $a$.}
$a \propto D / \ell$ (car $N$ exact, $\lambda$ exact).
\[
\frac{\Delta a}{a} = \sqrt{\left(\frac{\Delta D}{D}\right)^2 + \left(\frac{\Delta \ell}{\ell}\right)^2}
= \sqrt{\left(\frac{\num{2}}{\num{1200}}\right)^2 + \left(\frac{\num{0,2}}{\num{18,7}}\right)^2}
= \sqrt{(\num{1,67e-3})^2 + (\num{10,7e-3})^2}
\approx \num{10,8e-3}.
\]
$\Delta a = a \times \num{10,8e-3} \approx \SI{0,417}{mm} \times 0,0108 \approx \SI{0,0045}{mm}$.

\textbf{Résultat :} $a = \SI{0,417 \pm 0,005}{mm}$ (arrondi au micron).
L'incertitude est dominée par la mesure de $\ell$ (règle). Pour l'améliorer, augmenter $N$ (mesurer sur plus de franges) ou utiliser un capteur CCD.
\end{exoresolu}

\begin{savaistu}[Métrologie optique et applications camerounaises]
La mesure d'interfranges est au cœur de l'\textbf{interférométrie}, technique de mesure de longueur ultra-précise (nanomètre).
\begin{itemize}
\item \textbf{Contrôle qualité industriel} : Vérification du polissage de lentilles pour les caméras de surveillance (sécurité à Yaoundé/Douala) ou les objectifs de smartphones.
\item \textbf{Télédétection spatiale} : Les satellites (comme ceux utilisés pour la cartographie du Cameroun ou la surveillance agricole) utilisent l'interférométrie radar (InSAR) pour mesurer le relief du sol avec une précision centimétrique. Le principe repose sur la mesure de la différence de phase (équivalente à une fraction d'interfrange) entre deux signaux radar.
\item \textbf{Historique} : Albert Michelson (prix Nobel 1907) a mesuré le mètre étalon en longueurs d'onde de lumière (cadmium rouge, $\lambda = \SI{643,8}{nm}$) dès 1892, ouvrant la voie à la définition du mètre de 1960 (Krypton-86) à 1983 (vitesse de la lumière).
\end{itemize}
Aujourd'hui, au laboratoire de l'Université de Yaoundé I ou de Douala, des étudiants en Master utilisent des interféromètres de Michelson ou de Fabry-Pérot pour caractériser des couches minces (cellules solaires, capteurs gazeux), appliquant directement la relation $i = \lambda D/a$ ou sa généralisation en phase.
\end{savaistu}

\begin{exercice}[Entraînement : Franges colorées en lumière blanche]
On éclaire les fentes de Young ($a = \SI{0,30}{mm}$, $D = \SI{1,00}{m}$) par de la lumière blanche (spectre continu $[\SI{400}{nm}; \SI{800}{nm}]$).
\begin{enumerate}
\item Calculer l'interfrange $i_{\text{bleu}}$ pour $\lambda = \SI{400}{nm}$ et $i_{\text{rouge}}$ pour $\lambda = \SI{700}{nm}$.
\item Sur l'écran, la frange centrale ($k=0$) est blanche. Quelle est la couleur de la première frange brillante ($k=1$) du côté positif ? Justifier.
\item À partir de quel ordre $k$ les franges d'ordre $k$ pour le bleu et l'ordre $k+1$ pour le rouge se chevauchent-elles (coïncidence) ? Que voit-on alors ?
\end{enumerate}
\end{exercice}

\begin{corrige}[Exercice : Franges colorées en lumière blanche]
\begin{enumerate}
\item $i = \lambda D/a$.
$i_{\text{bleu}} = \frac{400 \times 10^{-9} \times 1,00}{0,30 \times 10^{-3}} = \SI{1,33}{mm}$.
$i_{\text{rouge}} = \frac{700 \times 10^{-9} \times 1,00}{0,30 \times 10^{-3}} = \SI{2,33}{mm}$.
\item En $k=1$, $x_1 = i$. Le bleu est à $\SI{1,33}{mm}$, le rouge à $\SI{2,33}{mm}$. La frange $k=1$ est donc \textbf{bleue/violette} sur son bord interne (proche du centre) et \textbf{rouge} sur son bord externe. Elle apparaît \textbf{colorée (irisée)}.
\item Coïncidence : on cherche $k_b \lambda_b = k_r \lambda_r$ avec $k_r < k_b$ (car $i_r > i_b$).
$k_b \times 400 = k_r \times 700 \Rightarrow 4 k_b = 7 k_r$.
Plus petite solution entière : $k_b = 7$, $k_r = 4$.
Position $x = 7 \times i_{\text{bleu}} = 7 \times \SI{1,33}{mm} = \SI{9,33}{mm}$ (ou $4 \times \SI{2,33}{mm} = \SI{9,33}{mm}$).
\textbf{Coïncidence exacte pour $k_{\text{bleu}}=7$ et $k_{\text{rouge}}=4$}.
Au-delà, les franges d'ordres différents se superposent pour toutes les couleurs : le mélange reconstruit la lumière blanche, les franges deviennent blanches et peu contrastées.
\end{enumerate}
\end{corrige}

\coursec{Interférences en lumière polychromatique et en lumière blanche}

Dans la section précédente, nous avons établi l'expression de l'interfrange $i = \dfrac{\lambda D}{a}$ pour une source \emph{monochromatique} : une seule longueur d'onde $\lambda$, un seul système de franges régulières. Mais la lumière qui nous entoure — celle du Soleil sur les marchés de Yaoundé, celle d'une lampe à incandescence — n'est presque jamais monochromatique. Que devient la figure d'interférences lorsque la source émet \emph{plusieurs} radiations à la fois ? C'est la question que nous allons résoudre ici, et sa réponse explique des spectacles familiers : les couleurs mouvantes d'une bulle de savon ou d'une flaque d'huile sur la route après la pluie.

\courssub{Lumière polychromatique : superposition des systèmes de franges}

\begin{definition}[Lumière polychromatique]
Une lumière \cle{polychromatique} est constituée de la superposition de plusieurs radiations monochromatiques de longueurs d'onde différentes. Si elle en contient un nombre fini (deux, trois...), on parle de lumière dichromatique, trichromatique... La \cle{lumière blanche} contient, elle, toutes les radiations visibles, de $\lambda \approx \SI{400}{nm}$ (violet) à $\lambda \approx \SI{800}{nm}$ (rouge).
\end{definition}

L'idée directrice est simple : chaque radiation de longueur d'onde $\lambda$ se comporte \emph{indépendamment} des autres (les ondes lumineuses de fréquences différentes ne peuvent pas interférer entre elles de façon stable). Chacune produit donc \emph{son propre} système de franges, avec \emph{son propre} interfrange :
\[ i = \frac{\lambda D}{a} \]
où $a$ est l'écartement des deux fentes de Young et $D$ la distance fentes--écran. Comme $D$ et $a$ sont fixés par le dispositif, l'interfrange est \textbf{proportionnelle à} $\lambda$ : les franges rouges ($\lambda$ grand) sont plus espacées que les franges bleues ($\lambda$ petit). Sur l'écran, on observe la superposition de tous ces systèmes.

\begin{propriete}[Frange centrale commune]
En lumière polychromatique, la \cle{frange centrale} est commune à toutes les radiations : en $x = 0$, la différence de marche est $\delta = 0$ \emph{quelle que soit} $\lambda$, donc toutes les radiations y sont en phase et donnent une frange brillante. C'est le seul point où tous les systèmes de franges coïncident rigoureusement.
\end{propriete}

En effet, $\delta = \dfrac{ax}{D}$ ne fait pas intervenir $\lambda$ : en $x = 0$, $\delta = 0 = 0 \times \lambda$ pour toute radiation. Partout ailleurs, les systèmes se décalent progressivement les uns par rapport aux autres, d'autant plus vite que les longueurs d'onde diffèrent.

\begin{popfigure}
\begin{center}
\begin{tikzpicture}[scale=1]
% Axe des positions
\draw[->,thick,popdark] (-5.2,0) -- (5.2,0) node[below left]{$x$};
\draw[thick,popdark] (0,-0.15) -- (0,0.15) node[above right=2pt and -2pt,popdark]{\footnotesize $0$};
% Franges bleues : interfrange 0.75
\foreach \k in {-6,-5,-4,-3,-2,-1,0,1,2,3,4,5,6}{
  \draw[line width=1.6pt,popblue] ({0.75*\k},0.35) -- ({0.75*\k},1.35);}
\node[popblue] at (0,1.6) {\footnotesize Système bleu : $i_2=\lambda_2 D/a$};
% Franges rouges : interfrange 1.0
\foreach \k in {-4,-3,-2,-1,0,1,2,3,4}{
  \draw[line width=1.6pt,poporange] ({1.0*\k},-1.35) -- ({1.0*\k},-0.35);}
\node[poporange] at (0,-1.7) {\footnotesize Système rouge : $i_1=\lambda_1 D/a$};
% Frange centrale commune
\draw[line width=2.2pt,popdark] (0,-1.35) -- (0,1.35);
\node[popdark,align=center] at (0,2.35) {\footnotesize frange centrale\\ \footnotesize commune ($\delta=0$)};
% Coïncidence k1=3 rouge = k2=4 bleu à x=3
\draw[dashed,popgreen,thick] (3,-1.5) -- (3,1.5);
\node[popgreen,align=center] at (3,2.35) {\footnotesize coïncidence\\ \footnotesize $3\lambda_1=4\lambda_2$};
\draw[dashed,popgreen,thick] (-3,-1.5) -- (-3,1.5);
\node[popgreen,align=center] at (-3,2.35) {\footnotesize coïncidence\\ \footnotesize symétrique};
\end{tikzpicture}
\end{center}
\legende{Superposition de deux systèmes de franges (ici $\lambda_1 = \SI{600}{nm}$ rouge et $\lambda_2 = \SI{450}{nm}$ bleu, donc $i_1/i_2 = 4/3$). Toutes les franges partent de la frange centrale commune ; les franges brillantes rouges d'ordre $3$ coïncident avec les bleues d'ordre $4$.}
\end{popfigure}

\courssub{Coïncidence des franges de deux radiations}

\begin{propriete}[Condition de coïncidence]
Deux radiations de longueurs d'onde $\lambda_1$ et $\lambda_2$ présentent une \cle{coïncidence} de franges brillantes en un point $M$ si leurs franges brillantes s'y superposent, c'est-à-dire si :
\[ \delta(M) = k_1 \lambda_1 = k_2 \lambda_2 \qquad \text{avec } k_1, k_2 \in \mathbb{Z} \]
Le point $M$ est alors situé à la distance $x = \dfrac{\delta \, D}{a}$ du centre de l'écran.
\end{propriete}

\begin{methode}[Exploiter une coïncidence de franges]
Pour trouver les points où deux radiations $\lambda_1$ et $\lambda_2$ coïncident :
\begin{enumerate}
\item Écrire la condition de coïncidence : $k_1 \lambda_1 = k_2 \lambda_2$.
\item Simplifier le rapport $\dfrac{\lambda_1}{\lambda_2} = \dfrac{k_2}{k_1}$ en une fraction d'entiers \emph{irréductible}.
\item En déduire les plus petits ordres entiers $k_1$ et $k_2$ solutions ; les autres coïncidences s'obtiennent en multipliant par $2, 3, \dots$
\item Calculer la différence de marche $\delta = k_1\lambda_1$ puis, si demandé, la position $x = \dfrac{\delta D}{a} = k_1 i_1 = k_2 i_2$.
\end{enumerate}
\end{methode}

\begin{exemplebox}[Coïncidence rouge--bleu]
Soient $\lambda_1 = \SI{600}{nm}$ (rouge) et $\lambda_2 = \SI{450}{nm}$ (bleu). On cherche les entiers tels que $k_1 \lambda_1 = k_2 \lambda_2$ :
\[ \frac{k_2}{k_1} = \frac{\lambda_1}{\lambda_2} = \frac{600}{450} = \frac{4}{3} \]
La plus petite solution est $k_1 = 3$ et $k_2 = 4$. Vérification :
\[ 3\lambda_1 = 3 \times 600 = \SI{1800}{nm} \qquad ; \qquad 4\lambda_2 = 4 \times 450 = \SI{1800}{nm} \]
Les franges brillantes d'ordre $3$ du rouge et d'ordre $4$ du bleu coïncident à la différence de marche $\delta = \SI{1,8}{\micro m}$. Les coïncidences suivantes ont lieu en $\delta = \SI{3,6}{\micro m}$, $\SI{5,4}{\micro m}$, etc.
\end{exemplebox}

\begin{exoresolu}[Distance entre deux coïncidences consécutives]
\textbf{Énoncé.} On éclaire des fentes de Young avec une lumière dichromatique : $\lambda_1 = \SI{600}{nm}$ et $\lambda_2 = \SI{450}{nm}$. On donne $a = \SI{0,50}{mm}$ et $D = \SI{1,5}{m}$. (a) Calculer les interfranges $i_1$ et $i_2$. (b) Déterminer la position de la première coïncidence de franges brillantes hors du centre. (c) En déduire la distance entre deux coïncidences consécutives.
\tcblower
\textbf{Solution.}
(a) Interfranges :
\[ i_1 = \frac{\lambda_1 D}{a} = \frac{600\times 10^{-9}\times 1{,}5}{0{,}50\times 10^{-3}} = \SI{1,8}{mm} \qquad ; \qquad i_2 = \frac{\lambda_2 D}{a} = \SI{1,35}{mm} \]
(b) La première coïncidence hors centre correspond à $k_1 = 3$, $k_2 = 4$ (voir l'exemple ci-dessus), soit $\delta = \SI{1,8}{\micro m}$. Sa position :
\[ x = \frac{\delta D}{a} = k_1 i_1 = 3 \times 1{,}8 = \SI{5,4}{mm} \]
Vérification avec le bleu : $x = k_2 i_2 = 4 \times 1{,}35 = \SI{5,4}{mm}$. Les deux calculs concordent.
(c) Les coïncidences se répètent pour $\delta = n \times \SI{1,8}{\micro m}$ ($n \in \mathbb{N}^*$), donc la distance entre deux coïncidences consécutives est constante et égale à $\Delta x = \SI{5,4}{mm}$ (c'est le pas spatial de la figure de coïncidence, égal au plus petit commun multiple des interfranges $i_1$ et $i_2$).
\end{exoresolu}

\courssub{Interférences en lumière blanche}

Éclairons maintenant les fentes de Young en \cle{lumière blanche} : c'est la superposition d'une \emph{infinité} de radiations, chacune avec son interfrange propre. Que voit-on sur l'écran ?

\begin{aretenir}[Aspect de la figure en lumière blanche]
\begin{itemize}
\item La \cle{frange centrale} est \textbf{blanche} : en $x = 0$, $\delta = 0$ pour toutes les radiations, qui y arrivent toutes en phase et se recomposent en lumière blanche.
\item De part et d'autre, on observe quelques franges \cle{irisées} : elles présentent des teintes colorées (rougeoyantes côté extérieur, bleutées côté centre).
\item Plus loin du centre, les systèmes de franges des différentes radiations se décalent tellement qu'ils finissent par se \cle{brouiller} : l'écran devient uniformément éclairé, d'une teinte blanchâtre appelée \emph{blanc d'ordre supérieur}.
\end{itemize}
\end{aretenir}

Pourquoi des couleurs ? En un point $M$ proche du centre, la différence de marche $\delta(M)$ est fixée. Les radiations pour lesquelles $\delta = k\lambda$ y sont \emph{renforcées} (frange brillante), tandis que celles pour lesquelles $\delta = \left(k + \tfrac{1}{2}\right)\lambda$ y sont \emph{éteintes} (frange sombre). La lumière reçue en $M$ est donc la lumière blanche \emph{privée} de certaines radiations et enrichie en d'autres : d'où une teinte particulière, appelée \cle{teinte sensible}. Comme $\delta$ varie continûment avec $x$, les teintes défilent : ce sont les irisations.

\begin{popfigure}
\begin{center}
\begin{tikzpicture}[scale=1]
% Frise : frange centrale blanche, franges irisées, brouillage
\fill[white] (-0.35,0) rectangle (0.35,1.6);
\draw[popdark,thick] (-0.35,0) rectangle (0.35,1.6);
% Franges irisées gauche et droite (dégradé simulé par bandes colorées)
\foreach \s in {1,-1}{
  \fill[popblue!70] ({\s*0.35},0) rectangle ({\s*0.85},1.6);
  \fill[popgreen!60] ({\s*0.85},0) rectangle ({\s*1.35},1.6);
  \fill[popgold!80] ({\s*1.35},0) rectangle ({\s*1.95},1.6);
  \fill[poporange!80] ({\s*1.95},0) rectangle ({\s*2.6},1.6);
  \fill[poporange!40] ({\s*2.6},0) rectangle ({\s*3.3},1.6);
  \fill[popmuted!25] ({\s*3.3},0) rectangle ({\s*5.0},1.6);
}
\draw[popdark] (-5.0,0) rectangle (5.0,1.6);
% Annotations
\node[align=center] at (0,-0.5) {\footnotesize frange centrale blanche};
\draw[->,popdark] (0,-0.35) -- (0,-0.05);
\node[align=center] at (2.4,-0.5) {\footnotesize franges irisées};
\draw[->,popdark] (2.4,-0.35) -- (2.4,-0.05);
\node[align=center] at (4.3,-0.5) {\footnotesize brouillage};
\draw[->,popdark] (4.3,-0.35) -- (4.3,-0.05);
\node[align=center] at (-2.4,-0.5) {\footnotesize franges irisées};
\draw[->,popdark] (-2.4,-0.35) -- (-2.4,-0.05);
\node[align=center] at (-4.3,-0.5) {\footnotesize brouillage};
\draw[->,popdark] (-4.3,-0.35) -- (-4.3,-0.05);
\end{tikzpicture}
\end{center}
\legende{Figure d'interférences en lumière blanche (représentation schématique) : frange centrale blanche, quelques franges irisées de part et d'autre, puis brouillage progressif vers le blanc d'ordre supérieur.}
\end{popfigure}

\begin{exemplebox}[Radiation éteinte dans une frange irisée]
En un point $M$ où $\delta = \SI{1,20}{\micro m}$, quelles radiations visibles sont éteintes ? Il faut $\delta = \left(k + \tfrac{1}{2}\right)\lambda$, soit $\lambda = \dfrac{2\delta}{2k+1}$ :
\begin{align*}
k=0 :& \quad \lambda = \SI{2400}{nm} \quad \text{(infrarouge, invisible)}\\
k=1 :& \quad \lambda = \frac{2{,}4}{3} = \SI{800}{nm} \quad \text{(limite du rouge)}\\
k=2 :& \quad \lambda = \frac{2{,}4}{5} = \SI{480}{nm} \quad \text{(bleu, éteint !)}\\
k=3 :& \quad \lambda = \frac{2{,}4}{7} \approx \SI{343}{nm} \quad \text{(ultraviolet)}
\end{align*}
Le bleu ($\lambda = \SI{480}{nm}$) est éteint en $M$ : la teinte y est donc jaune-orangée (couleur complémentaire du bleu). C'est exactement le mécanisme des teintes des bulles de savon.
\end{exemplebox}

\begin{attention}[Erreur fréquente : les franges en lumière blanche]
En lumière blanche, \textbf{seule la frange centrale est blanche}. Affirmer que « toutes les franges sont blanches » est faux : l'interfrange $i = \lambda D/a$ dépend de $\lambda$, donc chaque radiation place ses franges à des positions différentes. Dès qu'on quitte le centre, les systèmes se décalent : les franges sont irisées, puis brouillées. Autre piège : la frange centrale est blanche \emph{parce que} $\delta = 0$ ne dépend pas de $\lambda$ — c'est cette justification qui est attendue au Baccalauréat, pas seulement le constat.
\end{attention}

\begin{savaistu}[Les ailes du papillon Morpho]
Les irisations bleues métalliques des ailes du papillon \emph{Morpho}, que l'on peut admirer dans les forêts du Sud et du Centre du Cameroun, ne proviennent \emph{pas} de pigments ! Elles résultent d'interférences lumineuses en lumière blanche sur des microstructures en couches fines présentes sur les écailles des ailes : certaines longueurs d'onde sont renforcées, d'autres éteintes, selon l'angle d'observation. C'est le même phénomène que les couleurs des bulles de savon, des flaques d'huile sur l'eau ou de la face lisible d'un CD. Les ingénieurs s'en inspirent aujourd'hui pour créer des « couleurs structurelles » sans colorants, inaltérables dans le temps.
\end{savaistu}

\begin{exercice}[Coïncidence en lumière dichromatique]
Une source émet deux radiations $\lambda_1 = \SI{650}{nm}$ et $\lambda_2 = \SI{520}{nm}$. Le dispositif de Young a pour paramètres $a = \SI{0,40}{mm}$ et $D = \SI{2,0}{m}$.
\begin{enumerate}
\item Calculer les interfranges $i_1$ et $i_2$.
\item Montrer que la première coïncidence de franges brillantes hors du centre correspond à $k_1 = 4$ et $k_2 = 5$.
\item À quelle distance du centre se trouve-t-elle ?
\end{enumerate}
\end{exercice}

\begin{corrige}[Exercice : coïncidence en lumière dichromatique]
1. $i_1 = \dfrac{\lambda_1 D}{a} = \dfrac{650\times 10^{-9}\times 2{,}0}{0{,}40\times 10^{-3}} = \SI{3,25}{mm}$ et $i_2 = \dfrac{\lambda_2 D}{a} = \SI{2,60}{mm}$.

2. Condition de coïncidence : $k_1\lambda_1 = k_2\lambda_2$, soit $\dfrac{k_2}{k_1} = \dfrac{650}{520} = \dfrac{5}{4}$ (fraction irréductible). La plus petite solution est donc $k_1 = 4$, $k_2 = 5$. Vérification : $4 \times 650 = 5 \times 520 = \SI{2600}{nm}$.

3. $x = k_1 i_1 = 4 \times 3{,}25 = \SI{13}{mm}$ (vérification : $x = k_2 i_2 = 5 \times 2{,}60 = \SI{13}{mm}$).
\end{corrige}

\begin{aretenir}[L'essentiel de la section]
\begin{itemize}
\item En lumière polychromatique, chaque radiation produit son propre système de franges d'interfrange $i = \lambda D/a$ ; la frange centrale ($\delta = 0$) est commune à toutes.
\item Deux radiations coïncident en $M$ si $\delta(M) = k_1\lambda_1 = k_2\lambda_2$ ; on résout en simplifiant le rapport $\lambda_1/\lambda_2$ en fraction d'entiers.
\item En lumière blanche : frange centrale blanche, franges irisées (radiations éteintes quand $\delta = (k+\tfrac12)\lambda$), puis brouillage en blanc d'ordre supérieur.
\item Bulles de savon, flaques d'huile, ailes de Morpho : manifestations quotidiennes des interférences en lumière blanche.
\end{itemize}
\end{aretenir}

\coursec{Effet Doppler : étude qualitative et applications}

Dans la section précédente, nous avons vu que la lumière blanche, composée de radiations de longueurs d'onde différentes, produit des systèmes de franges qui se superposent : la position d'une frange dépend directement de la longueur d'onde $\lambda$. Autrement dit, la couleur perçue est une signature de $\lambda$. Or que se passerait-il si la longueur d'onde \emph{reçue} par un observateur n'était pas exactement celle \emph{émise} par la source ? C'est précisément ce qui arrive lorsque la source et le récepteur sont en mouvement l'un par rapport à l'autre : c'est l'\cle{effet Doppler}, phénomène universel qui concerne aussi bien le son d'une mototaxi à Douala que la lumière des galaxies lointaines.

\courssub{Une observation quotidienne : la sirène d'ambulance}

Tu l'as certainement remarqué dans les rues de Yaoundé : lorsqu'une ambulance s'approche de toi, le son de sa sirène paraît \emph{aigu} ; dès qu'elle te dépasse et s'éloigne, le même son devient soudainement plus \emph{grave}. Pourtant, le conducteur n'a rien changé : la sirène émet en permanence un son de fréquence $f_E$ constante.

\begin{experience}[La sirène qui change de hauteur sans changer de fréquence]
\textbf{Observation.} Un observateur immobile au bord de la route entend :
\begin{itemize}
\item un son \textbf{plus aigu} (fréquence $f_R > f_E$) quand l'ambulance \emph{s'approche} ;
\item un son \textbf{plus grave} (fréquence $f_R < f_E$) quand elle \emph{s'éloigne}.
\end{itemize}
\textbf{Interprétation.} La hauteur d'un son est liée à sa fréquence : plus la fréquence est élevée, plus le son est aigu. La fréquence \emph{reçue} par l'oreille n'est donc pas égale à la fréquence \emph{émise} par la sirène. Le seul ingrédient nouveau est le \cle{mouvement relatif} entre la source et l'observateur.
\end{experience}

\courssub{Explication qualitative : le resserrement des fronts d'onde}

Modélisons la situation. Une sirène immobile émet des \cle{fronts d'onde} sphériques (des cercles en vue de dessus) espacés régulièrement de la longueur d'onde $\lambda = v_{\text{son}}\,T$, où $v_{\text{son}} \approx \SI{340}{m.s^{-1}}$ est la célérité du son dans l'air et $T$ la période d'émission.

Lorsque la source \emph{se déplace} vers la droite à la vitesse $v_S$, chaque front d'onde est émis depuis une position différente : les cercles ne sont plus concentriques.

\begin{itemize}
\item \textbf{Devant la source} (du côté où elle avance) : chaque nouveau front est émis un peu plus près du précédent. Les fronts se \textbf{resserrent} : la longueur d'onde apparente $\lambda'$ diminue. Comme $f = \dfrac{v_{\text{son}}}{\lambda'}$, la fréquence reçue \textbf{augmente} : le son est plus aigu.
\item \textbf{Derrière la source} : les fronts s'\textbf{écartent} : $\lambda''$ augmente, la fréquence reçue \textbf{diminue} : le son est plus grave.
\end{itemize}

\begin{popfigure}
\begin{center}
\begin{tikzpicture}[scale=0.85]
% Fronts d'onde décalés (source se déplaçant vers la droite)
\foreach \i/\r in {0/0.35, 1/0.75, 2/1.15, 3/1.55, 4/1.95, 5/2.35}{
  \draw[popblue, thick] ({0.28*\i},0) circle (\r);
}
% Source
\fill[poporange] (1.68,0) circle (0.09) node[above=2pt, popdark] {\footnotesize source $S$};
% Vecteur vitesse
\draw[-{Stealth}, very thick, poporange] (1.68,-0.35) -- (3.0,-0.35) node[midway, below, popdark] {\footnotesize $\vec{v}_S$};
% Observateurs
\node[popgreen, font=\footnotesize\bfseries] at (4.6,1.6) {Observateur avant};
\node[popdark, font=\footnotesize] at (4.6,1.15) {son plus \textbf{aigu}};
\node[popgreen, font=\footnotesize\bfseries] at (-3.3,1.6) {Observateur arrière};
\node[popdark, font=\footnotesize] at (-3.3,1.15) {son plus \textbf{grave}};
\fill[popgreen] (4.1,0) circle (0.09);
\fill[popgreen] (-2.9,0) circle (0.09);
% Annotations resserrement / écartement
\draw[{Stealth}-{Stealth}, poppink, thick] (2.15,-1.15) -- (2.45,-1.15) node[midway, below, poppink] {\footnotesize $\lambda'$ court};
\draw[{Stealth}-{Stealth}, poppurple, thick] (-1.6,-1.9) -- (-0.9,-1.9) node[midway, below, poppurple] {\footnotesize $\lambda''$ long};
\end{tikzpicture}
\end{center}
\legende{Fronts d'onde émis par une source en mouvement vers la droite : ils se resserrent devant (fréquence reçue plus élevée) et s'écartent derrière (fréquence reçue plus faible).}
\end{popfigure}

\begin{definition}[Effet Doppler]
L'\cle{effet Doppler} est la modification de la \textbf{fréquence} (et donc de la \textbf{longueur d'onde}) perçue par un récepteur lorsque la \textbf{source} et le \textbf{récepteur} sont en \textbf{mouvement relatif} l'un par rapport à l'autre. Ce phénomène concerne \emph{toutes} les ondes : ondes sonores, ondes radio, lumière.
\end{definition}

\begin{propriete}[Sens du décalage de fréquence (admise)]
Soit $f_E$ la fréquence émise par la source et $f_R$ la fréquence reçue par le récepteur :
\begin{itemize}
\item si la source et le récepteur \textbf{se rapprochent}, alors $f_R > f_E$ (son plus aigu, longueur d'onde reçue plus courte) ;
\item si la source et le récepteur \textbf{s'éloignent}, alors $f_R < f_E$ (son plus grave, longueur d'onde reçue plus longue).
\end{itemize}
De plus, le décalage relatif $\dfrac{|f_R - f_E|}{f_E}$ est \textbf{d'autant plus grand que la vitesse relative est élevée} : la mesure du décalage permet donc de remonter à la vitesse.
\end{propriete}

\begin{propriete}[Expression quantitative pour le son (complément, non exigible)]
Pour une source se déplaçant à la vitesse $v_S$ vers un observateur immobile, dans un milieu où le son se propage à la célérité $v_{\text{son}}$, la fréquence reçue s'écrit :
\[ f_R = f_E \times \dfrac{v_{\text{son}}}{v_{\text{son}} - v_S} \qquad (\text{rapprochement}), \]
et, lorsque la source s'éloigne :
\[ f_R = f_E \times \dfrac{v_{\text{son}}}{v_{\text{son}} + v_S} \qquad (\text{éloignement}). \]
Ces formules confirment le raisonnement qualitatif : au dénominateur, $v_{\text{son}} - v_S < v_{\text{son}}$ donne $f_R > f_E$ à l'approche. Elles ne sont pas exigibles au Bac, mais elles permettent de vérifier la cohérence des ordres de grandeur.
\end{propriete}

\begin{attention}[La fréquence émise ne change pas !]
Erreur très fréquente au Bac : croire que la sirène « émet plus aigu » quand elle approche. \textbf{Non !} La fréquence $f_E$ \emph{émise} par la source reste rigoureusement constante. C'est la fréquence $f_R$ \emph{reçue} qui diffère, à cause du mouvement relatif qui modifie l'espacement des fronts d'onde dans le milieu. De même, en lumière, l'étoile n'émet pas « plus rouge » : c'est la lumière \emph{reçue} qui est décalée.
\end{attention}

\courssub{L'effet Doppler lumineux : décalage vers le rouge ou vers le bleu}

La lumière est une onde : elle subit donc aussi l'effet Doppler. Dans le spectre visible, le \textbf{rouge} correspond aux grandes longueurs d'onde ($\lambda \approx \SI{650}{nm}$) et le \textbf{bleu} aux petites ($\lambda \approx \SI{450}{nm}$).

\begin{propriete}[Effet Doppler lumineux]
\begin{itemize}
\item Une source lumineuse qui \textbf{s'éloigne} de l'observateur : la longueur d'onde reçue \emph{augmente} ; la lumière est décalée vers le \textbf{rouge} : on parle de \cle{redshift}.
\item Une source lumineuse qui \textbf{se rapproche} : la longueur d'onde reçue \emph{diminue} ; la lumière est décalée vers le \textbf{bleu} : on parle de \cle{blueshift}.
\end{itemize}
\end{propriete}

Comment le mesure-t-on concrètement ? La lumière d'une étoile ou d'une galaxie présente des \cle{raies d'absorption} caractéristiques (raies de l'hydrogène, du calcium...), dont les longueurs d'onde sont parfaitement connues grâce aux mesures en laboratoire sur Terre. En comparant le spectre de la galaxie au spectre de référence, on mesure le décalage $\Delta\lambda$ de chaque raie, et on en déduit la vitesse d'éloignement de la galaxie.

\begin{popfigure}
\begin{center}
\begin{tikzpicture}[scale=1]
% Spectre de référence
\draw[thick, popdark] (0,2.6) -- (8,2.6);
\node[left, popdark, font=\footnotesize\bfseries] at (0,2.6) {Référence (laboratoire)};
\foreach \x/\c in {1.2/poppurple, 2.6/popblue, 4.4/popgreen, 6.0/popgold}{
  \draw[line width=2.5pt, \c] (\x,2.25) -- (\x,2.95);
}
% Spectre galaxie
\draw[thick, popdark] (0,0.8) -- (8,0.8);
\node[left, popdark, font=\footnotesize\bfseries] at (0,0.8) {Galaxie lointaine};
\foreach \x/\c in {2.0/poppurple, 3.4/popblue, 5.2/popgreen, 6.8/popgold}{
  \draw[line width=2.5pt, \c] (\x,0.45) -- (\x,1.15);
}
% Flèches de décalage
\foreach \xa/\xb in {1.2/2.0, 2.6/3.4, 4.4/5.2, 6.0/6.8}{
  \draw[-{Stealth}, poppink, thick] (\xa,1.75) -- (\xb,1.75);
}
\node[poppink, font=\footnotesize\bfseries] at (4,1.45) {décalage vers le rouge (redshift)};
% Graduations
\node[popmuted, font=\footnotesize] at (1.0,-0.15) {bleu ($\lambda$ petit)};
\node[popmuted, font=\footnotesize] at (7.0,-0.15) {rouge ($\lambda$ grand)};
\draw[-{Stealth}, popmuted] (0.5,-0.45) -- (7.5,-0.45) node[midway, below, popmuted, font=\footnotesize] {longueur d'onde $\lambda$ croissante};
\end{tikzpicture}
\end{center}
\legende{Comparaison des raies d'absorption : dans le spectre d'une galaxie lointaine, chaque raie est décalée vers les grandes longueurs d'onde (vers le rouge), signe que la galaxie s'éloigne de nous.}
\end{popfigure}

\begin{savaistu}[Hubble, 1929 : l'Univers en expansion]
En mesurant le décalage vers le rouge de la lumière de nombreuses galaxies, l'astronome américain \textbf{Edwin Hubble} découvrit en 1929 que \emph{presque toutes} les galaxies s'éloignent de nous, et que leur vitesse d'éloignement est d'autant plus grande qu'elles sont lointaines (loi de Hubble). C'est la première preuve observationnelle de l'\textbf{expansion de l'Univers}, fondement de la cosmologie moderne et du modèle du Big Bang. Le même principe permet aujourd'hui de détecter des exoplanètes : l'étoile, entraînée par sa planète, oscille légèrement, et son spectre oscille entre redshift et blueshift !
\end{savaistu}

\courssub{Applications technologiques}

L'effet Doppler est partout dans la technologie moderne :

\begin{itemize}
\item \textbf{Radar routier} (contrôle de vitesse) : le radar émet une onde électromagnétique de fréquence $f_E$ vers le véhicule ; l'onde réfléchie revient avec une fréquence $f_R$ différente, car le véhicule est en mouvement. Le décalage $|f_R - f_E|$ donne directement la vitesse du véhicule.
\item \textbf{Échographie Doppler médicale} : dans les hôpitaux (par exemple au CHU de Yaoundé), on mesure le décalage de fréquence des ultrasons réfléchis par les globules rouges en mouvement pour évaluer le \textbf{débit sanguin} et détecter des anomalies cardiaques ou vasculaires.
\item \textbf{Radar météorologique} : le décalage Doppler des ondes réfléchies par les gouttes de pluie renseigne sur la \textbf{vitesse du vent} dans les nuages, utile pour la prévision des intempéries.
\end{itemize}

\begin{exemplebox}[Le radar routier]
Un radar fixe au bord de l'autoroute Yaoundé--Douala émet une onde de fréquence $f_E = \SI{24,125}{GHz}$. Une voiture qui s'approche renvoie l'onde avec une fréquence $f_R$ légèrement supérieure. L'électronique du radar compare $f_R$ et $f_E$ : le décalage mesuré, proportionnel à la vitesse $v$ du véhicule, permet d'afficher instantanément $v = \SI{112}{km.h^{-1}}$... et de dresser le procès-verbal !
\end{exemplebox}

\begin{exoresolu}[Ambulance et effet Doppler]
Une ambulance se déplace à $v_S = \SI{72}{km.h^{-1}}$ en ligne droite vers un observateur immobile. Sa sirène émet un son de fréquence $f_E = \SI{500}{Hz}$. On donne $v_{\text{son}} = \SI{340}{m.s^{-1}}$.

\textbf{1.} Convertir la vitesse de l'ambulance en \SI{}{m.s^{-1}}.

\textbf{2.} Sans calcul, indiquer si la fréquence perçue par l'observateur est supérieure ou inférieure à $\SI{500}{Hz}$. Justifier.

\textbf{3.} Vérifier, à l'aide de la formule de rapprochement, que la fréquence reçue vaut environ $f_R = \SI{531}{Hz}$.

\textbf{4.} Que devient qualitativement cette fréquence une fois que l'ambulance a dépassé l'observateur ?

\tcblower
\textbf{1.} Conversion :
\[ v_S = \dfrac{72}{3{,}6} = \SI{20}{m.s^{-1}}. \]

\textbf{2.} L'ambulance \emph{se rapproche} de l'observateur : les fronts d'onde se resserrent devant la source, la longueur d'onde apparente diminue, donc la fréquence reçue est \textbf{supérieure} à $f_E = \SI{500}{Hz}$ : le son est perçu plus aigu.

\textbf{3.} Application numérique :
\[ f_R = f_E \times \dfrac{v_{\text{son}}}{v_{\text{son}} - v_S} = 500 \times \dfrac{340}{340 - 20} = 500 \times \dfrac{340}{320} \approx \SI{531}{Hz}. \]
On retrouve bien $f_R > f_E$, cohérent avec le raisonnement qualitatif.

\textbf{4.} Après le dépassement, l'ambulance \emph{s'éloigne} : les fronts d'onde s'écartent derrière elle. La fréquence reçue devient alors \textbf{inférieure} à $\SI{500}{Hz}$ (la formule d'éloignement donne $f_R = 500 \times \dfrac{340}{360} \approx \SI{472}{Hz}$) : l'observateur entend brutalement un son plus grave. C'est exactement la transition « aigu $\to$ grave » que l'on perçoit au passage d'un véhicule d'urgence.
\end{exoresolu}

\begin{methode}[Raisonner sur un exercice d'effet Doppler (niveau Bac)]
\begin{enumerate}
\item \textbf{Identifier} la source, le récepteur, et le sens du mouvement relatif : rapprochement ou éloignement ?
\item \textbf{Conclure sur le sens du décalage} : rapprochement $\Rightarrow f_R > f_E$ (son aigu, blueshift en lumière) ; éloignement $\Rightarrow f_R < f_E$ (son grave, redshift).
\item \textbf{Exploiter la proportionnalité} : plus la vitesse relative est grande, plus le décalage est grand ; réciproquement, la mesure du décalage permet de calculer la vitesse.
\item \textbf{Vérifier la cohérence} : la fréquence émise $f_E$ ne change jamais ; seule la fréquence reçue est modifiée.
\end{enumerate}
\end{methode}

\begin{exercice}[Radar météorologique]
Un radar météorologique émet une onde électromagnétique vers un nuage de gouttes de pluie qui se rapproche de la station. L'onde réfléchie par les gouttes revient au radar avec une fréquence $f_R$.

\textbf{1.} Comparer qualitativement $f_R$ à la fréquence émise $f_E$. Justifier.

\textbf{2.} Le décalage mesuré est deux fois plus grand pour un second nuage que pour le premier. Que peut-on dire des vitesses de rapprochement des deux nuages ?
\end{exercice}

\begin{corrige}[Exercice 1]
\textbf{1.} Le nuage \emph{se rapproche} du radar : par effet Doppler, la fréquence de l'onde réfléchie est \textbf{supérieure} à la fréquence émise : $f_R > f_E$ (les fronts d'onde sont « resserrés » par le mouvement de rapprochement).

\textbf{2.} Le décalage $|f_R - f_E|$ est proportionnel à la vitesse relative de rapprochement. Un décalage deux fois plus grand signifie donc que le second nuage se rapproche \textbf{deux fois plus vite} que le premier.
\end{corrige}

\begin{aretenir}[L'essentiel de la section]
\begin{itemize}
\item L'\cle{effet Doppler} est le décalage de la fréquence perçue dû au \textbf{mouvement relatif} source--récepteur ; il concerne toutes les ondes.
\item Rapprochement $\Rightarrow$ fréquence reçue \textbf{plus élevée} ; éloignement $\Rightarrow$ fréquence reçue \textbf{plus faible}. La fréquence émise, elle, reste constante.
\item En lumière : éloignement $\Rightarrow$ décalage vers le \textbf{rouge} (redshift) ; rapprochement $\Rightarrow$ décalage vers le \textbf{bleu} (blueshift).
\item Le redshift des galaxies prouve l'\textbf{expansion de l'Univers} (Hubble, 1929).
\item Applications : radar routier, échographie Doppler, radar météorologique — le décalage mesuré donne accès à la \textbf{vitesse}.
\end{itemize}
\end{aretenir}

\newpage

\coursec{Activité d'intégration}

\coursec{Activité d'intégration : Diffraction, interférences et nature ondulatoire de la lumière}

\courssub{Objectifs de l'activité}

À l'issue de cette activité d'intégration, tu seras capable de :
\begin{itemize}
\item mobiliser simultanément les phénomènes de \cle{diffraction} et d'\cle{interférences} pour déterminer une même grandeur physique par deux méthodes indépendantes ;
\item exploiter la relation $L = \dfrac{2\lambda D}{a}$ issue de l'écart angulaire $\theta = \dfrac{\lambda}{a}$ ;
\item exploiter l'expression de l'interfrange $i = \dfrac{\lambda D}{a}$ dans un dispositif de fentes de Young ;
\item traiter une situation de \cle{coïncidence} de franges en lumière polychromatique ;
\item comparer deux résultats expérimentaux, calculer un \cle{écart relatif} et rédiger une conclusion argumentée sur la nature ondulatoire de la lumière.
\end{itemize}

\courssub{Situation}

Le laboratoire de physique du Lycée de Nkol-Eton, à Yaoundé, vient de recevoir un pointeur laser destiné aux travaux pratiques d'optique des classes de Terminale C. Malheureusement, l'étiquette du fabricant, qui indiquait la longueur d'onde $\lambda$ de la radiation émise, s'est décollée pendant le transport. Le professeur de physique propose à ses élèves de retrouver cette valeur par eux-mêmes, en exploitant les deux phénomènes ondulatoires étudiés dans la leçon : la diffraction et les interférences. Pour garantir la fiabilité du résultat, il impose que la mesure soit faite par \textbf{deux méthodes indépendantes}, puis comparée.

\textbf{Première expérience : diffraction par un fil.} Le faisceau laser éclaire un fil fin de diamètre $a_1 = \SI{80}{\micro m}$, placé perpendiculairement au faisceau. Sur un écran situé à la distance $D = \SI{2,00}{m}$ du fil, les élèves observent une figure de diffraction identique à celle d'une fente de même largeur (théorème de Babinet) : une tache centrale brillante, deux fois plus large que les taches latérales. Ils mesurent à la règle la largeur de la tache centrale : $L = \SI{3,2}{cm}$.

\begin{popfigure}
\begin{tikzpicture}[scale=1.5, >=Stealth]
% Laser source
\draw[popred, thick, ->] (-1, 0) -- (0, 0) node[midway, above] {Faisceau laser};
% Wire (obstacle)
\draw[fill=popdark] (-0.05, -0.3) rectangle (0.05, 0.3) node[midway, rotate=90, white, font=\small] {Fil $a_1$};
% Screen
\draw[thick, popblue] (2.5, -1) -- (2.5, 1) node[midway, right] {\'Ecran};
% Central bright spot (wire diffraction)
\draw[poporange, very thick] (2.5, -0.4) -- (2.5, 0.4) node[midway, right=2mm] {$L$};
\draw[<->, thin] (2.7, -0.4) -- (2.7, 0) node[midway, right] {$L/2$};
% Rays to first minima (edges of central spot)
\draw[dashed, thin, popmuted] (0, 0) -- (2.5, 0.4);
\draw[dashed, thin, popmuted] (0, 0) -- (2.5, -0.4);
% Angle theta
\draw[->, poppurple] (0.3, 0) arc (0:9:0.3) node[right=1mm] {$\theta$};
% Distance D
\draw[<->, thin] (-0.5, -0.7) -- (2.5, -0.7) node[midway, below] {$D$};
\end{tikzpicture}
\legende{Diffraction par un fil (équivalent fente par Babinet) : tache centrale brillante de largeur $L = 2D\tan\theta \approx 2\lambda D/a_1$.}
\end{popfigure}

\textbf{Deuxième expérience : fentes de Young.} Le même laser éclaire maintenant deux fentes fines parallèles, distantes de $a_2 = \SI{0,25}{mm}$. L'écran d'observation reste placé à $D = \SI{2,00}{m}$ des fentes. Les élèves observent des franges rectilignes, équidistantes, alternativement brillantes et sombres. Pour améliorer la précision, ils mesurent la longueur occupée par $8$ interfranges consécutifs : $\ell = \SI{4,0}{cm}$.

\begin{popfigure}
\begin{tikzpicture}[scale=1.2, >=Stealth]
% Slits
\draw[thick, popdark] (-0.1, 0.15) -- (0, 0.15) node[midway, left] {$S_1$};
\draw[thick, popdark] (-0.1, -0.15) -- (0, -0.15) node[midway, left] {$S_2$};
% Screen
\draw[thick, popblue] (3, -1) -- (3, 1) node[midway, right] {\'Ecran};
% Central axis
\draw[dashed, thin] (-0.5, 0) -- (3.5, 0);
% Point M (bright fringe order p)
\draw[poporange, thick] (3, 0.6) node[right] {$M$};
\draw[poporange, ->] (3, 0) -- (3, 0.6) node[midway, right] {$x$};
% Rays from slits to M
\draw[popgreen, ->] (0, 0.15) -- (3, 0.6);
\draw[popgreen, ->] (0, -0.15) -- (3, 0.6);
% Distance D
\draw[<->, thin] (-0.5, -0.5) -- (3, -0.5) node[midway, below] {$D$};
% Separation a
\draw[<->, thin] (0, 0.15) -- (0, -0.15) node[midway, left] {$a$};
% Angle theta
\draw[->, poppurple] (0.3, 0) arc (0:11:0.3) node[right=1mm] {$\theta$};
% Path difference construction
\draw[dashed, thin] (0, 0.15) -- (0.3, 0.15);
\draw[<->, thin, popred] (0.3, -0.15) -- (0.3, 0.15) node[midway, right] {$\delta \approx a\sin\theta$};
\end{tikzpicture}
\legende{Dispositif des fentes de Young : géométrie pour l'établissement de $\delta = \frac{a x}{D}$ et $i = \frac{\lambda D}{a}$.}
\end{popfigure}

\textbf{Troisième expérience : source polychromatique.} On remplace le laser par une source émettant deux radiations de longueurs d'onde connues : $\lambda_1 = \SI{650}{nm}$ (rouge) et $\lambda_2 = \SI{520}{nm}$ (verte). Le dispositif de fentes de Young est conservé. Les élèves constatent que les deux systèmes de franges se superposent et que, à certains endroits de l'écran, les franges brillantes des deux radiations coïncident exactement.

\textbf{Problème posé :} quelle est la longueur d'onde du laser ? Les deux méthodes donnent-elles des résultats cohérents ? Que conclure sur la nature de la lumière ?

\courssub{Consignes}

\begin{enumerate}
\item \textbf{Diffraction.} Rappeler l'expression de l'écart angulaire $\theta$ de la tache centrale de diffraction. En déduire, dans l'approximation des petits angles, la relation entre $L$, $\lambda$, $D$ et $a_1$, puis calculer la valeur $\lambda_{\text{dif}}$ de la longueur d'onde du laser.
\item \textbf{Interférences.} Déterminer la valeur de l'interfrange $i$ à partir de la mesure des $8$ interfranges. En déduire la valeur $\lambda_{\text{int}}$ de la longueur d'onde du laser.
\item \textbf{Comparaison.} Calculer l'écart relatif entre les deux valeurs obtenues. Conclure sur la cohérence des deux méthodes, sachant que l'écart relatif toléré en TP est de $5\,\%$.
\item \textbf{Coïncidences.} Écrire la condition de coïncidence des franges brillantes des deux radiations en fonction des ordres d'interférence $p_1$ et $p_2$. Déterminer les plus petits entiers $p_1$ et $p_2$ qui la réalisent, puis la plus petite différence de marche $\delta$ correspondante (hors la position centrale).
\item \textbf{Synthèse.} Rédiger une conclusion de quelques lignes expliquant pourquoi ces trois observations prouvent la nature ondulatoire de la lumière, et pourquoi un modèle purement corpusculaire ne peut pas les expliquer.
\end{enumerate}

\courssub{Ressources à mobiliser}

\begin{aretenir}[Boîte à outils de l'activité]
\begin{itemize}
\item \textbf{Diffraction} par une fente ou un fil de largeur $a$ : écart angulaire de la tache centrale $\theta = \dfrac{\lambda}{a}$ (en radians) ; largeur de la tache centrale sur l'écran : $L = 2D\tan\theta \approx \dfrac{2\lambda D}{a}$.
\item \textbf{Fentes de Young} : différence de marche $\delta = \dfrac{a x}{D}$ ; frange brillante si $\delta = p\lambda$ ($p \in \mathbb{Z}$), frange sombre si $\delta = \left(p + \tfrac{1}{2}\right)\lambda$.
\item \textbf{Interfrange} : $i = \dfrac{\lambda D}{a}$, distance entre deux franges consécutives de même nature.
\item \textbf{Coïncidence} en lumière polychromatique : les franges brillantes coïncident lorsque $p_1 \lambda_1 = p_2 \lambda_2$.
\item \textbf{Écart relatif} entre deux mesures : $\varepsilon = \dfrac{|\lambda_{\text{dif}} - \lambda_{\text{int}}|}{\lambda_{\text{moy}}}$, exprimé en pourcentage.
\end{itemize}
\end{aretenir}

\courssub{Démarche et corrigé commenté}

\begin{exoresolu}[Corrigé complet de l'activité]
\textbf{Énoncé rappelé.} Déterminer la longueur d'onde du laser par diffraction ($a_1 = \SI{80}{\micro m}$, $D = \SI{2,00}{m}$, $L = \SI{3,2}{cm}$) puis par interférences ($a_2 = \SI{0,25}{mm}$, $D = \SI{2,00}{m}$, $8$ interfranges sur $\ell = \SI{4,0}{cm}$), comparer les résultats, puis étudier les coïncidences pour $\lambda_1 = \SI{650}{nm}$ et $\lambda_2 = \SI{520}{nm}$.
\tcblower
\textbf{1. Exploitation de la diffraction.}

L'écart angulaire de la tache centrale (premiers minima) est $\theta = \dfrac{\lambda}{a_1}$. Sur l'écran, la demi-largeur de la tache vaut $D\tan\theta$. Comme l'angle est très petit (vérification \textit{a posteriori} : $\theta \approx \lambda/a_1 \sim 10^{-2}$ rad), on utilise l'approximation $\tan\theta \approx \theta$. La largeur totale de la tache centrale est donc :
\[ L = 2D\theta = \frac{2\lambda D}{a_1} \quad \Longrightarrow \quad \lambda_{\text{dif}} = \frac{L\,a_1}{2D}. \]
Vérification dimensionnelle : $\dfrac{[L]\times[a_1]}{[D]} = \dfrac{\text{m}\times\text{m}}{\text{m}} = \text{m}$, c'est bien une longueur. Application numérique, avec toutes les grandeurs en mètres (SI) :
\[ \lambda_{\text{dif}} = \frac{\num{3,2e-2} \times \num{80e-6}}{2 \times \num{2,00}} = \frac{\num{2,56e-6}}{\num{4,00}} = \SI{6,40e-7}{m} = \SI{640}{nm}. \]

\textbf{2. Exploitation des interférences.}

La mesure de $8$ interfranges sur $\ell = \SI{4,0}{cm}$ donne la valeur de l'interfrange :
\[ i = \frac{\ell}{8} = \frac{\SI{4,0}{cm}}{8} = \SI{0,50}{cm} = \SI{5,0e-3}{m}. \]
De l'expression de l'interfrange $i = \dfrac{\lambda D}{a_2}$, on tire la longueur d'onde :
\[ \lambda_{\text{int}} = \frac{i\,a_2}{D} = \frac{\num{5,0e-3} \times \num{0,25e-3}}{\num{2,00}} = \frac{\num{1,25e-6}}{\num{2,00}} = \SI{6,25e-7}{m} = \SI{625}{nm}. \]

\textbf{3. Comparaison des deux méthodes.}

La valeur moyenne est $\lambda_{\text{moy}} = \dfrac{640 + 625}{2} = \SI{632,5}{nm} \approx \SI{633}{nm}$. L'écart relatif vaut :
\[ \varepsilon = \frac{|640 - 625|}{632,5} = \frac{15}{632,5} \approx \num{0,0237} \approx 2,4\,\%. \]
Cet écart est inférieur à la tolérance de $5\,\%$ : les deux méthodes sont \textbf{cohérentes}. Le laser est très probablement un laser hélium-néon standard de longueur d'onde $\SI{632,8}{nm}$ (rouge). Les écarts résiduels proviennent des incertitudes de lecture à la règle (bords de la tache de diffraction peu nets, finesse des franges d'interférence).

\textbf{4. Coïncidences en lumière polychromatique.}

Une frange brillante de la radiation $\lambda_1$ d'ordre $p_1$ coïncide avec une frange brillante de la radiation $\lambda_2$ d'ordre $p_2$ lorsque les différences de marche sont égales :
\[ \delta = p_1 \lambda_1 = p_2 \lambda_2 \quad \Longrightarrow \quad \frac{p_1}{p_2} = \frac{\lambda_2}{\lambda_1} = \frac{520}{650} = \frac{4}{5}. \]
Les plus petits entiers solutions (hors $p_1 = p_2 = 0$, qui correspond à la frange centrale blanche) sont donc $p_1 = 4$ et $p_2 = 5$. La plus petite différence de marche de coïncidence est :
\[ \delta = 4 \times \SI{650}{nm} = 5 \times \SI{520}{nm} = \SI{2600}{nm} = \SI{2,6}{\micro m}. \]
La quatrième frange brillante rouge coïncide exactement avec la cinquième frange brillante verte : à cet endroit de l'écran, on observe une frange jaune-orangé (synthèse additive rouge + vert).

\textbf{5. Synthèse.}

La diffraction (étalement du faisceau derrière un fil, tache centrale élargie) et les interférences (franges alternativement brillantes et sombres) sont deux phénomènes \emph{spécifiquement ondulatoires}. Un modèle corpusculaire, où la lumière serait faite de particules se propageant en ligne droite, prédirait une simple ombre géométrique du fil et deux taches lumineuses derrière les fentes : il est incapable d'expliquer les zones sombres là où deux ondes se superposent (interférences destructrices) ni l'élargissement de la tache centrale. Le fait que deux expériences indépendantes, fondées sur deux phénomènes distincts, conduisent à la \emph{même} longueur d'onde (à $2,4\,\%$ près) conforte la validité du modèle ondulatoire : la lumière se comporte comme une onde périodique caractérisée par sa longueur d'onde $\lambda$.
\end{exoresolu}

\courssub{Prolongement : Effet Doppler (approche qualitative)}

Le programme exige une compréhension qualitative de l'effet Doppler pour les ondes sonores et lumineuses.

\begin{definition}[Effet Doppler]
L'\cle{effet Doppler} est la modification apparente de la fréquence (ou de la longueur d'onde) d'une onde perçue par un observateur lorsque la source et l'observateur sont en mouvement relatif l'un par rapport à l'autre.
\end{definition}

\begin{propriete}[Règles qualitatives]
\begin{itemize}
\item \textbf{Rapprochement} (source et observateur se rapprochent) : la fréquence perçue \textbf{augmente} (aigu pour le son, \textbf{bleu} pour la lumière : décalage vers le bleu).
\item \textbf{Eloignement} (source et observateur s'éloignent) : la fréquence perçue \textbf{diminue} (grave pour le son, \textbf{rouge} pour la lumière : décalage vers le rouge).
\end{itemize}
\end{propriete}

\begin{savaistu}[Applications au Cameroun et dans le monde]
\begin{itemize}
\item \textbf{Radar routier (SONAR/POLICE) :} Une onde radio émise par le radar est réfléchie par le véhicule. Le décalage de fréquence de l'écho permet de calculer la vitesse du véhicule (contrôle de vitesse sur l'autoroute Yaoundé-Douala).
\item \textbf{Échographie Doppler (Hôpitaux) :} Mesure de la vitesse du sang dans les artères (détection de sténoses, suivi grossesse). Les ultrasons réfléchis sur les globules rouges subissent un décalage proportionnel à la vitesse du flux sanguin.
\item \textbf{Astronomie (Décalage vers le rouge) :} La lumière des galaxies lointaines présente un décalage systématique vers le rouge ($z > 0$). Cela prouve qu'elles s'éloignent de nous : l'Univers est en expansion (Loi de Hubble-Lemaître). C'est une preuve majeure du Big Bang.
\item \textbf{Télécommunications (MTN/Orange/CAMTEL) :} La compensation du Doppler est essentielle pour les liaisons satellites (orbite basse) et la 5G (mobilité rapide).
\end{itemize}
\end{savaistu}

\begin{attention}[Pièges à éviter pour le Bac]
\begin{itemize}
\item Ne pas confondre \textbf{vitesse de l'onde} (célérité $c$ ou $v$, constante dans un milieu donné) et \textbf{fréquence} ($f$, modifiée par l'effet Doppler).
\item Pour la lumière, l'effet Doppler est \emph{relativiste} : il ne dépend que de la vitesse relative source/observateur, pas d'un milieu de propagation (pas d'éther). La formule approchée pour $v \ll c$ est $\dfrac{\Delta \lambda}{\lambda} \approx \dfrac{v}{c}$.
\item Le décalage vers le rouge cosmologique n'est \emph{pas} un effet Doppler classique (vitesse dans l'espace) mais une conséquence de l'expansion de l'espace lui-même (Relativité Générale). Au programme de Terminale, on l'aborde comme une analogie pédagogique.
\end{itemize}
\end{attention}

\courssub{Bilan de l'activité}

Cette activité t'a permis de mettre en œuvre une démarche complète de physicien :
\begin{itemize}
\item \textbf{Deux mesures indépendantes d'une même grandeur.} La diffraction ($\lambda_{\text{dif}} = \SI{640}{nm}$) et les interférences ($\lambda_{\text{int}} = \SI{625}{nm}$) conduisent à des valeurs voisines : l'écart relatif de $2,4\,\%$, inférieur à la tolérance, valide les deux modèles théoriques $\theta = \lambda/a$ et $i = \lambda D/a$.
\item \textbf{L'intérêt des mesures multiples.} Mesurer $8$ interfranges au lieu d'un seul divise l'incertitude de lecture par $8$ : c'est une stratégie expérimentale à retenir pour le Bac et pour les TP.
\item \textbf{La coïncidence des franges} illustre concrètement la condition $p_1\lambda_1 = p_2\lambda_2$ : la plus petite différence de marche commune, $\delta = \SI{2,6}{\micro m}$, est le plus petit commun multiple « optique » des deux longueurs d'onde.
\item \textbf{L'effet Doppler} montre que la fréquence d'une onde n'est pas une propriété intrinsèque absolue mais dépend du référentiel d'observation, confirmant la nature ondulatoire (relation $c = \lambda f$).
\item \textbf{Une conclusion argumentée} : seule l'hypothèse ondulatoire rend compte simultanément de la diffraction, des interférences, des coïncidences et de l'effet Doppler. C'est ce type de raisonnement — observations, modèle, prédictions quantitatives vérifiées — qui a conduit Fresnel à imposer la théorie ondulatoire de la lumière au XIX\textsuperscript{e} siècle.
\end{itemize}

\begin{attention}[Pour aller plus loin sans se tromper]
Dans les calculs, convertis \emph{systématiquement} toutes les longueurs en mètres avant l'application numérique ($\SI{80}{\micro m} = \num{80e-6}{m}$, $\SI{0,25}{mm} = \num{2,5e-4}{m}$, $\SI{3,2}{cm} = \num{3,2e-2}{m}$, $\SI{4,0}{cm} = \num{4,0e-2}{m}$). La confusion entre micromètres, millimètres et nanomètres est la première cause d'erreur dans ce type d'exercice. Pense aussi à distinguer la largeur $L$ de la tache centrale (qui contient un facteur $2$) de l'écart angulaire $\theta$ (qui n'en contient pas). Vérifie toujours l'homogénéité dimensionnelle de tes formules.
\end{attention}

\newpage

\coursec{Méthodes et exercices résolus}

\coursec{Aspect ondulatoire de la lumière : diffraction, interférences et effet Doppler}

\courssub{La diffraction de la lumière}

\begin{definition}[Diffraction]
La \cle{diffraction} est l'étalement d'une onde lorsqu'elle rencontre un obstacle ou une ouverture de taille comparable à sa longueur d'onde. Pour la lumière, ce phénomène prouve qu'elle ne se propage pas strictement en ligne droite (optique géométrique) mais se comporte comme une onde.
\end{definition}

\begin{popfigure}
\begin{tikzpicture}[scale=1.2, >=Stealth]
% Slit
\draw[thick, popdark] (-0.1,-0.6) -- (-0.1,0.6);
\draw[thick, popdark] (0.1,-0.6) -- (0.1,0.6);
\draw[fill=popblueL] (-0.1,-0.6) rectangle (0.1,0.6);
\draw[<->] (-0.15,-0.6) -- (-0.15,0.6) node[midway, left] {$a$};
% Screen
\draw[thick, popdark] (4,-1.5) -- (4,1.5);
\draw[fill=popgray] (3.95,-1.5) rectangle (4.05,1.5);
% Central spot
\draw[fill=popyellow!80!poporange] (3.95,-0.8) rectangle (4.05,0.8);
\draw[<->] (4.1,-0.8) -- (4.1,0.8) node[midway, right] {$L$};
% Rays
\draw[popblue, thick] (0,0) -- (4,0) node[midway, above] {axe optique};
\draw[popblue, dashed] (0,0) -- (4,0.8);
\draw[popblue, dashed] (0,0) -- (4,-0.8);
% Angle theta
\draw[->] (0.3,0) arc (0:11.3:0.3) node[right, pos=0.5] {$\theta$};
% Distance D
\draw[<->] (-0.1,-1) -- (4,-1) node[midway, below] {$D$};
% Labels
\node at (0,-1.3) {Fente};
\node at (4,-1.8) {Écran};
\end{tikzpicture}
\legende{Dispositif de diffraction par une fente : géométrie de la tache centrale.}
\end{popfigure}

\begin{propriete}[Écart angulaire et largeur de la tache centrale]
Pour une fente de largeur $a$ éclairée par une onde plane monochromatique de longueur d'onde $\lambda$ ($\lambda \ll a$), l'écart angulaire (demi-largeur angulaire) de la tache centrale est :
\[ \theta = \frac{\lambda}{a} \qquad (\theta \text{ en radians}). \]
Sur un écran placé à la distance $D$ ($D \gg a$), la largeur linéaire $L$ de la tache centrale (distance entre les deux premiers minima) est :
\[ L = 2D\tan\theta \approx 2D\theta = \frac{2\lambda D}{a}. \]
\end{propriete}

\begin{aretenir}[Condition d'observation]
La diffraction est nette si la largeur de la fente $a$ est de l'ordre de grandeur de la longueur d'onde $\lambda$ (typiquement $a \sim 10^2 \text{ à } 10^3 \lambda$). Si $a \gg \lambda$, $\theta$ est très petit et la tache centrale se confond avec l'image géométrique de la fente.
\end{aretenir}

\begin{methode}[Décrire et modéliser le phénomène de diffraction]
\begin{enumerate}
\item \cle{Identifier} le diffracteur (fente de largeur $a$ ou fil de diamètre $a$) et la longueur d'onde $\lambda$ de la source (laser le plus souvent).
\item \cle{Observer} la figure : une tache centrale brillante, large et intense, encadrée de taches latérales plus étroites et moins lumineuses (franges de diffraction).
\item \cle{Écrire} la relation fondamentale pour l'écart angulaire : $\theta = \lambda/a$.
\item \cle{Relier} $\theta$ à la largeur $L$ sur l'écran à distance $D$ : $L \approx 2\lambda D/a$ (petits angles : $\tan\theta \approx \theta$).
\item \cle{Conclure} sur la nature ondulatoire : la diffraction est inexplicable par l'optique géométrique ; elle est une signature du comportement ondulatoire.
\end{enumerate}
\begin{attention}[Réflexe unités]
Convertis toujours $\lambda$ (souvent en nm) et $a$ (souvent en mm ou $\mu$m) en \textbf{mètres} avant tout calcul : $1\ \text{nm} = 10^{-9}\ \text{m}$, $1\ \text{mm} = 10^{-3}\ \text{m}$. L'angle $\theta$ s'exprime en radians (sans unité).
\end{attention}
\end{methode}

\begin{exoresolu}[Largeur d'une fente par diffraction (type Bac C)]
On éclaire une fente fine de largeur $a$ avec un laser hélium-néon de longueur d'onde $\lambda = 632,8\ \text{nm}$. L'écran d'observation est placé à $D = 2,00\ \text{m}$ de la fente, perpendiculairement au faisceau. On mesure la largeur de la tache centrale de diffraction : $L = 4,00\ \text{cm}$.
\begin{enumerate}
\item Faire un schéma du dispositif et y faire apparaître l'écart angulaire $\theta$.
\item Établir la relation entre $L$, $\lambda$, $D$ et $a$.
\item Calculer la largeur $a$ de la fente.
\item On remplace la fente par un cheveu de diamètre $d$. La tache centrale mesure alors $L' = 2,53\ \text{cm}$. Calculer $d$.
\end{enumerate}
\tcblower
\textbf{1. Schéma.} Voir figure ci-dessus. Le faisceau laser traverse la fente de largeur $a$ ; sur l'écran situé à la distance $D$, la tache centrale a pour largeur $L$. L'écart angulaire $\theta$ est l'angle entre le centre de la figure et le bord de la tache centrale (premier minimum).

\textbf{2. Relation.} Dans le triangle rectangle formé par le centre de la fente, le centre de la tache et le bord de la tache centrale :
\[ \tan\theta = \frac{L/2}{D} = \frac{L}{2D}. \]
La diffraction s'observe dans les petits angles, donc $\tan\theta \approx \theta$ (en rad). Or la théorie de la diffraction donne $\theta = \dfrac{\lambda}{a}$. On en déduit :
\[ \frac{L}{2D} = \frac{\lambda}{a} \quad\Longrightarrow\quad \boxed{L = \frac{2\lambda D}{a}}. \]

\textbf{3. Calcul de $a$.} On isole $a$ :
\[ a = \frac{2\lambda D}{L} = \frac{2 \times 632,8\times 10^{-9} \times 2,00}{4,00\times 10^{-2}}. \]
\[ a = \frac{2,5312\times 10^{-6}}{4,00\times 10^{-2}} = 6,33\times 10^{-5}\ \text{m} = 63,3\ \mu\text{m}. \]
La fente a une largeur $a \approx 63,3\ \mu\text{m}$, de l'ordre de cent fois la longueur d'onde : les conditions de la diffraction sont bien réunies.

\textbf{4. Diamètre du cheveu.} Un fil produit la même figure de diffraction qu'une fente de même largeur (théorème de Babinet) :
\[ d = \frac{2\lambda D}{L'} = \frac{2 \times 632,8\times 10^{-9} \times 2,00}{2,53\times 10^{-2}} = 1,00\times 10^{-4}\ \text{m}. \]
Le cheveu a un diamètre $d \approx 100\ \mu\text{m}$, valeur tout à fait réaliste.

\textbf{Conclusion :} la diffraction permet de mesurer optiquement, sans contact, des dimensions de l'ordre du dixième de millimètre : c'est une preuve expérimentale de la nature ondulatoire de la lumière.
\end{exoresolu}

\courssub{Montrer la nature ondulatoire de la lumière à partir des observations}

\begin{methode}[Rédiger une argumentation « nature ondulatoire »]
\begin{enumerate}
\item \cle{Décrire} l'observation : tache de diffraction étalée alors que l'optique géométrique prévoit une image de la fente ; ou franges d'interférences (zones alternativement brillantes et sombres).
\item \cle{Montrer l'échec du modèle corpusculaire} : si la lumière était un flux de particules se propageant en ligne droite, on observerait une tache de la taille de la fente (diffraction) ou simplement l'addition des éclairements (interférences), jamais des zones sombres là où deux faisceaux se superposent.
\item \cle{Invoquer le modèle ondulatoire} : la lumière est une onde ; elle subit la diffraction quand $a \sim \lambda$, et deux ondes lumineuses cohérentes se superposent selon le principe de superposition (interférences constructives ou destructives).
\item \cle{Conclure} : seuls les phénomènes de diffraction et d'interférences, propres aux ondes, expliquent les observations ; la lumière possède donc un aspect ondulatoire, caractérisé par sa longueur d'onde $\lambda$.
\end{enumerate}
\end{methode}

\begin{exoresolu}[Analyse d'une expérience de cours]
Au laboratoire du lycée, on réalise deux expériences avec un pointeur laser vert ($\lambda = 532\ \text{nm}$).
\begin{itemize}
\item Expérience A : le faisceau traverse une fente de $0,10\ \text{mm}$ ; sur l'écran, la tache observée est bien plus large que la fente et entourée de taches secondaires.
\item Expérience B : le faisceau éclaire deux fentes fines parallèles distantes de $a = 0,50\ \text{mm}$ ; sur l'écran à $D = 1,50\ \text{m}$, on observe une succession de franges brillantes et sombres équidistantes.
\end{itemize}
\begin{enumerate}
\item Pourquoi l'expérience A est-elle incompatible avec la propagation rectiligne de la lumière ?
\item Pourquoi l'existence de franges sombres dans l'expérience B est-elle inexplicable si la lumière était un flux de particules ?
\item Conclure sur la nature de la lumière.
\end{enumerate}
\tcblower
\textbf{1.} Si la lumière se propageait rigoureusement en ligne droite, l'écran montrerait une image géométrique de la fente : une tache rectangulaire de $0,10\ \text{mm}$ de large. Or la tache centrale mesure $L = \dfrac{2\lambda D}{a}$, très supérieure à $a$, et s'accompagne de taches latérales : le faisceau s'est \emph{étalé} après la fente. La propagation rectiligne est mise en défaut.

\textbf{2.} Un flux de particules lumineuses arrivant par les deux fentes produirait sur l'écran une simple superposition des deux éclairements : partout où de la lumière arrive, l'écran serait éclairé, et l'éclairement total serait la somme des éclairements. Il serait impossible d'obtenir des zones \emph{sombres} là où les deux faisceaux se superposent. Or on observe des franges noires : en certains points, « lumière + lumière = obscurité ». Seule une onde peut s'annuler par superposition (onde en opposition de phase : interférences destructives).

\textbf{3. Conclusion.} Diffraction et interférences sont des phénomènes caractéristiques des ondes. Leur observation avec la lumière démontre que celle-ci possède un \cle{aspect ondulatoire} : elle se comporte comme une onde électromagnétique de longueur d'onde $\lambda$.
\end{exoresolu}

\courssub{Interférences lumineuses : conditions et différence de marche}

\begin{popfigure}
\begin{tikzpicture}[scale=1.1, >=Stealth]
% Slits
\draw[thick, popdark] (-0.15,-0.25) -- (-0.15,0.25);
\draw[thick, popdark] (0.15,-0.25) -- (0.15,0.25);
\draw[fill=popblueL] (-0.15,-0.25) rectangle (0.15,0.25);
\node at (0,-0.5) {$S_1$ et $S_2$};
\draw[<->] (-0.15,-0.4) -- (0.15,-0.4) node[midway, below] {$a$};
% Screen
\draw[thick, popdark] (5,-2) -- (5,2);
\draw[fill=popgray] (4.95,-2) rectangle (5.05,2);
\node at (5,-2.3) {Écran};
% Point M
\draw[popblue, thick] (0,0) -- (5,1.2);
\draw[popblue, thick] (0,0.25) -- (5,1.2);
\draw[popblue, thick] (0,-0.25) -- (5,1.2);
\node[popblue] at (5.3,1.2) {$M(x)$};
% Axis
\draw[dashed] (0,0) -- (5,0) node[midway, above] {axe optique};
% Distances
\draw[<->] (-0.15,-1) -- (5,-1) node[midway, below] {$D$};
\draw[<->] (5.1,0) -- (5.1,1.2) node[midway, right] {$x$};
% Path difference construction
\draw[poporange, thick, dashed] (0,0.25) -- (0.2,0.25) -- (0.2,-0.25);
\draw[<->, poporange] (0.2,0) -- (0.2,0.25) node[midway, right] {$\delta$};
\node[poporange] at (0.5,0) {$\delta \approx \frac{ax}{D}$};
\end{tikzpicture}
\legende{Géométrie des fentes de Young : différence de marche $\delta$ en un point $M$.}
\end{popfigure}

\begin{propriete}[Différence de marche et conditions d'interférences]
Deux fentes $S_1$ et $S_2$ distantes de $a$, éclairées par une même source monochromatique (cohérence spatiale), produisent sur un écran distant de $D$ ($D \gg a$) des franges d'interférence.
En un point $M$ d'abscisse $x$ (repère origine au centre de la figure), la différence de marche est :
\[ \delta = S_2M - S_1M \approx \frac{ax}{D}. \]
\begin{itemize}
\item \textbf{Frange brillante (interférence constructive) :} $\delta = k\lambda$, avec $k \in \mathbb{Z}$ (ordre d'interférence).
\item \textbf{Frange sombre (interférence destructive) :} $\delta = \left(k + \frac{1}{2}\right)\lambda$, avec $k \in \mathbb{Z}$.
\end{itemize}
\end{propriete}

\begin{aretenir}[Position des franges et interfrange]
Les positions des franges brillantes sont $x_k = k\,\dfrac{\lambda D}{a}$.
L'\cle{interfrange} $i$ (distance entre deux franges brillantes consécutives, ou deux franges sombres consécutives) est :
\[ i = x_{k+1} - x_k = \frac{\lambda D}{a}. \]
\end{aretenir}

\begin{methode}[Établissement de l'interfrange $i = \lambda D / a$ (démonstration exigible au Bac)]
\begin{enumerate}
\item \cle{Schématiser} : deux fentes $S_1$ et $S_2$ distantes de $a$, écran à la distance $D$ ($D \gg a$), point $M$ d'abscisse $x$ sur l'écran.
\item \cle{Calculer la différence de marche} : avec $S_1M^2 = D^2 + \left(x + \tfrac{a}{2}\right)^2$ et $S_2M^2 = D^2 + \left(x - \tfrac{a}{2}\right)^2$, écrire
\[ S_1M^2 - S_2M^2 = (S_1M - S_2M)(S_1M + S_2M) = 2ax. \]
Comme $S_1M + S_2M \approx 2D$, il vient $\delta = S_1M - S_2M = \dfrac{ax}{D}$.
\item \cle{Traduire les conditions d'interférence} : frange brillante $\delta = k\lambda$ ; frange sombre $\delta = (k+\tfrac12)\lambda$.
\item \cle{Positionner les franges brillantes} : $x_k = \dfrac{k\lambda D}{a}$.
\item \cle{Conclure} : l'interfrange est $i = x_{k+1} - x_k = \dfrac{\lambda D}{a}$.
\end{enumerate}
\textbf{Réflexes d'exploitation :} pour mesurer $i$ avec précision, mesure la largeur $\ell$ de $n$ interfranges : $i = \ell/n$. L'ordre d'interférence en $M$ est $p = \dfrac{\delta}{\lambda} = \dfrac{ax}{\lambda D}$.
\end{methode}

\begin{exoresolu}[Fentes de Young : détermination d'une longueur d'onde (type Bac C)]
On réalise l'expérience des fentes de Young avec deux fentes $S_1$ et $S_2$ distantes de $a = 0,40\ \text{mm}$, éclairées par une lumière monochromatique. L'écran est placé à $D = 1,20\ \text{m}$. On mesure la distance entre la 4$^{\text{e}}$ frange brillante située d'un côté de la frange centrale et la 4$^{\text{e}}$ frange brillante située de l'autre côté : $\ell = 14,4\ \text{mm}$.
\begin{enumerate}
\item Établir l'expression de la différence de marche $\delta$ en un point $M$ d'abscisse $x$.
\item En déduire l'expression de l'interfrange $i$.
\item Calculer $i$ puis la longueur d'onde $\lambda$ de la lumière utilisée.
\item Quelle est la nature (brillante ou sombre) de la frange située à $x = 5,4\ \text{mm}$ de la frange centrale ?
\end{enumerate}
\tcblower
\textbf{1. Différence de marche.} $S_1M^2 = D^2 + \left(x + \dfrac{a}{2}\right)^2$ et $S_2M^2 = D^2 + \left(x - \dfrac{a}{2}\right)^2$. En soustrayant :
\[ (S_1M - S_2M)(S_1M + S_2M) = \left(x + \tfrac{a}{2}\right)^2 - \left(x - \tfrac{a}{2}\right)^2 = 2ax. \]
Comme $D \gg a$ et $D \gg x$, on a $S_1M + S_2M \approx 2D$, d'où :
\[ \boxed{\delta = S_1M - S_2M = \frac{ax}{D}}. \]

\textbf{2. Interfrange.} Franges brillantes : $\delta = k\lambda \Leftrightarrow x_k = \dfrac{k\lambda D}{a}$. Deux franges consécutives ($k$ et $k+1$) sont distantes de :
\[ \boxed{i = \frac{\lambda D}{a}}. \]

\textbf{3. Calculs.} Entre la frange d'ordre $-4$ et la frange d'ordre $+4$, il y a $8$ interfranges :
\[ i = \frac{\ell}{8} = \frac{14,4\ \text{mm}}{8} = 1,80\ \text{mm}. \]
On isole $\lambda$ dans l'expression de $i$ :
\[ \lambda = \frac{ia}{D} = \frac{1,80\times 10^{-3} \times 0,40\times 10^{-3}}{1,20} = 6,0\times 10^{-7}\ \text{m}. \]
La longueur d'onde est $\lambda = 600\ \text{nm}$ (lumière orange).

\textbf{4. Nature de la frange en $x = 5,4\ \text{mm}$.} Calculons l'ordre d'interférence :
\[ p = \frac{x}{i} = \frac{5,4}{1,80} = 3,0. \]
$p$ est un entier : $\delta = 3\lambda$, les ondes arrivent en phase. La frange en $x = 5,4\ \text{mm}$ est \cle{brillante} (c'est la 3$^{\text{e}}$ frange brillante).

\textbf{Conclusion :} la mesure d'une distance macroscopique ($\ell$) permet de déterminer une longueur microscopique ($\lambda$) : c'est tout l'intérêt du dispositif interférométrique.
\end{exoresolu}

\courssub{Interférences en lumière polychromatique et en lumière blanche}

\begin{popfigure}
\begin{tikzpicture}[scale=1.5]
% White light fringes schematic
% Central white fringe
\draw[fill=white, thick] (-0.2,-0.3) rectangle (0.2,0.3);
\node at (0,0) {\textbf{Blanc}};
% Colored fringes
\foreach \x/\col/\lab in {-0.5/violet/Violet, -0.7/blue/Bleu, -0.9/cyan/Cyan, -1.1/green/Vert, -1.3/yellow/Jaune, -1.5/orange/Orange, -1.7/red/Rouge,
                          0.5/violet/Violet, 0.7/blue/Bleu, 0.9/cyan/Cyan, 1.1/green/Vert, 1.3/yellow/Jaune, 1.5/orange/Orange, 1.7/red/Rouge} {
    \draw[fill=\col!80!white, thick] (\x-0.1,-0.3) rectangle (\x+0.1,0.3);
    \node[rotate=90, font=\tiny] at (\x,0) {\lab};
}
% Higher order blur
\draw[fill=popgray!50, thick, dashed] (-2.2,-0.3) rectangle (-1.9,0.3);
\draw[fill=popgray!50, thick, dashed] (1.9,-0.3) rectangle (2.2,0.3);
\node[rotate=90, font=\small] at (-2.05,0) {Blanc d'ordre sup.};
\node[rotate=90, font=\small] at (2.05,0) {Blanc d'ordre sup.};
\draw[<->] (-2.3,-0.5) -- (2.3,-0.5) node[midway, below] {$x$};
\node at (0,-0.8) {Frange centrale blanche ($k=0$)};
\end{tikzpicture}
\legende{Aspect qualitatif des franges en lumière blanche : frange centrale blanche, franges irisées (violet côté centre), puis blanc d'ordre supérieur.}
\end{popfigure}

\begin{methode}[Coïncidence des franges de deux radiations]
\begin{enumerate}
\item \cle{Écrire} les positions des franges brillantes de chaque radiation : $x_{k_1} = \dfrac{k_1\lambda_1 D}{a}$ et $x_{k_2} = \dfrac{k_2\lambda_2 D}{a}$ (les interfranges $i_1 = \dfrac{\lambda_1 D}{a}$ et $i_2 = \dfrac{\lambda_2 D}{a}$ sont différentes).
\item \cle{Traduire la coïncidence} : une frange brillante de chaque système se superpose en $x$ si
\[ k_1\lambda_1 = k_2\lambda_2, \qquad k_1, k_2 \in \mathbb{Z}. \]
\item \cle{Chercher les plus petits entiers} vérifiant la condition (la frange centrale $k_1 = k_2 = 0$ est toujours une coïncidence, commune à toutes les radiations).
\item \cle{En lumière blanche} : la frange centrale est blanche (toutes les radiations coïncident en $x = 0$) ; de part et d'autre, les franges sont irisées car les systèmes de franges se décalent, puis le blanc d'ordre supérieur apparaît quand les franges se recouvrent.
\end{enumerate}
\end{methode}

\begin{exoresolu}[Lampe à vapeur de mercure : coïncidence de deux raies]
On éclaire des fentes de Young ($a = 0,50\ \text{mm}$, $D = 1,00\ \text{m}$) avec une lumière contenant deux radiations : le violet $\lambda_1 = 436\ \text{nm}$ et le vert $\lambda_2 = 546\ \text{nm}$.
\begin{enumerate}
\item Calculer les interfranges $i_1$ et $i_2$ des deux systèmes de franges.
\item À quelle distance minimale $x$ de la frange centrale (hors $x = 0$) une frange brillante violette coïncide-t-elle avec une frange brillante verte ?
\item Décrire qualitativement l'aspect de l'écran si on utilise de la lumière blanche.
\end{enumerate}
\tcblower
\textbf{1. Interfranges.}
\[ i_1 = \frac{\lambda_1 D}{a} = \frac{436\times 10^{-9} \times 1,00}{0,50\times 10^{-3}} = 8,72\times 10^{-4}\ \text{m} = 0,872\ \text{mm}. \]
\[ i_2 = \frac{\lambda_2 D}{a} = \frac{546\times 10^{-9} \times 1,00}{0,50\times 10^{-3}} = 1,092\times 10^{-3}\ \text{m} = 1,092\ \text{mm}. \]

\textbf{2. Coïncidence.} Il faut $k_1\lambda_1 = k_2\lambda_2$, soit :
\[ \frac{k_1}{k_2} = \frac{\lambda_2}{\lambda_1} = \frac{546}{436} = \frac{273}{218}. \]
$273 = 3\times 7\times 13$ et $218 = 2\times 109$ sont premiers entre eux : les plus petits entiers sont $k_1 = 273$ et $k_2 = 218$. Alors :
\[ x = k_1 i_1 = 273 \times 0,872\ \text{mm} \approx 238\ \text{mm}. \]
Vérification : $x = k_2 i_2 = 218 \times 1,092 \approx 238\ \text{mm}$. La première coïncidence hors centre a lieu à $x \approx 23,8\ \text{cm}$ : en pratique, les deux systèmes de franges se brouillent bien avant, car les franges proches du centre ne coïncident jamais exactement.

\textbf{3. Lumière blanche.} Chaque radiation du spectre donne son propre système de franges, d'interfrange proportionnelle à $\lambda$. En $x = 0$, toutes les franges brillantes coïncident : la \cle{frange centrale est blanche}. En s'éloignant, les franges se décalent : on observe des franges \cle{irisées} (teintées, bordées de violet côté centre car $\lambda_{\text{violet}} < \lambda_{\text{rouge}}$), puis, lorsque le décalage devient de l'ordre d'une demi-interfrange pour certaines radiations, les franges se brouillent en un éclairement blanchâtre appelé \emph{blanc d'ordre supérieur}.

\textbf{Conclusion :} la coïncidence exige que le rapport des longueurs d'onde soit un rapport d'entiers ; la frange centrale blanche est le repère incontournable des interférences en lumière blanche.
\end{exoresolu}

\begin{savaistu}[Le spectromètre à réseau et l'astrophysique au Cameroun]
Les réseaux de diffraction (grande série de fentes) remplacent les fentes de Young en spectroscopie professionnelle. Ils dispersent la lumière selon la relation $a\sin\theta = k\lambda$ (plus précise que $i=\lambda D/a$). Au Cameroun, l'analyse spectrale est cruciale : contrôle de la qualité de l'eau (laboratoires SNV/ANOR), diagnostic médical (dosage sanguin), et bientôt astronomie avec le projet de radiotélescope africain (African VLBI Network) qui utilisera l'effet Doppler sur la raie à 21 cm de l'hydrogène pour cartographier la Voie lactée.
\end{savaistu}

\courssub{Effet Doppler : étude qualitative et applications}

\begin{popfigure}
\begin{tikzpicture}[scale=1.2, >=Stealth]
% Source moving right -> Observer on right
% Wavefronts compressed on right
\foreach \r/\c in {0.3/popblue, 0.6/popblue, 0.9/popblue, 1.2/popblue, 1.5/popblue} {
    \draw[\c, thick] (2,0) circle (\r);
}
% Wavefronts stretched on left
\foreach \r/\c in {0.3/popred, 0.7/popred, 1.1/popred, 1.5/popred} {
    \draw[\c, thick] (2,0) circle (\r);
}
% Source
\draw[fill=poporange, thick] (2,0) circle (0.15) node {S};
\draw[->, thick, poporange] (2.2,0) -- (2.8,0) node[midway, above] {$v_S$};
% Observer right
\node[draw, circle, thick, fill=popgreenL, minimum size=0.8cm] at (4.5,0) {O};
\node at (4.5,-0.7) {Récepteur};
% Observer left
\node[draw, circle, thick, fill=popgreenL, minimum size=0.8cm] at (-0.5,0) {O'};
\node at (-0.5,-0.7) {Récepteur};
% Labels
\node[popblue] at (3.5,1.5) {Fronts resserrés : $\lambda' \downarrow, f' \uparrow$};
\node[popred] at (0.5,1.5) {Fronts étirés : $\lambda' \uparrow, f' \downarrow$};
\end{tikzpicture}
\legende{Effet Doppler : source mobile $S$ s'approchant de $O$ (fréquence perçue plus aiguë) et s'éloignant de $O'$ (fréquence perçue plus grave).}
\end{popfigure}

\begin{propriete}[Effet Doppler (qualitatif)]
Lorsqu'une source d'ondes et un récepteur sont en mouvement relatif l'un par rapport à l'autre (ou par rapport au milieu de propagation pour les ondes mécaniques), la fréquence perçue $f'$ diffère de la fréquence émise $f$ :
\begin{itemize}
\item \textbf{Rapprochement :} les fronts d'onde s'accumulent $\Rightarrow$ longueur d'onde perçue $\lambda'$ diminue, fréquence $f'$ augmente (son plus \emph{aigu}, lumière décalée vers le \emph{bleu}).
\item \textbf{Eloignement :} les fronts d'onde s'espacent $\Rightarrow$ $\lambda'$ augmente, $f'$ diminue (son plus \emph{grave}, lumière décalée vers le \emph{rouge}).
\end{itemize}
\end{propriete}

\begin{methode}[Raisonnement sur l'effet Doppler]
\begin{enumerate}
\item \cle{Identifier} la source (émetteur) et le récepteur, et préciser lequel est en mouvement par rapport au milieu (son) ou relativement à l'autre (lumière).
\item \cle{Raisonner sur les fronts d'onde} : source qui s'approche $\rightarrow$ fronts resserrés $\rightarrow$ $\lambda \downarrow, f \uparrow$ ; source qui s'éloigne $\rightarrow$ fronts étirés $\rightarrow$ $\lambda \uparrow, f \downarrow$.
\item \cle{Citer des applications} :
\begin{itemize}
\item \textbf{Son :} sirène d'ambulance, radar routier (gendarmerie sur l'axe Yaoundé-Douala), échographie Doppler (hôpitaux de Yaoundé/Douala pour débit sanguin), radar météorologique (prévision pluies).
\item \textbf{Lumière :} décalage vers le rouge (\emph{redshift}) des galaxies $\rightarrow$ expansion de l'Univers (loi de Hubble) ; spectroscopie stellaire (vitesse radiale d'étoiles, détection d'exoplanètes).
\end{itemize}
\item \cle{Ne pas quantifier} : au programme de Terminale C, l'étude est qualitative ; aucune formule de décalage n'est exigible.
\end{enumerate}
\end{methode}

\begin{exoresolu}[De la sirène d'ambulance aux galaxies lointaines]
\begin{enumerate}
\item Une ambulance traverse Douala en émettant une sirène de fréquence constante. Un piéton immobile entend un son plus aigu quand le véhicule s'approche, puis plus grave dès qu'il le dépasse. Expliquer ce phénomène.
\item Le spectre de la lumière d'une galaxie lointaine présente les raies caractéristiques de l'hydrogène, mais décalées vers de grandes longueurs d'onde. Que peut-on en déduire ?
\item Le radar routier utilisé par la gendarmerie émet une onde électromagnétique vers un véhicule et compare la fréquence de l'onde réfléchie à celle de l'onde émise. Expliquer le principe de la mesure de vitesse.
\end{enumerate}
\tcblower
\textbf{1. Sirène.} La sirène émet des ondes sonores sphériques à fréquence fixe. Quand l'ambulance \emph{se rapproche} du piéton, chaque front d'onde est émis un peu plus près de lui que le précédent : les fronts se resserrent, la longueur d'onde perçue diminue et, la célérité du son dans l'air étant constante ($c = \lambda f$), la fréquence perçue \emph{augmente} : le son paraît plus aigu. Quand l'ambulance \emph{s'éloigne}, les fronts s'espacent : la fréquence perçue diminue, le son paraît plus grave. C'est l'\cle{effet Doppler}.

\textbf{2. Galaxie.} Les raies de l'hydrogène sont décalées vers les grandes longueurs d'onde, c'est-à-dire vers le rouge (\emph{redshift}) : d'après l'effet Doppler lumineux, la galaxie \emph{s'éloigne} de la Terre. Ce décalage, observé pour quasiment toutes les galaxies et d'autant plus grand qu'elles sont lointaines (loi de Hubble), est la preuve observationnelle de l'\cle{expansion de l'Univers}.

\textbf{3. Radar routier.} L'onde émise se réfléchit sur le véhicule en mouvement, qui joue tour à tour le rôle de récepteur en mouvement puis de source en mouvement : l'onde réfléchie revient avec une fréquence légèrement différente de la fréquence émise. L'écart de fréquence (battement) est d'autant plus grand que le véhicule roule vite : l'appareil, étalonné, en déduit directement la vitesse. Le même principe sert en médecine (échographie Doppler pour mesurer la vitesse du sang) et en météorologie (radars de précipitations).

\textbf{Conclusion :} l'effet Doppler relie le décalage de fréquence perçue au mouvement relatif source-récepteur ; c'est un outil universel de mesure de vitesse, de la route aux galaxies.
\end{exoresolu}

\courssub{Synthèse : choisir et combiner les bonnes formules}

\begin{aretenir}[Formulaire récapitulatif]
\begin{center}
\begin{tabularx}{\linewidth}{|l|X|X|}\hline
\textbf{Phénomène} & \textbf{Grandeurs clés} & \textbf{Formules} \\ \hline
Diffraction (fente/fil) & $a$ : largeur fente ou diam. fil & $\theta = \dfrac{\lambda}{a}$ \\ 
& $L$ : largeur tache centrale & $L = \dfrac{2\lambda D}{a}$ \\ \hline
Interférences (Young) & $a$ : distance entre fentes & $\delta = \dfrac{ax}{D}$ \\
& $i$ : interfrange & $i = \dfrac{\lambda D}{a}$ \\
& Frange brillante & $\delta = k\lambda \ (k\in\mathbb{Z})$ \\
& Frange sombre & $\delta = (k+\tfrac12)\lambda$ \\ \hline
Coïncidence (2 $\lambda$) & $k_1, k_2$ entiers & $k_1\lambda_1 = k_2\lambda_2$ \\ \hline
\end{tabularx}
\end{center}
\end{aretenir}

\begin{methode}[Réflexes de synthèse diffraction-interférences]
\begin{enumerate}
\item \cle{Faire le tri} dans les données : diffraction $\Rightarrow$ $\theta = \lambda/a$, $L = 2\lambda D/a$ ; interférences $\Rightarrow$ $\delta = ax/D$, $i = \lambda D/a$, $\delta = k\lambda$ (brillante) ou $\delta = (k+\tfrac12)\lambda$ (sombre).
\item \cle{Attention au sens de $a$} : en diffraction, $a$ est la largeur de la fente ; en interférences, $a$ est la \emph{distance entre les deux fentes}. Ce ne sont pas les mêmes grandeurs !
\item \cle{Vérifier l'homogénéité} : $\lambda D/a$ est bien homogène à une longueur ($\text{m}\times\text{m}/\text{m}$).
\item \cle{Exploiter un comptage de franges} : $n$ interfranges occupent la longueur $\ell = n\,i$.
\item \cle{Conclure chaque question} par une phrase avec valeur numérique, unité SI et chiffres significatifs cohérents avec les données.
\end{enumerate}
\end{methode}

\begin{exoresolu}[Exercice de synthèse : du laser aux fentes de Young]
Un professeur utilise un laser de longueur d'onde inconnue $\lambda$.
\begin{itemize}
\item \textbf{Partie 1 :} il place devant le faisceau une fente de largeur $a_1 = 0,080\ \text{mm}$ ; sur un écran à $D_1 = 1,50\ \text{m}$, la tache centrale de diffraction a pour largeur $L = 2,4\ \text{cm}$.
\item \textbf{Partie 2 :} il remplace la fente par deux fentes de Young distantes de $a_2 = 0,25\ \text{mm}$, l'écran étant à $D_2 = 2,00\ \text{m}$.
\end{itemize}
\begin{enumerate}
\item Calculer $\lambda$ à partir de la partie 1.
\item Calculer l'interfrange $i$ de la figure d'interférences.
\item Combien d'interfranges compte-t-on sur une largeur de $2,08\ \text{cm}$ de l'écran ?
\item Déterminer la différence de marche en un point situé à $x = 10,4\ \text{mm}$ de la frange centrale et préciser la nature de la frange en ce point.
\end{enumerate}
\tcblower
\textbf{1. Longueur d'onde.} Diffraction : $L = \dfrac{2\lambda D_1}{a_1}$, donc
\[ \lambda = \frac{L\,a_1}{2D_1} = \frac{2,4\times 10^{-2} \times 0,080\times 10^{-3}}{2 \times 1,50} = \frac{1,92\times 10^{-6}}{3,00} = 6,4\times 10^{-7}\ \text{m}. \]
$\lambda = 640\ \text{nm}$ (lumière rouge).

\textbf{2. Interfrange.}
\[ i = \frac{\lambda D_2}{a_2} = \frac{6,4\times 10^{-7} \times 2,00}{0,25\times 10^{-3}} = 5,12\times 10^{-3}\ \text{m} = 5,12\ \text{mm}. \]

\textbf{3. Nombre d'interfranges.}
\[ n = \frac{\ell}{i} = \frac{20,8\ \text{mm}}{5,12\ \text{mm}} \approx 4,06. \]
On compte donc environ \cle{4 interfranges} sur cette largeur (soit 5 franges brillantes si la frange centrale est au bord, 4 ou 5 selon la position).

\textbf{4. Différence de marche en $x = 10,4\ \text{mm}$.}
\[ \delta = \frac{a_2 x}{D_2} = \frac{0,25\times 10^{-3} \times 10,4\times 10^{-3}}{2,00} = 1,30\times 10^{-6}\ \text{m}. \]
Ordre d'interférence : $p = \dfrac{\delta}{\lambda} = \dfrac{1,30\times 10^{-6}}{6,4\times 10^{-7}} \approx 2,03 \approx 2$. On a $\delta \approx 2\lambda$, multiple entier de $\lambda$ : les ondes sont en phase, la frange est \cle{brillante} (2$^{\text{e}}$ frange brillante, cohérent avec $x \approx 2i = 10,24\ \text{mm}$).

\textbf{Conclusion :} les deux phénomènes, diffraction et interférences, exploitent la même longueur d'onde et confirment l'aspect ondulatoire de la lumière ; la cohérence des résultats ($x \approx 2i$) valide l'ensemble des calculs.
\end{exoresolu}

\begin{attention}[Les 5 erreurs qui coûtent des points au Bac]
\begin{enumerate}
\item \textbf{Oublier les conversions.} $\lambda$ en nm, $a$ en mm, $i$ en mm : tout doit passer en \emph{mètres} avant le calcul. Une longueur d'onde de $600\ \text{nm}$ n'est pas $600\ \text{m}$ ! Écris systématiquement la conversion sur ta copie.
\item \textbf{Confondre les deux significations de $a$.} Dans $\theta = \lambda/a$, $a$ est la largeur de la fente diffractante ; dans $i = \lambda D/a$, $a$ est la \emph{distance entre les deux fentes} de Young. Mélanger les deux formules dans le mauvais contexte est l'erreur la plus fréquente.
\item \textbf{Oublier le facteur 2 en diffraction.} L'écart angulaire $\theta = \lambda/a$ est la \emph{demi}-largeur angulaire de la tache centrale : la largeur totale sur l'écran est $L = 2D\theta = \dfrac{2\lambda D}{a}$. Utiliser $L = \lambda D/a$ donne un résultat faux d'un facteur 2.
\item \textbf{Mal compter les interfranges.} Entre la frange d'ordre $-n$ et la frange d'ordre $+n$, il y a $2n$ interfranges, pas $2n+1$ ni $n$. De même, $N$ franges délimitent $N-1$ interfranges. Toujours préciser le comptage dans la rédaction.
\item \textbf{Écrire les conditions d'interférence de travers.} Frange \emph{brillante} : $\delta = k\lambda$ avec $k \in \mathbb{Z}$ (et non $k \in \mathbb{N}$ : les ordres négatifs existent !) ; frange \emph{sombre} : $\delta = (k + \tfrac{1}{2})\lambda$. Inverser les deux conditions, ou oublier de préciser que $k$ est un entier relatif, fait perdre des points de rigueur. Ajoute enfin que les interférences n'existent qu'entre sources \emph{cohérentes} (issues de la même source ponctuelle) : c'est la condition d'application souvent oubliée.
\end{enumerate}
\end{attention}

\newpage

\coursec{Exercices}

\courssub{Je vérifie mes connaissances}

\begin{exercice}[QCM : notions fondamentales]
Parmi les affirmations suivantes, une seule est correcte. Laquelle ?
\begin{enumerate}
    \item L'écart angulaire de la tache centrale de diffraction par une fente de largeur $a$ éclairée par une onde de longueur d'onde $\lambda$ vaut $\theta = \dfrac{a}{\lambda}$.
    \item Dans le dispositif des fentes de Young, les franges d'interférence sont localisées au plan des fentes.
    \item L'interfrange $i$ sur un écran placé à une distance $D$ des fentes d'écartement $a$, éclairées par une lumière monochromatique de longueur d'onde $\lambda$, s'exprime par $i = \dfrac{\lambda D}{a}$.
    \item L'effet Doppler-Fizeau pour la lumière ne dépend que de la vitesse relative de la source par rapport à l'observateur, et non de la vitesse de la lumière dans le vide.
\end{enumerate}
\end{exercice}

\begin{exercice}[Vrai ou Faux ? Justifiez]
Dire si les affirmations suivantes sont vraies ou fausses, en justifiant brièvement votre réponse.
\begin{enumerate}
    \item Si l'on diminue la largeur $a$ d'une fente éclairée par un laser, la tache centrale de diffraction observée sur un écran éloigné devient plus étroite.
    \item En lumière blanche, la frange centrale (ordre 0) observée sur l'écran des fentes de Young est blanche, tandis que les franges d'ordres non nuls sont colorées.
    \item Pour un observateur fixe, une source lumineuse qui s'éloigne voit son spectre décalé vers le bleu (longueurs d'onde plus courtes).
    \item La condition d'obtention d'une frange brillante d'ordre $k$ aux fentes de Young est $\delta = k \lambda$, où $\delta$ est la différence de marche.
\end{enumerate}
\end{exercice}

\begin{exercice}[Définitions et calculs directs]
\begin{enumerate}
    \item Définir l'écart angulaire $\theta$ de la tache centrale de diffraction par une fente de largeur $a$. Donner sa relation avec $\lambda$ et $a$.
    \item Une fente de largeur $a = \SI{0,20}{mm}$ est éclairée par un laser He-Ne ($\lambda = \SI{632,8}{nm}$). Calculer l'écart angulaire $\theta$ (en radians et en degrés) de la tache centrale.
    \item Dans un montage de Young, les fentes sont écartées de $a = \SI{0,50}{mm}$ et l'écran est à $D = \SI{2,00}{m}$. La lumière utilisée a une longueur d'onde $\lambda = \SI{589}{nm}$ (lampe sodium). Calculer l'interfrange $i$ en millimètres.
    \item Quelle est la différence de marche $\delta$ pour la troisième frange brillante ($k=3$) ? Et pour la deuxième frange sombre ?
\end{enumerate}
\end{exercice}

\begin{exercice}[Identification d'ordres d'interférence]
On observe sur un écran les franges d'interférence obtenues avec le dispositif des fentes de Young ($a = \SI{0,40}{mm}$, $D = \SI{1,50}{m}$, $\lambda = \SI{600}{nm}$). Un point M de l'écran se trouve à une distance $x_M = \SI{4,5}{mm}$ de la frange centrale.
\begin{enumerate}
    \item Calculer l'interfrange $i$.
    \item Déterminer l'ordre d'interférence $k$ correspondant au point M. Est-ce une frange brillante ou sombre ?
    \item Calculer la différence de marche $\delta_M$ en ce point.
\end{enumerate}
\end{exercice}

\courssub{Je m'entraîne}

\begin{exercice}[Diffraction par une fente : détermination de $\lambda$ et prévision]
Un faisceau laser traverse une fente unique de largeur $a_1 = \SI{0,15}{mm}$. Sur un écran placé à $D = \SI{2,50}{m}$ de la fente, on mesure la largeur de la tache centrale (distance entre les deux premiers minima) : $L_1 = \SI{22,0}{mm}$.
\begin{enumerate}
    \item En supposant les petits angles, établir la relation entre $L_1$, $D$, $\lambda$ et $a_1$.
    \item Déduire la longueur d'onde $\lambda$ du laser (en nm).
    \item On remplace la fente par une fente deux fois plus fine ($a_2 = a_1/2$). Prédire la nouvelle largeur $L_2$ de la tache centrale. Vérifier par le calcul.
\end{enumerate}
\end{exercice}

\begin{exercice}[Fentes de Young : analyse complète]
On utilise un dispositif des fentes de Young avec $a = \SI{0,60}{mm}$ et $D = \SI{3,00}{m}$. La source est une lampe à vapeur de mercure filtrée pour ne garder que la raie verte ($\lambda = \SI{546,1}{nm}$).
\begin{enumerate}
    \item Calculer l'interfrange $i$ (en mm).
    \item Déterminer l'abscisse $x_4$ de la frange brillante d'ordre 4.
    \item Déterminer l'abscisse $x'_{3}$ de la frange sombre d'ordre 3 (la troisième frange sombre à partir du centre).
    \item Calculer la distance entre la frange brillante d'ordre 2 et la frange sombre d'ordre 3.
    \item Si l'on immerge tout le dispositif dans l'eau ($n = 1,33$), que devient l'interfrange ? Calculer la nouvelle valeur.
\end{enumerate}
\end{exercice}

\begin{exercice}[Mesure de l'écartement des fentes et incertitude]
Pour déterminer l'écartement $a$ entre les deux fentes de Young, on mesure la distance $L_N$ qui sépare la frange brillante d'ordre 0 de la frange brillante d'ordre $N$ sur l'écran. On relève : $N = 20$, $L_{20} = \SI{36,4}{mm}$, $D = \SI{2,000 \pm 0,005}{m}$, $\lambda = \SI{632,8 \pm 0,1}{nm}$.
\begin{enumerate}
    \item Exprimer $a$ en fonction de $N$, $L_N$, $D$ et $\lambda$.
    \item Calculer la valeur de $a$ (en mm).
    \item Calculer l'incertitude relative sur $a$ ($\frac{\Delta a}{a}$) en supposant que l'incertitude sur $L_N$ est $\pm \SI{0,2}{mm}$ et que les incertitudes sont indépendantes.
    \item Donner le résultat final sous la forme $a = a_0 \pm \Delta a$ (en mm, avec un nombre approprié de chiffres significatifs).
\end{enumerate}
\end{exercice}

\begin{exercice}[Interférences en lumière polychromatique]
On éclaire les fentes de Young ($a = \SI{0,50}{mm}$, $D = \SI{2,00}{m}$) avec une source bichromatique émettant deux longueurs d'onde : $\lambda_1 = \SI{589,0}{nm}$ (jaune, raie D du sodium) et $\lambda_2 = \SI{460,0}{nm}$ (bleu).
\begin{enumerate}
    \item Calculer les interfranges $i_1$ et $i_2$ pour chaque composante.
    \item Déterminer les ordres $k_1$ et $k_2$ des premières franges brillantes (non nulles) qui coïncident. Quelle est la différence de marche $\delta$ correspondante ?
    \item Calculer l'abscisse $x_c$ de cette première coïncidence sur l'écran.
    \item Décrire qualitativement l'aspect de la figure d'interférence si l'on utilise une source de lumière blanche (spectre continu). Quelle est la couleur de la frange centrale ?
\end{enumerate}
\end{exercice}

\courssub{Je me prépare au Bac}

\begin{exercice}[Bac Type : Nature ondulatoire de la lumière et expansion de l'Univers]
\textbf{Contexte :} Lors d'une séance de travaux pratiques au Lycée de Nkolbisson, les élèves observent deux phénomènes :
\begin{itemize}
    \item \textbf{Expérience 1 :} Un fil de cheveux (diamètre $d \approx \SI{70}{\mu m}$) est placé sur le trajet d'un faisceau laser ($\lambda = \SI{650}{nm}$). Sur un écran à $D = \SI{2,00}{m}$, on observe une figure de diffraction composée d'une tache centrale brillante encadrée de franges alternatives sombres et brillantes. La largeur de la tache centrale mesurée est $L = \SI{37,1}{mm}$.
    \item \textbf{Expérience 2 :} Le même laser éclaire deux fentes jeunes très fines, écartées de $a = \SI{0,30}{mm}$. L'écran est toujours à $D = \SI{2,00}{m}$. On observe des franges d'interférence régulières d'interfrange $i = \SI{4,33}{mm}$.
\end{itemize}
\textbf{Questions :}
\begin{enumerate}
    \item \textbf{Diffraction par le fil :} Modéliser le fil comme une fente de largeur $a = d$. Vérifier que la mesure de $L$ est cohérente avec la théorie de la diffraction (petits angles). En déduire une valeur expérimentale de $\lambda$ et la comparer à la valeur nominale.
    \item \textbf{Interférences de Young :} Vérifier que la valeur de l'interfrange $i$ mesurée est cohérente avec la longueur d'onde $\lambda = \SI{650}{nm}$.
    \item \textbf{Synthèse :} En vous appuyant sur ces deux observations (diffraction par un obstacle et interférences à deux fentes), rédiger un argumentaire structuré (10-15 lignes) montrant que la lumière présente un comportement ondulatoire. Préciser le rôle de la cohérence dans l'expérience 2.
    \item \textbf{Application astrophysique :} On observe le spectre d'une galaxie lointaine. La raie H$_\alpha$ de l'hydrogène, qui est à $\lambda_0 = \SI{656,3}{nm}$ au laboratoire, est observée à $\lambda' = \SI{695,7}{nm}$.
    \begin{enumerate}
        \item Calculer le décalage spectral $z = \frac{\lambda' - \lambda_0}{\lambda_0}$.
        \item En déduire la vitesse de récession $v$ de la galaxie (on utilisera l'approximation $v = zc$ pour $z \ll 1$, avec $c = \SI{3,00 \times 10^8}{m/s}$).
        \item Expliquer pourquoi ce décalage vers le rouge, observé pour la quasi-totalité des galaxies, est une preuve majeure de l'expansion de l'Univers.
    \end{enumerate}
\end{enumerate}
\end{exercice}

\begin{exercice}[Bac Type : Radar routier et effet Doppler-Fizeau]
\textbf{Contexte :} Sur la route nationale N°3 (Yaoundé - Douala), un automobiliste est verbalisé pour excès de vitesse. Le radar embarqué dans le véhicule de la police a émis une onde électromagnétique de fréquence $f_0 = \SI{24,15}{GHz}$ (bande K). L'onde réfléchie par la voiture a été reçue avec une fréquence $f'' = f_0 + \Delta f$, où $\Delta f = \SI{3200}{Hz}$. L'automobiliste conteste, arguant que le radar était mal calibré.
\textbf{Données :} $c = \SI{3,00 \times 10^8}{m/s}$. Vitesse autorisée : $\SI{90}{km/h}$.
\textbf{Questions :}
\begin{enumerate}
    \item \textbf{Principe qualitatif :} Expliquer le principe de la mesure de vitesse par radar à effet Doppler. Pourquoi y a-t-il \emph{deux} décalages de fréquence successifs (émission vers la cible, puis réflexion vers le radar) ? Schématiser les vecteurs vitesse et ondes.
    \item \textbf{Relation quantitative :} Montrer que pour une vitesse $v \ll c$, le décalage de fréquence total $\Delta f$ entre l'onde émise et l'onde reçue après réflexion sur une cible en approche à la vitesse $v$ vaut approximativement $\Delta f \approx \dfrac{2v}{\lambda_0} = \dfrac{2v f_0}{c}$.
    \item \textbf{Calcul de la vitesse :} Calculer la vitesse $v$ du véhicule en $\SI{}{m/s}$ puis en $\SI{}{km/h}$.
    \item \textbf{Analyse du contestation :} L'automobiliste a-t-il raison de contester ? Justifier.
    \item \textbf{Erreur de mesure :} Si l'incertitude sur la mesure de $\Delta f$ est $\pm \SI{50}{Hz}$, calculer l'incertitude sur la vitesse $\Delta v$ (en $\SI{}{km/h}$). La verbalisation est-elle robuste ?
\end{enumerate}
\end{exercice}

\begin{exercice}[Bac Type : Caractérisation d'une source laser par diffraction et interférences]
\textbf{Contexte :} Dans le laboratoire de physique de l'ENSET de Douala, on cherche à caractériser une diode laser destinée aux télécommunications par fibre optique. On dispose d'un banc d'optique, d'une fente unique de largeur connue $a = \SI{80,0 \pm 0,5}{\mu m}$, de deux fentes de Young d'écartement inconnu $a_Y$, et d'un écran sur rail gradué.
\textbf{Expérience A (Diffraction) :} La fente unique est éclairée par le laser. On mesure la distance $D_A = \SI{1,500 \pm 0,002}{m}$ fente-écran. La largeur de la tache centrale (premiers minima) est $L_A = \SI{24,6 \pm 0,2}{mm}$.
\textbf{Expérience B (Interférences) :} On remplace par les fentes de Young. On mesure $D_B = \SI{2,000 \pm 0,002}{m}$. On compte $N = 50$ interfranges sur une longueur $L_B = \SI{78,5 \pm 0,3}{mm}$.
\textbf{Questions :}
\begin{enumerate}
    \item \textbf{Longueur d'onde :} À partir de l'expérience A, établir l'expression de $\lambda$ en fonction de $a$, $D_A$, $L_A$. Calculer $\lambda$ et son incertitude absolue $\Delta \lambda$. Exprimer le résultat en nm.
    \item \textbf{Écartement des fentes de Young :} À partir de l'expérience B, établir l'expression de $a_Y$ en fonction de $\lambda$, $D_B$, $N$, $L_B$. Calculer $a_Y$ et son incertitude $\Delta a_Y$. Exprimer en mm.
    \item \textbf{Cohérence :} Les deux valeurs de $\lambda$ (si on en déduisait une depuis l'expérience B en supposant $a_Y$ connu) sont-elles compatibles ? Discuter.
    \item \textbf{Lumière blanche :} Si l'on éclaire les fentes de Young avec une lampe à incandescence (lumière blanche) munie d'un filtre rouge ($\lambda \approx \SI{650}{nm}$) puis bleu ($\lambda \approx \SI{450}{nm}$), décrire l'évolution de l'interfrange et de la visibilité des franges. Pourquoi utilise-t-on une source monochromatique (laser) pour ce type de mesure de précision ?
    \item \textbf{Application télécom :} La diode laser testée émet à $\lambda = \SI{1550}{nm}$ (fenêtre de transparence de la silice). Calculer l'interfrange attendu pour $D = \SI{2,00}{m}$ et $a_Y = \SI{0,40}{mm}$. Discuter de la faisabilité de la mesure (taille des photodétecteurs, bruit).
\end{enumerate}
\end{exercice}

\newpage

\coursec{Corrigés des exercices}

\courssub{Corrigés de la rubrique « Je vérifie mes connaissances »}

\begin{corrige}[Exercice 1 — QCM : notions fondamentales]
\textbf{Réponse 3.}
\begin{itemize}
    \item[1.] \textbf{Faux.} L'écart angulaire de la tache centrale de diffraction vaut $\theta \approx \dfrac{\lambda}{a}$ : il est \emph{inversement} proportionnel à la largeur $a$ de la fente. La formule $\theta = a/\lambda$ est en outre non homogène si l'on y lit un angle sans dimension : $[a/\lambda] = 1$ certes, mais elle prédit que la tache s'élargit quand la fente s'élargit, contrairement à l'expérience.
    \item[2.] \textbf{Faux.} Les franges des fentes de Young sont \textbf{non localisées} : elles sont observables sur tout écran placé derrière les fentes (ou au foyer image d'une lentille). Ce ne sont pas des franges localisées comme celles des anneaux de Newton.
    \item[3.] \textbf{Vrai.} L'interfrange vaut $i = \dfrac{\lambda D}{a}$ (dans l'approximation des petits angles, $D \gg a$). Analyse dimensionnelle : $[\lambda D/a] = \dfrac{L \times L}{L} = L$, homogène à une longueur.
    \item[4.] \textbf{Faux.} La formule approchée $\dfrac{\Delta f}{f_0} \approx \dfrac{v}{c}$ fait intervenir explicitement la célérité $c$ de la lumière. Il est vrai que seule la vitesse \emph{relative} source-observateur intervient (contrairement au son, où le milieu compte), mais $c$ apparaît dans l'expression du décalage.
\end{itemize}
\end{corrige}

\begin{corrige}[Exercice 2 — Vrai ou Faux ? Justifiez]
\begin{enumerate}
    \item \textbf{Faux.} L'écart angulaire vaut $\theta \approx \dfrac{\lambda}{a}$. Si $a$ diminue, $\theta$ \emph{augmente} : la tache centrale s'\textbf{élargit}. C'est d'ailleurs le test expérimental le plus simple de la diffraction : plus l'ouverture est fine, plus la figure est étalée.
    \item \textbf{Vrai.} En $x = 0$, la différence de marche est $\delta = 0$ pour \emph{toutes} les longueurs d'onde : toutes les radiations interfèrent constructivement, la frange centrale est \textbf{blanche}. Pour $k \neq 0$, la position $x_k = k\dfrac{\lambda D}{a}$ dépend de $\lambda$ : les systèmes de franges de chaque couleur sont décalés, d'où des franges \textbf{colorées} (bleu vers le centre, rouge vers l'extérieur).
    \item \textbf{Faux.} Une source qui \emph{s'éloigne} voit ses longueurs d'onde \emph{allongées} : décalage vers le \textbf{rouge} (\emph{redshift}, $z > 0$). Le décalage vers le bleu (\emph{blueshift}) correspond à une source qui \emph{s'approche}.
    \item \textbf{Vrai.} La condition d'interférence constructive est $\delta = k\lambda$ avec $k \in \mathbb{Z}$, ordre d'interférence. La frange centrale correspond à $k = 0$ ($\delta = 0$) et est toujours brillante.
\end{enumerate}
\end{corrige}

\begin{corrige}[Exercice 3 — Définitions et calculs directs]
\begin{enumerate}
    \item L'\textbf{écart angulaire} $\theta$ de la tache centrale est le demi-angle au sommet du cône sous-tendu par la tache centrale depuis la fente, c'est-à-dire l'angle entre l'axe optique et la direction du premier minimum ($k = 1$). Dans l'approximation des petits angles :
    \[ \theta \approx \frac{\lambda}{a} \]
    (la largeur angulaire totale de la tache centrale, entre les deux premiers minima, vaut $2\theta = 2\lambda/a$).
    \item $a = \SI{0,20}{mm} = 2,0\times10^{-4}\,\text{m}$ ; $\lambda = \SI{632,8}{nm} = 6,328\times10^{-7}\,\text{m}$.
    \[ \theta = \frac{\lambda}{a} = \frac{6,328\times10^{-7}}{2,0\times10^{-4}} = \SI{3,16e-3}{rad} \]
    Conversion en degrés : $\theta = 3,16\times10^{-3} \times \dfrac{180}{\pi} \approx \SI{0,181}{\degree}$. L'approximation des petits angles est largement justifiée ($\theta \ll 1$ rad).
    \item $i = \dfrac{\lambda D}{a} = \dfrac{5,89\times10^{-7} \times 2,00}{5,0\times10^{-4}} = \dfrac{1,178\times10^{-6}}{5,0\times10^{-4}} = \SI{2,356e-3}{m} \approx \SI{2,36}{mm}$.
    \item Troisième frange brillante ($k = 3$) : $\delta = k\lambda = 3 \times 589 = \SI{1767}{nm} \approx \SI{1,77}{\mu m}$.
    Franges sombres : $\delta = \left(k+\tfrac{1}{2}\right)\lambda$. En partant du centre : 1$^{re}$ sombre $\delta = \lambda/2$ ($k=0$) ; 2$^{e}$ sombre $\delta = \dfrac{3\lambda}{2}$ ($k=1$). Donc $\delta = 1,5 \times 589 = \SI{883,5}{nm} \approx \SI{884}{nm}$.
\end{enumerate}
\end{corrige}

\begin{corrige}[Exercice 4 — Identification d'ordres d'interférence]
\begin{enumerate}
    \item Interfrange : $i = \dfrac{\lambda D}{a} = \dfrac{600\times10^{-9} \times 1,50}{0,40\times10^{-3}} = \dfrac{9,00\times10^{-7}}{4,0\times10^{-4}} = \SI{2,25e-3}{m} = \SI{2,25}{mm}$.
    \item Rapport : $\dfrac{x_M}{i} = \dfrac{4,5}{2,25} = 2$, nombre \textbf{entier}. Le point M est donc sur une \textbf{frange brillante} d'ordre $k = 2$.
    (Si le rapport avait été un demi-entier, M aurait été sur une frange sombre ; s'il n'était ni entier ni demi-entier, M serait en un point d'éclairement intermédiaire.)
    \item Différence de marche : $\delta_M = k\lambda = 2 \times 600 = \SI{1200}{nm} = \SI{1,20}{\mu m}$.
    Vérification directe : $\delta_M = \dfrac{a\,x_M}{D} = \dfrac{0,40\times10^{-3} \times 4,5\times10^{-3}}{1,50} = \SI{1,20e-6}{m}$. Cohérent.
\end{enumerate}
\end{corrige}

\courssub{Corrigés de la rubrique « Je m'entraîne »}

\begin{corrige}[Exercice 5 — Diffraction par une fente : détermination de $\lambda$ et prévision]
\begin{enumerate}
    \item Les premiers minima vérifient $a_1 \sin\theta \approx a_1 \theta = \pm\lambda$, soit $\theta = \pm\dfrac{\lambda}{a_1}$. Sur l'écran à distance $D$, chaque minimum est à la position $x = \pm D\tan\theta \approx \pm D\theta = \pm \dfrac{D\lambda}{a_1}$. La largeur totale de la tache centrale est donc :
    \[ \boxed{L_1 = \frac{2D\lambda}{a_1}} \]
    \item On isole $\lambda$ : $\lambda = \dfrac{L_1 a_1}{2D} = \dfrac{22,0\times10^{-3} \times 0,15\times10^{-3}}{2 \times 2,50} = \dfrac{3,30\times10^{-6}}{5,00} = \SI{6,60e-7}{m} = \boxed{\SI{660}{nm}}$.
    (Laser rouge, valeur tout à fait plausible pour une diode laser.)
    \item $L$ est inversement proportionnelle à $a$ à $D$ et $\lambda$ fixés : $L_2 = L_1 \times \dfrac{a_1}{a_2} = 2 L_1 = \boxed{\SI{44,0}{mm}}$.
    Vérification : $L_2 = \dfrac{2 \times 2,50 \times 6,60\times10^{-7}}{7,5\times10^{-5}} = \dfrac{3,30\times10^{-6}}{7,5\times10^{-5}} = \SI{4,40e-2}{m} = \SI{44,0}{mm}$. Conforme.
\end{enumerate}
\textbf{Point de vigilance :} ne pas oublier le facteur 2 : la tache centrale s'étend \emph{de part et d'autre} de l'axe, sa largeur est $2D\lambda/a$ et non $D\lambda/a$.
\end{corrige}

\begin{corrige}[Exercice 6 — Fentes de Young : analyse complète]
\begin{enumerate}
    \item $i = \dfrac{\lambda D}{a} = \dfrac{546,1\times10^{-9} \times 3,00}{0,60\times10^{-3}} = \dfrac{1,6383\times10^{-6}}{6,0\times10^{-4}} = \SI{2,73e-3}{m} = \boxed{\SI{2,73}{mm}}$.
    \item Frange brillante d'ordre 4 : $x_4 = 4i = 4 \times 2,7305 = \boxed{\SI{10,9}{mm}}$.
    \item Franges sombres : $x'_k = \left(k+\tfrac{1}{2}\right)i$. La \emph{troisième} frange sombre à partir du centre correspond à $k = 2$ (la première étant $k=0$) :
    $x'_3 = 2,5\,i = 2,5 \times 2,7305 = \boxed{\SI{6,83}{mm}}$.
    \item Frange brillante d'ordre 2 : $x_2 = 2i = \SI{5,46}{mm}$. Distance cherchée :
    $d = x'_3 - x_2 = 2,5i - 2i = 0,5\,i = \boxed{\SI{1,37}{mm}}$.
    Résultat général utile : une frange sombre est toujours à une \emph{demi-interfrange} de la frange brillante voisine.
    \item Dans un milieu d'indice $n$, la célérité devient $c/n$ et la longueur d'onde $\lambda' = \dfrac{\lambda}{n}$ (la fréquence, elle, est inchangée). L'interfrange devient :
    \[ i' = \frac{\lambda' D}{a} = \frac{i}{n} = \frac{2,7305}{1,33} = \boxed{\SI{2,05}{mm}} \]
    Les franges se resserrent : le dispositif est plus « compact » dans l'eau.
\end{enumerate}
\textbf{Point de vigilance :} la numérotation des franges sombres commence à $k = 0$ : la « troisième frange sombre » est d'ordre $k = 2$, d'où le facteur $2{,}5$ et non $3{,}5$.
\end{corrige}

\begin{corrige}[Exercice 7 — Mesure de l'écartement des fentes et incertitude]
\begin{enumerate}
    \item La distance entre la frange d'ordre 0 et la frange d'ordre $N$ vaut $L_N = N i = N\dfrac{\lambda D}{a}$, d'où :
    \[ \boxed{a = \frac{N\lambda D}{L_N}} \]
    \item $a = \dfrac{20 \times 632,8\times10^{-9} \times 2,000}{36,4\times10^{-3}} = \dfrac{2,5312\times10^{-5}}{3,64\times10^{-2}} = \SI{6,954e-4}{m} \approx \SI{0,695}{mm}$.
    \item Incertitudes relatives :
    \begin{align*}
    \frac{\Delta \lambda}{\lambda} &= \frac{0,1}{632,8} = 1,6\times10^{-4} \\
    \frac{\Delta D}{D} &= \frac{0,005}{2,000} = 2,5\times10^{-3} \\
    \frac{\Delta L_N}{L_N} &= \frac{0,2}{36,4} = 5,5\times10^{-3}
    \end{align*}
    Propagation quadratique (incertitudes indépendantes) :
    \[ \frac{\Delta a}{a} = \sqrt{(1,6\times10^{-4})^2 + (2,5\times10^{-3})^2 + (5,5\times10^{-3})^2} = \sqrt{3,65\times10^{-5}} \approx 6,0\times10^{-3} \]
    soit $\dfrac{\Delta a}{a} \approx 0,60\,\%$. C'est la mesure de $L_N$ qui domine l'incertitude.
    \item $\Delta a = 0,695 \times 6,0\times10^{-3} \approx \SI{0,004}{mm}$.
    Résultat final : $\boxed{a = \SI{0,695 \pm 0,004}{mm}}$ (incertitude arrondie à un chiffre significatif, valeur arrondie au même rang décimal).
\end{enumerate}
\textbf{Point de vigilance :} mesurer $N$ interfranges puis diviser par $N$ réduit l'incertitude relative sur $i$ d'un facteur $N$ : c'est la raison d'être de cette méthode.
\end{corrige}

\begin{corrige}[Exercice 8 — Interférences en lumière polychromatique]
\begin{enumerate}
    \item $i_1 = \dfrac{\lambda_1 D}{a} = \dfrac{589,0\times10^{-9} \times 2,00}{0,50\times10^{-3}} = \boxed{\SI{2,356}{mm}}$ ;
    $i_2 = \dfrac{\lambda_2 D}{a} = \dfrac{460,0\times10^{-9} \times 2,00}{0,50\times10^{-3}} = \boxed{\SI{1,840}{mm}}$.
    \item Condition de coïncidence de deux franges brillantes : $k_1 i_1 = k_2 i_2 \iff k_1 \lambda_1 = k_2 \lambda_2$, soit $\dfrac{k_1}{k_2} = \dfrac{\lambda_2}{\lambda_1} = \dfrac{460}{589}$.
    La fraction est irréductible ($460 = 2^2 \times 5 \times 23$ ; $589 = 19 \times 31$ : aucun facteur commun). Les plus petits entiers sont donc $\boxed{k_1 = 460}$ et $\boxed{k_2 = 589}$.
    Différence de marche commune : $\delta_c = k_1\lambda_1 = 460 \times 589,0 = \SI{270940}{nm} \approx \SI{0,271}{mm}$.
    \item Abscisse : $x_c = k_1 i_1 = 460 \times 2,356 = \SI{1083,8}{mm} \approx \boxed{\SI{1,08}{m}}$.
    Cette coïncidence est située très loin de la frange centrale : en pratique, hors de l'écran et dans une zone où les franges ont perdu toute visibilité. La coïncidence exacte est donc ici un résultat \emph{théorique}.
    \item En lumière blanche : la frange centrale ($k=0$) est \textbf{blanche} car $\delta = 0$ pour toutes les longueurs d'onde. Les franges voisines ($k = \pm 1, \pm 2$) sont \textbf{colorées} : teinte bleutée côté centre (petites $\lambda$, petit interfrange), rougeoyante côté extérieur. Dès que les ordres se recouvrent ($k \gtrsim 3$), toutes les couleurs se superposent : l'écran devient uniformément éclairé (franges « lavées », visibilité nulle).
\end{enumerate}
\end{corrige}

\courssub{Corrigés de la rubrique « Je me prépare au Bac »}

\begin{corrige}[Exercice 9 — Bac Type : Nature ondulatoire de la lumière et expansion de l'Univers]
\begin{enumerate}
    \item \textbf{Diffraction par le fil.} D'après le principe de Babinet, le fil de diamètre $d$ donne la même figure de diffraction qu'une fente de largeur $a = d$. Largeur théorique de la tache centrale :
    \[ L_{\text{th}} = \frac{2D\lambda}{d} = \frac{2 \times 2,00 \times 650\times10^{-9}}{70\times10^{-6}} = \SI{3,71e-2}{m} = \SI{37,1}{mm} \]
    La mesure $L = \SI{37,1}{mm}$ coïncide avec la théorie à mieux que $0,1\,\%$ : accord excellent.
    Longueur d'onde expérimentale : $\lambda_{\text{exp}} = \dfrac{L\,d}{2D} = \dfrac{37,1\times10^{-3} \times 70\times10^{-6}}{2 \times 2,00} = \SI{6,49e-7}{m} \approx \SI{649}{nm}$, compatible avec la valeur nominale $\SI{650}{nm}$ (écart relatif $\approx 0,1\,\%$).
    \item \textbf{Interférences de Young.} $i_{\text{th}} = \dfrac{\lambda D}{a} = \dfrac{650\times10^{-9} \times 2,00}{0,30\times10^{-3}} = \SI{4,33e-3}{m} = \SI{4,33}{mm}$, identique à la valeur mesurée. Cohérence parfaite.
    \item \textbf{Synthèse (argumentaire).} Dans l'expérience 1, l'optique géométrique prévoirait une simple ombre nette du fil ; or on observe une tache centrale \emph{élargie} encadrée de franges : la lumière se répand au-delà de l'ombre géométrique, ce qui est inexplicable avec un modèle purement corpusculaire en propagation rectiligne. Dans l'expérience 2, deux faisceaux issus de deux fentes se superposent et produisent des zones alternativement brillantes et sombres : additionner deux lumières peut donner de l'obscurité, phénomène uniquement explicable par l'addition d'\emph{amplitudes} ondulatoires pouvant s'annuler (interférences destructives, $\delta = (k+\tfrac12)\lambda$). Les deux expériences, quantitatives, sont décrites par la \emph{même} grandeur $\lambda = \SI{650}{nm}$, caractéristique d'une onde. La lumière présente donc un comportement ondulatoire, modélisé par le principe de Huygens-Fresnel. Dans l'expérience 2, la \textbf{cohérence} des deux sources secondaires est indispensable : elles doivent être issues de la même source primaire (même fréquence, différence de phase constante en chaque point) ; sans cohérence, le terme d'interférence s'annule en moyenne temporelle et aucune frange n'est observable.
    \item \textbf{Astrophysique.}
    \begin{enumerate}
        \item $z = \dfrac{\lambda' - \lambda_0}{\lambda_0} = \dfrac{695,7 - 656,3}{656,3} = \dfrac{39,4}{656,3} = \boxed{6,00\times10^{-2}}$ (sans unité).
        \item $v = zc = 0,0600 \times 3,00\times10^{8} = \boxed{\SI{1,80e7}{m/s}}$, soit $\SI{18000}{km/s}$ (environ $6\,\%$ de $c$ : l'approximation $v = zc$ reste acceptable).
        \item Le décalage est systématiquement vers le \emph{rouge} ($z > 0$) pour la quasi-totalité des galaxies, et d'autant plus grand qu'elles sont lointaines (loi de Hubble-Lemaître $v = H_0 d$). Tout se passe comme si l'Univers entier s'étirait : ce n'est pas un simple effet Doppler dû à des mouvements dans un espace figé, mais l'\textbf{expansion de l'espace lui-même} qui étire la longueur d'onde des photons durant leur trajet. C'est l'une des preuves observationnelles majeures du modèle du Big Bang.
    \end{enumerate}
\end{enumerate}
\textbf{Barème indicatif (sur 10) :} Q1 : 2 pts (formule, calcul, comparaison) ; Q2 : 1 pt ; Q3 : 3 pts (deux arguments + rôle de la cohérence) ; Q4a : 1 pt ; Q4b : 1 pt ; Q4c : 2 pts.
\textbf{Points de vigilance :} citer le principe de Babinet pour justifier l'assimilation fil/fente ; vérifier l'hypothèse des petits angles avant d'écrire $L = 2D\lambda/a$ ; dans l'argumentaire, expliciter que « lumière + lumière = obscurité » est la signature des interférences ; préciser que $z$ est sans dimension.
\end{corrige}

\begin{corrige}[Exercice 10 — Bac Type : Radar routier et effet Doppler-Fizeau]
\begin{enumerate}
    \item \textbf{Principe qualitatif.} Le radar émet une onde de fréquence $f_0$ vers le véhicule en approche. \emph{Premier décalage} : le véhicule joue le rôle d'observateur mobile se rapprochant de la source ; il reçoit une fréquence $f_1 > f_0$. \emph{Second décalage} : le véhicule réémet (réfléchit) cette onde ; il joue alors le rôle de source mobile se rapprochant du radar, qui reçoit $f'' > f_1$. Le décalage total $\Delta f = f'' - f_0$, proportionnel à la vitesse $v$, permet la mesure.
    \begin{popfigure}
    \begin{tikzpicture}[scale=0.9, >=Stealth]
        \draw[fill=popblueL] (-4,0) rectangle (-3.4,0.8);
        \node[below] at (-3.7,-0.1) {\small Radar (fixe)};
        \draw[fill=poporangeL] (2.6,0) rectangle (3.8,0.7);
        \node[below] at (3.2,-0.1) {\small Voiture};
        \draw[->, popdark, thick] (2.6,0.35) -- (1.6,0.35) node[midway, above]{$\vec{v}$};
        \draw[->, popblue, thick, decorate, decoration={snake, amplitude=1.5pt}] (-3.3,0.9) -- (2.4,0.9) node[midway, above]{$f_0$};
        \draw[->, poppurple, thick, decorate, decoration={snake, amplitude=1.5pt}] (2.4,-0.5) -- (-3.3,-0.5) node[midway, below]{$f'' > f_0$};
    \end{tikzpicture}
    \legende{Double effet Doppler : onde incidente $f_0$, onde réfléchie $f''$ par la cible en approche.}
    \end{popfigure}
    \item \textbf{Démonstration.} Pour $v \ll c$, le décalage relatif par réception mobile est $\dfrac{\Delta f}{f} \approx \dfrac{v}{c}$.
    Réception par la cible : $f_1 = f_0\left(1 + \dfrac{v}{c}\right)$.
    Réémission par la cible (source mobile) : $f'' = f_1\left(1 + \dfrac{v}{c}\right) = f_0\left(1 + \dfrac{v}{c}\right)^2 = f_0\left(1 + \dfrac{2v}{c} + \dfrac{v^2}{c^2}\right)$.
    Comme $v \ll c$, le terme $\dfrac{v^2}{c^2}$ est négligeable : $f'' \approx f_0\left(1 + \dfrac{2v}{c}\right)$, d'où :
    \[ \boxed{\Delta f = f'' - f_0 \approx \frac{2v f_0}{c} = \frac{2v}{\lambda_0}} \]
    \item \textbf{Calcul.} $v = \dfrac{\Delta f \cdot c}{2 f_0} = \dfrac{3200 \times 3,00\times10^{8}}{2 \times 24,15\times10^{9}} = \dfrac{9,60\times10^{11}}{4,83\times10^{10}} = \SI{19,9}{m/s}$.
    Conversion : $v = 19,9 \times 3,6 \approx \boxed{\SI{71,6}{km/h}}$.
    \item \textbf{Contestation.} $\SI{71,6}{km/h} < \SI{90}{km/h}$ : le véhicule respectait la limitation. \textbf{L'automobiliste a raison de contester} : la verbalisation est infondée (défaut de calibration ou erreur de lecture du radar).
    \item \textbf{Incertitude.} $\Delta v = \dfrac{\Delta(\Delta f) \cdot c}{2 f_0} = \dfrac{50 \times 3,00\times10^{8}}{4,83\times10^{10}} = \SI{0,31}{m/s} \approx \SI{1,1}{km/h}$.
    Résultat : $v = \SI{71,6 \pm 1,1}{km/h}$. Même la borne supérieure ($\SI{72,7}{km/h}$) reste très inférieure à $\SI{90}{km/h}$ : la verbalisation n'est \textbf{pas robuste}.
\end{enumerate}
\textbf{Barème indicatif (sur 10) :} Q1 : 2 pts (deux décalages identifiés + schéma) ; Q2 : 3 pts (deux étapes + approximation du terme en $v^2/c^2$ explicitement négligé) ; Q3 : 2 pts ; Q4 : 1 pt ; Q5 : 2 pts.
\textbf{Points de vigilance :} le facteur 2 (aller-retour) est la clef de l'exercice ; justifier l'approximation en écrivant le développement complet avant de négliger $v^2/c^2$ ; toujours convertir en km/h ($\times 3,6$) pour conclure dans le contexte routier ; conclure par une phrase de synthèse argumentée.
\end{corrige}

\begin{corrige}[Exercice 11 — Bac Type : Caractérisation d'une source laser par diffraction et interférences]
\begin{enumerate}
    \item \textbf{Longueur d'onde (diffraction).} $L_A = \dfrac{2D_A\lambda}{a} \implies \boxed{\lambda = \dfrac{a L_A}{2 D_A}}$.
    \[ \lambda = \frac{80,0\times10^{-6} \times 24,6\times10^{-3}}{2 \times 1,500} = \frac{1,968\times10^{-6}}{3,000} = \SI{6,56e-7}{m} = \SI{656}{nm} \]
    Incertitudes relatives : $\dfrac{\Delta a}{a} = \dfrac{0,5}{80,0} = 6,3\times10^{-3}$ ; $\dfrac{\Delta L_A}{L_A} = \dfrac{0,2}{24,6} = 8,1\times10^{-3}$ ; $\dfrac{\Delta D_A}{D_A} = \dfrac{0,002}{1,500} = 1,3\times10^{-3}$.
    \[ \frac{\Delta \lambda}{\lambda} = \sqrt{(6,3\times10^{-3})^2 + (8,1\times10^{-3})^2 + (1,3\times10^{-3})^2} \approx 1,0\times10^{-2} \]
    $\Delta\lambda \approx 656 \times 0,010 \approx \SI{7}{nm}$. Résultat : $\boxed{\lambda = \SI{656 \pm 7}{nm}}$.
    \item \textbf{Écartement des fentes de Young.} $L_B = N i = N\dfrac{\lambda D_B}{a_Y} \implies \boxed{a_Y = \dfrac{N\lambda D_B}{L_B}}$.
    \[ a_Y = \frac{50 \times 656\times10^{-9} \times 2,000}{78,5\times10^{-3}} = \frac{6,56\times10^{-5}}{7,85\times10^{-2}} = \SI{8,36e-4}{m} = \SI{0,836}{mm} \]
    Incertitudes : $\dfrac{\Delta\lambda}{\lambda} = 1,0\times10^{-2}$ ; $\dfrac{\Delta D_B}{D_B} = \dfrac{0,002}{2,000} = 1,0\times10^{-3}$ ; $\dfrac{\Delta L_B}{L_B} = \dfrac{0,3}{78,5} = 3,8\times10^{-3}$.
    \[ \frac{\Delta a_Y}{a_Y} = \sqrt{(1,0\times10^{-2})^2 + (1,0\times10^{-3})^2 + (3,8\times10^{-3})^2} \approx 1,1\times10^{-2} \]
    $\Delta a_Y \approx 0,836 \times 0,011 \approx \SI{0,009}{mm}$. Résultat : $\boxed{a_Y = \SI{0,836 \pm 0,009}{mm}}$.
    \item \textbf{Cohérence des deux déterminations.} En recalculant $\lambda_B = \dfrac{a_Y L_B}{N D_B}$ avec la valeur de $a_Y$ issue de la question 2, on retrouve $\lambda_B = \SI{656}{nm}$ par construction : ce n'est pas un test indépendant. Le test de cohérence pertinent consiste à comparer $\lambda = \SI{656 \pm 7}{nm}$ à une valeur de référence (fiche constructeur de la diode, ou mesure au spectromètre) : si l'intervalle $[\SI{649}{nm} ; \SI{663}{nm}]$ contient la référence, les déterminations sont compatibles. On remarque que l'incertitude est dominée par la lecture de $L_A$ ($8,1\times10^{-3}$) et par celle de $a$ ($6,3\times10^{-3}$) : pour améliorer la précision, il faudrait mesurer la tache centrale avec plus de soin (ou utiliser $D$ plus grand pour élargir la figure) et connaître $a$ plus finement.
    \item \textbf{Lumière blanche et filtres.} Avec le filtre rouge ($\lambda \approx \SI{650}{nm}$), l'interfrange est $i_R = \dfrac{\lambda_R D}{a_Y}$, plus grande qu'avec le filtre bleu ($\lambda \approx \SI{450}{nm}$) : $i_B/i_R = 450/650 \approx 0,69$, les franges bleues sont plus serrées. La visibilité reste bonne dans les deux cas car le filtre sélectionne une bande étroite (cohérence temporelle suffisante). Sans filtre (lumière blanche) : frange centrale blanche, quelques franges colorées, puis brouillage complet. On utilise un \textbf{laser} pour la métrologie de précision car sa grande cohérence temporelle ($\Delta\lambda$ très faible) et spatiale garantit des franges de visibilité maximale sur tout l'écran, ce qui permet de mesurer un grand nombre $N$ d'interfranges et de réduire l'incertitude.
    \item \textbf{Application télécom.} On considère un dispositif de Young d'écartement $a_Y = \SI{0,40}{mm}$. $i = \dfrac{\lambda D}{a_Y} = \dfrac{1550\times10^{-9} \times 2,00}{0,40\times10^{-3}} = \boxed{\SI{7,75}{mm}}$.
    Faisabilité : un interfrange de $\SI{7,75}{mm}$ est parfaitement résolu par une photodiode InGaAs (zone sensible $\sim \SI{1}{mm}$) ou une caméra infrarouge (pixels $\sim \SI{10}{\mu m}$). Les sources de bruit (lumière ambiante, bruit thermique du détecteur) sont maîtrisées par un filtre passe-bande centré sur $\SI{1550}{nm}$ et par modulation du faisceau. La mesure est donc tout à fait réalisable en laboratoire de caractérisation optique.
\end{enumerate}
\textbf{Barème indicatif (sur 12) :} Q1 : 3 pts (expression, calcul, incertitude) ; Q2 : 3 pts ; Q3 : 2 pts ; Q4 : 2 pts ; Q5 : 2 pts.
\textbf{Points de vigilance :} dans la propagation des incertitudes, combiner les incertitudes \emph{relatives} en quadrature (somme des carrés) ; ne pas oublier le facteur 2 dans $L_A = 2D_A\lambda/a$ ; arrondir l'incertitude à un chiffre significatif et la valeur au même rang décimal ; à la question 3, reconnaître explicitement que le recalcul de $\lambda$ n'est pas un test indépendant — c'est le raisonnement critique attendu au Bac.
\end{corrige}

\newpage

\coursec{Fiche bilan}

\begin{popfigure}
\begin{center}\tcbox[colback=popdark,colframe=popdark,coltext=white,arc=9pt,boxsep=4pt,left=6pt,right=6pt]{\bfseries\small Lumière : aspect ondulatoire}\end{center}
\vspace{-2pt}
\begin{tcbraster}[raster columns=3,raster equal height=rows,raster column skip=6pt,raster row skip=6pt,enhanced,blankest]
\begin{tcolorbox}[enhanced,colback=white,colframe=popblue,boxrule=0.9pt,arc=3pt,title={Diffraction},fonttitle=\bfseries\scriptsize,coltitle=white,colbacktitle=popblue,left=4pt,right=4pt,top=2pt,bottom=2pt,boxsep=1pt,toptitle=1.5pt,bottomtitle=1.5pt,halign=left]
\begin{itemize}[nosep,leftmargin=*,itemsep=1pt,font=\scriptsize]
\item $\theta = \dfrac{\lambda}{a}$
\item $L = \dfrac{2\lambda D}{a}$
\end{itemize}
\end{tcolorbox}
\begin{tcolorbox}[enhanced,colback=white,colframe=poporange,boxrule=0.9pt,arc=3pt,title={Interférences},fonttitle=\bfseries\scriptsize,coltitle=white,colbacktitle=poporange,left=4pt,right=4pt,top=2pt,bottom=2pt,boxsep=1pt,toptitle=1.5pt,bottomtitle=1.5pt,halign=left]
\begin{itemize}[nosep,leftmargin=*,itemsep=1pt,font=\scriptsize]
\item Sources cohérentes
\item $\delta = \dfrac{ax}{D}$
\end{itemize}
\end{tcolorbox}
\begin{tcolorbox}[enhanced,colback=white,colframe=popgreen,boxrule=0.9pt,arc=3pt,title={Franges de Young},fonttitle=\bfseries\scriptsize,coltitle=white,colbacktitle=popgreen,left=4pt,right=4pt,top=2pt,bottom=2pt,boxsep=1pt,toptitle=1.5pt,bottomtitle=1.5pt,halign=left]
\begin{itemize}[nosep,leftmargin=*,itemsep=1pt,font=\scriptsize]
\item Brillante : $\delta = k\lambda$
\item Sombre : $\delta = (k+\tfrac{1}{2})\lambda$
\end{itemize}
\end{tcolorbox}
\begin{tcolorbox}[enhanced,colback=white,colframe=poppurple,boxrule=0.9pt,arc=3pt,title={Interfrange},fonttitle=\bfseries\scriptsize,coltitle=white,colbacktitle=poppurple,left=4pt,right=4pt,top=2pt,bottom=2pt,boxsep=1pt,toptitle=1.5pt,bottomtitle=1.5pt,halign=left]
\begin{itemize}[nosep,leftmargin=*,itemsep=1pt,font=\scriptsize]
\item $i = \dfrac{\lambda D}{a}$
\item Mesure de $\lambda$
\end{itemize}
\end{tcolorbox}
\begin{tcolorbox}[enhanced,colback=white,colframe=poppink,boxrule=0.9pt,arc=3pt,title={Lumière blanche},fonttitle=\bfseries\scriptsize,coltitle=white,colbacktitle=poppink,left=4pt,right=4pt,top=2pt,bottom=2pt,boxsep=1pt,toptitle=1.5pt,bottomtitle=1.5pt,halign=left]
\begin{itemize}[nosep,leftmargin=*,itemsep=1pt,font=\scriptsize]
\item Coïncidences $k_1\lambda_1 = k_2\lambda_2$
\item Frange centrale blanche
\end{itemize}
\end{tcolorbox}
\begin{tcolorbox}[enhanced,colback=white,colframe=popgold,boxrule=0.9pt,arc=3pt,title={Effet Doppler},fonttitle=\bfseries\scriptsize,coltitle=white,colbacktitle=popgold,left=4pt,right=4pt,top=2pt,bottom=2pt,boxsep=1pt,toptitle=1.5pt,bottomtitle=1.5pt,halign=left]
\begin{itemize}[nosep,leftmargin=*,itemsep=1pt,font=\scriptsize]
\item Approche : $\lambda$ perçue $\searrow$
\item Redshift des galaxies
\end{itemize}
\end{tcolorbox}
\end{tcbraster}

\legende{Carte mentale de la leçon : les cinq piliers de l'aspect ondulatoire de la lumière.}
\end{popfigure}

\begin{bilan}
\textbf{1. Définitions clés}
\begin{itemize}
  \item \cle{Diffraction} : modification de la propagation d'une onde (étalement angulaire) lorsqu'elle rencontre une ouverture ou un obstacle dont la dimension $a$ est de l'ordre de grandeur de la longueur d'onde $\lambda$.
  \item \cle{Interférences} : superposition de deux ondes lumineuses \textbf{cohérentes} (même fréquence, déphasage constant) donnant lieu à une répartition alternative de zones brillantes et sombres.
  \item \cle{Différence de marche} $\delta$ : différence des chemins optiques parcourus par les deux ondes depuis les sources jusqu'au point d'observation $M$.
  \item \cle{Ordre d'interférence} : $p = \dfrac{\delta}{\lambda}$ ; $p = k \in \mathbb{Z}$ (frange brillante), $p = k + \tfrac{1}{2}$ (frange sombre).
  \item \cle{Interfrange} $i$ : distance séparant deux franges consécutives de même nature.
  \item \cle{Effet Doppler} : modification de la fréquence perçue lorsque la source et le récepteur sont en mouvement relatif.
\end{itemize}

\textbf{2. Formules à connaître par c\oe ur}
\begin{center}
\begin{tabularx}{\linewidth}{|l|X|X|}
\hline
\rowcolor{popblueL} \textbf{Grandeur} & \textbf{Expression} & \textbf{Conditions / remarques} \\
\hline
Écart angulaire (diffraction) & $\theta = \dfrac{\lambda}{a}$ (rad) & $a$ : largeur de la fente ; valable car $\theta$ petit \\
\hline
Tache centrale de diffraction & $L = \dfrac{2\lambda D}{a}$ & $L$ : largeur sur l'écran à la distance $D$ \\
\hline
Différence de marche (Young) & $\delta = \dfrac{ax}{D}$ & $a$ : écart des fentes ; $x$ : position de $M$ \\
\hline
Frange brillante & $\delta = k\lambda$, $k \in \mathbb{Z}$ & Frange centrale : $k = 0$ \\
\hline
Frange sombre & $\delta = \left(k + \dfrac{1}{2}\right)\lambda$ & \\
\hline
Interfrange & $i = \dfrac{\lambda D}{a}$ & Homogène à une longueur : $\dfrac{\text{m}\times\text{m}}{\text{m}}$ \\
\hline
Coïncidence (2 couleurs) & $k_1\lambda_1 = k_2\lambda_2$ & Franges brillantes superposées \\
\hline
\end{tabularx}
\end{center}

\textbf{3. Méthodes express}
\begin{itemize}
  \item \textbf{Mesurer $\lambda$ par interférences} : mesurer la largeur de $n$ interfranges ($\ell = ni$), puis $\lambda = \dfrac{ia}{D} = \dfrac{\ell a}{nD}$. Grand $n$ $\Rightarrow$ meilleure précision.
  \item \textbf{Nature d'un point $M$} : calculer $\delta = \dfrac{ax}{D}$ puis $p = \dfrac{\delta}{\lambda}$. $p$ entier $\Rightarrow$ brillante ; $p$ demi-entier $\Rightarrow$ sombre.
  \item \textbf{Diffraction} : $\lambda$ inconnu $\Rightarrow$ $\lambda = \dfrac{La}{2D}$ ; toujours convertir $\lambda$ en mètres ($1\ \text{nm} = 10^{-9}\ \text{m}$).
  \item \textbf{Doppler (qualitatif)} : la source \emph{approche} $\Rightarrow$ fréquence perçue plus grande (son aigu, lumière décalée vers le bleu) ; la source \emph{s'éloigne} $\Rightarrow$ fréquence plus faible (son grave, décalage vers le rouge).
\end{itemize}

\textbf{4. Vocabulaire et pièges}
\begin{itemize}
  \item Diffraction et interférences prouvent la \cle{nature ondulatoire} de la lumière (le modèle particulaire ne les explique pas).
  \item Les franges de Young sont \textbf{non localisées} et équidistantes ; la frange centrale ($\delta = 0$) est brillante pour \emph{toutes} les longueurs d'onde : en lumière blanche elle est blanche, irisée de part et d'autre.
  \item Ne pas confondre $a$ (écart des fentes ou largeur de la fente) avec $D$ (distance fentes--écran).
\end{itemize}
\end{bilan}

\begin{aretenir}[Checklist avant le Bac]
\begin{itemize}
  \item[$\square$] Je sais définir la diffraction et citer la condition sur $a$ et $\lambda$.
  \item[$\square$] Je connais $\theta = \dfrac{\lambda}{a}$ et je sais en déduire la largeur de la tache centrale $L = \dfrac{2\lambda D}{a}$.
  \item[$\square$] Je sais expliquer pourquoi deux sources lumineuses indépendantes ne donnent pas d'interférences (cohérence).
  \item[$\square$] Je sais établir $\delta = \dfrac{ax}{D}$ pour les fentes de Young (approximation des petits angles).
  \item[$\square$] Je connais les conditions de franges brillantes ($\delta = k\lambda$) et sombres ($\delta = (k+\tfrac{1}{2})\lambda$).
  \item[$\square$] Je sais établir et exploiter l'interfrange $i = \dfrac{\lambda D}{a}$ (démonstration exigible).
  \item[$\square$] Je vérifie l'homogénéité de chaque formule (analyse dimensionnelle).
  \item[$\square$] Je sais déterminer la nature (brillante ou sombre) d'un point $M$ via l'ordre $p = \delta/\lambda$.
  \item[$\square$] Je sais exploiter une coïncidence en lumière polychromatique : $k_1\lambda_1 = k_2\lambda_2$.
  \item[$\square$] Je sais décrire l'aspect des franges en lumière blanche (centrale blanche, irisées).
  \item[$\square$] Je sais expliquer qualitativement l'effet Doppler et citer deux applications (radar routier, expansion de l'Univers).
  \item[$\square$] Je convertis toujours les longueurs d'onde en mètres et je donne le résultat avec son unité SI.
\end{itemize}
\end{aretenir}

\findecours

\end{document}
